\input{preamble}

% OK, start here.
%
\begin{document}

\title{Limites des espaces algébriques}


\maketitle

\phantomsection
\label{section-phantom}

\tableofcontents

\section{Introduction}
\label{section-introduction}

\noindent
Dans ce chapitre, nous rassemblons des résultats relatifs aux limites des espaces algébriques.
Un premier thème est la caractérisation des espaces algébriques $F$ localement
de présentation finie sur la base $S$ comme foncteurs compatibles aux limites.
Nous poursuivons par l'étude des limites de systèmes projectifs indexés par des
ensembles ordonnés filtrants (Catégories, définition \ref{categories-definition-directed-set})
dont les morphismes de transition sont affines. Nous discutons l'approximation
noethérienne absolue des espaces algébriques quasi-compacts et quasi-séparés,
suivant \cite{CLO}. Une autre approche est due à David Rydh (voir
\cite{rydh_approx}); ses résultats couvrent également l'approximation noethérienne
absolue de certains champs algébriques.


\section{Conventions}
\label{section-conventions}

\noindent
L'hypothèse générale est que tous les schémas sont contenus dans
un grand site fppf $\Sch_{fppf}$. Tous les anneaux $A$ considérés
ont la propriété que $\Spec(A)$ est (isomorphe à) un
objet de ce grand site.

\medskip\noindent
Soit $S$ un schéma et soit $X$ un espace algébrique sur $S$.
Dans ce chapitre et dans le suivant, nous écrirons $X \times_S X$
pour le produit de $X$ par lui-même (dans la catégorie des espaces algébriques
sur $S$), au lieu de $X \times X$.











\section{Morphismes de présentation finie}
\label{section-finite-presentation}

\noindent
Dans cette section, nous généralisons
Limites, proposition
\ref{limits-proposition-characterize-locally-finite-presentation}
aux morphismes d'espaces algébriques.
La définition suivante est motivée par
la proposition que nous venons de citer.

\begin{definition}
\label{definition-locally-finite-presentation}
Soit $S$ un schéma.
\begin{enumerate}
\item Un foncteur $F : (\Sch/S)_{fppf}^{opp} \to \textit{Ens}$
est dit {\it compatible aux limites} ou {\it localement de présentation finie} si
pour tout schéma affine $T$ sur $S$ qui est une limite $T = \lim T_i$
d'un système projectif filtrant de schémas affines $T_i$ sur $S$, on a
$$
F(T) = \colim F(T_i).
$$
Nous disons parfois que $F$ est {\it localement de présentation finie sur $S$}.
\item Soient $F, G : (\Sch/S)_{fppf}^{opp} \to \textit{Ens}$.
Une transformation de foncteurs $a : F \to G$
est {\it compatible aux limites} ou {\it localement de présentation finie}
si, pour tout schéma $T$ sur $S$ et tout $y \in G(T)$, le foncteur
$$
F_y : (\Sch/T)_{fppf}^{opp} \longrightarrow \textit{Ens}, \quad
T'/T \longmapsto \{x \in F(T') \mid a(x) = y|_{T'}\}
$$
est localement de présentation finie sur $T$\footnote{La caractérisation (2) du
lemme \ref{lemma-characterize-relative-limit-preserving}
est peut-être plus facile à lire.}. Nous disons parfois que
$F$ est {\it compatible aux limites relativement à $G$}.
\end{enumerate}
\end{definition}

\noindent
Le foncteur $F_y$ est en un certain sens la fibre de
$a : F \to G$ au-dessus de $y$, sauf qu'il s'agit d'un préfaisceau sur le grand site fppf
de $T$. Une formule pour ce foncteur est la suivante :
\begin{equation}
\label{equation-fibre-map-functors}
F_y =
F|_{(\Sch/T)_{fppf}}
{\times}_{G|_{(\Sch/T)_{fppf}}}
*
\end{equation}
Ici, $*$ est l'objet final de la catégorie des (pré)faisceaux
sur $(\Sch/T)_{fppf}$ (voir
Sites, exemple \ref{sites-example-singleton-sheaf})
et le morphisme $* \to G|_{(\Sch/T)_{fppf}}$ est donné par $y$.
Notons que si $j : (\Sch/T)_{fppf} \to (\Sch/S)_{fppf}$
est le foncteur de localisation, alors la formule ci-dessus devient
$F_y = j^{-1}F \times_{j^{-1}G} *$ et $j_!F_y$ est simplement le produit fibré
$F \times_{G, y} T$. (Voir
Sites, section \ref{sites-section-localize},
pour des informations sur la localisation, et en particulier
Sites, remarque \ref{sites-remark-localize-presheaves}
pour des informations sur $j_!$ pour les préfaisceaux.)

\medskip\noindent
À ce stade, nous disposons temporairement de deux définitions de ce que signifie,
pour un morphisme $X \to Y$ d'espaces algébriques sur $S$, être localement de présentation
finie. L'une est donnée par
Morphismes d'espaces,
définition \ref{spaces-morphisms-definition-locally-finite-presentation}
et l'autre utilise le fait que $X \to Y$ est une transformation de foncteurs, de sorte que
la définition \ref{definition-locally-finite-presentation}
s'applique (nous emploierons autant que possible la terminologie « compatible aux limites »
pour cette notion). Nous montrerons dans
la proposition \ref{proposition-characterize-locally-finite-presentation}
que ces deux définitions coïncident.

\begin{lemma}
\label{lemma-characterize-relative-limit-preserving}
Soit $S$ un schéma. Soit $a : F \to G$ une transformation de foncteurs
$(\Sch/S)_{fppf}^{opp} \to \textit{Ens}$.
Les conditions suivantes sont équivalentes :
\begin{enumerate}
\item $a : F \to G$ est compatible aux limites ;
\item pour tout schéma affine $T$ sur $S$ qui est une
limite $T = \lim T_i$ d'un système projectif filtrant de
schémas affines $T_i$ sur $S$, le diagramme d'ensembles
$$
\xymatrix{
\colim_i F(T_i) \ar[r] \ar[d]_a & F(T) \ar[d]^a \\
\colim_i G(T_i) \ar[r] & G(T)
}
$$
est cartésien.
\end{enumerate}
\end{lemma}

\begin{proof}
Supposons (1). Considérons $T = \lim_{i \in I} T_i$ comme dans (2). Soit
$(y, x_T)$ un élément du produit fibré
$\colim_i G(T_i) \times_{G(T)} F(T)$.
Alors $y$ provient d'un $y_i \in G(T_i)$ pour un certain $i$.
Considérons le foncteur $F_{y_i}$ sur $(\Sch/T_i)_{fppf}$ comme dans
la définition \ref{definition-locally-finite-presentation}.
On voit que $x_T \in F_{y_i}(T)$. De plus, $T = \lim_{i' \geq i} T_{i'}$
est un système projectif filtrant de schémas affines sur $T_i$. Ainsi (1) implique
que $x_T$ est l'image d'un unique élément $x$ de
$\colim_{i' \geq i} F_{y_i}(T_{i'})$. Ainsi $x$ est l'unique
élément de $\colim F(T_i)$ ayant pour image le couple $(y, x_T)$.
Ceci prouve (2).

\medskip\noindent
Supposons (2). Soit $T$ un schéma et $y_T \in G(T)$. Nous devons montrer que
$F_{y_T}$ est compatible aux limites. Soit $T' = \lim_{i \in I} T'_i$ un
schéma affine sur $T$, limite d'un système projectif filtrant de schémas affines $T'_i$
sur $T$. Soit $x_{T'} \in F_{y_T}$. Choisissons $i \in I$, ce qui est possible car
$I$ est un ensemble ordonné filtrant. Notons $y_i \in F(T'_i)$ l'image
de $y_{T'}$. Alors $(y_i, x_{T'})$ est un
élément du produit fibré
$\colim_i G(T'_i) \times_{G(T')} F(T')$.
Par (2), il existe donc un unique élément $x$ de $\colim_i F(T'_i)$
ayant pour image $(y_i, x_{T'})$. Il est clair que $x$ définit un élément
de $\colim_i F_y(T'_i)$ ayant pour image $x_{T'}$, ce qui conclut.
\end{proof}

\begin{lemma}
\label{lemma-composition-locally-finite-presentation}
Soit $S$ un schéma contenu dans $\Sch_{fppf}$.
Soient $F, G, H : (\Sch/S)_{fppf}^{opp} \to \textit{Ens}$.
Soient $a : F \to G$, $b : G \to H$ des transformations de foncteurs.
Si $a$ et $b$ sont compatibles aux limites, alors
$$
b \circ a : F \longrightarrow H
$$
est compatible aux limites.
\end{lemma}

\begin{proof}
Soit $T = \lim_{i \in I} T_i$ comme dans la caractérisation (2) du
lemme \ref{lemma-characterize-relative-limit-preserving}.
Considérons le diagramme d'ensembles
$$
\xymatrix{
\colim_i F(T_i) \ar[r] \ar[d]_a & F(T) \ar[d]^a \\
\colim_i G(T_i) \ar[r] \ar[d]_b & G(T) \ar[d]^b \\
\colim_i H(T_i) \ar[r] & H(T)
}
$$
Par hypothèse, les deux carrés sont cartésiens. Le
rectangle extérieur est donc lui aussi cartésien, ce qui prouve le lemme.
\end{proof}

\begin{lemma}
\label{lemma-locally-finite-presentation-permanence}
Soit $S$ un schéma contenu dans $\Sch_{fppf}$.
Soient $F, G, H : (\Sch/S)_{fppf}^{opp} \to \textit{Ens}$.
Soient $a : F \to G$, $b : G \to H$ des transformations de foncteurs.
Si $b \circ a$ et $b$ sont compatibles aux limites, alors $a$
est compatible aux limites.
\end{lemma}

\begin{proof}
Soit $T = \lim_{i \in I} T_i$ comme dans la caractérisation (2) du
lemme \ref{lemma-characterize-relative-limit-preserving}.
Considérons le diagramme d'ensembles
$$
\xymatrix{
\colim_i F(T_i) \ar[r] \ar[d]_a & F(T) \ar[d]^a \\
\colim_i G(T_i) \ar[r] \ar[d]_b & G(T) \ar[d]^b \\
\colim_i H(T_i) \ar[r] & H(T)
}
$$
Par hypothèse, le carré inférieur et le rectangle extérieur
sont cartésiens. Le carré supérieur
est donc lui aussi cartésien, ce qui prouve le lemme.
\end{proof}

\begin{lemma}
\label{lemma-base-change-locally-finite-presentation}
Soit $S$ un schéma contenu dans $\Sch_{fppf}$.
Soient $F, G, H : (\Sch/S)_{fppf}^{opp} \to \textit{Ens}$.
Soient $a : F \to G$, $b : H \to G$ des transformations de foncteurs.
Considérons le diagramme cartésien
$$
\xymatrix{
H \times_{b, G, a} F \ar[r]_-{b'} \ar[d]_{a'} & F \ar[d]^a \\
H \ar[r]^b & G
}
$$
Si $a$ est compatible aux limites, alors le changement de base $a'$ l'est aussi.
\end{lemma}

\begin{proof}
Preuve omise. Indication : c'est formel.
\end{proof}

\begin{lemma}
\label{lemma-fibre-product-locally-finite-presentation}
Soit $S$ un schéma contenu dans $\Sch_{fppf}$.
Soient $E, F, G, H : (\Sch/S)_{fppf}^{opp} \to \textit{Ens}$.
Soient $a : F \to G$, $b : H \to G$ et $c : G \to E$
des transformations de foncteurs. Si $c$, $c \circ a$ et $c \circ b$
sont compatibles aux limites, alors $F \times_G H \to E$ l'est aussi.
\end{lemma}

\begin{proof}
Soit $T = \lim_{i \in I} T_i$ comme dans la caractérisation (2) du
lemme \ref{lemma-characterize-relative-limit-preserving}.
Nous avons alors
$$
\colim (F \times_G H)(T_i) =
\colim F(T_i) \times_{\colim G(T_i)} \colim H(T_i)
$$
car les limites inductives filtrantes commutent aux produits finis. Notre but est donc de
montrer que
$$
\xymatrix{
\colim F(T_i) \times_{\colim G(T_i)} \colim H(T_i) \ar[r] \ar[d] &
F(T) \times_{G(T)} H(T) \ar[d] \\
\colim_i E(T_i) \ar[r] & E(T)
}
$$
est cartésien. Cela découle de l'observation que,
pour des applications d'ensembles $E' \to E$, $F \to G$, $H \to G$ et $G \to E$,
nous avons
$$
E' \times_E (F \times_G H) =
(E' \times_E F) \times_{(E' \times_E G)} (E' \times_E H)
$$
Quelques détails sont omis.
\end{proof}

\begin{lemma}
\label{lemma-sheafify-finite-presentation}
Soit $S$ un schéma contenu dans $\Sch_{fppf}$.
Soit $F : (\Sch/S)_{fppf}^{opp} \to \textit{Ens}$ un foncteur.
Si $F$ est compatible aux limites, alors le faisceau associé $F^\#$ l'est aussi.
\end{lemma}

\begin{proof}
Supposons que $F$ soit compatible aux limites.
Il suffit de montrer que $F^+$ est compatible aux limites, puisque
$F^\# = (F^+)^+$, voir
Sites, théorème \ref{sites-theorem-plus}.
Soit $T$ un schéma affine sur $S$, et écrivons $T = \lim T_i$
comme limite d'un système projectif filtrant de schémas affines sur $S$.
Rappelons que $F^+(T)$ est la limite inductive des $\check H^0(\mathcal{V}, F)$
où cette limite inductive porte sur tous les recouvrements de $T$ dans $(\Sch/S)_{fppf}$.
Tout recouvrement fppf d'un schéma affine admet un raffinement qui est un recouvrement
fppf standard, voir
Topologies, lemme \ref{topologies-lemma-fppf-affine}.
Nous pouvons donc écrire
$$
F^+(T)
=
\colim_{\mathcal{V}\text{ recouvrement standard de }T}
\check H^0(\mathcal{V}, F).
$$
Tout recouvrement $\mathcal{V} = \{T_k \to T\}_{k = 1, \ldots, n}$
intervenant dans cette limite inductive peut s'écrire
$\mathcal{V}_i \times_{T_i} T$ pour un certain $i$ et un recouvrement fppf standard
$\mathcal{V}_i = \{T_{i, k} \to T_i\}_{k = 1, \ldots, n}$ de $T_i$.
Notons $\mathcal{V}_{i'} = \{T_{i', k} \to T_{i'}\}_{k = 1, \ldots, n}$
le recouvrement obtenu par changement de base pour $i' \geq i$. On a alors
\begin{align*}
\colim_{i' \geq i} \check H^0(\mathcal{V}_i, F)
& =
\colim_{i' \geq i}
\text{Noyau du couple}
\left(
\xymatrix{
\prod F(T_{i', k})
\ar@<1ex>[r] \ar@<-1ex>[r] &
\prod F(T_{i', k} \times_{T_{i'}} T_{i', l})
}
\right)
\\
& =
\text{Noyau du couple}
\left(
\xymatrix{
\colim_{i' \geq i}
\prod F(T_{i', k})
\ar@<1ex>[r] \ar@<-1ex>[r] &
\colim_{k' \geq k}
\prod F(T_{i', k} \times_{T_{i'}} T_{i', l})
}
\right)
\\
& =
\text{Noyau du couple}
\left(
\xymatrix{
\prod F(T_k)
\ar@<1ex>[r] \ar@<-1ex>[r] &
\prod F(T_k \times_T T_l)
}
\right)
\\
& =
\check H^0(\mathcal{V}, F)
\end{align*}
La deuxième égalité résulte de l'exactitude des limites inductives filtrantes.
La troisième résulte de ce que $F$ est compatible aux limites et de ce que
$\lim_{i' \geq i} T_{i', k} = T_k$ et
$\lim_{i' \geq i} T_{i', k} \times_{T_{i'}} T_{i', l} = T_k \times_T T_l$
d'après Limites, lemme \ref{limits-lemma-scheme-over-limit}.
En utilisant ceci simultanément pour tous les recouvrements, nous obtenons
\begin{align*}
F^+(T)
& =
\colim_{\mathcal{V}\text{ recouvrement standard de }T} \check H^0(\mathcal{V}, F) \\
& =
\colim_{i \in I}
\colim_{\mathcal{V}_i\text{ recouvrement standard de }T_i}
\check H^0(T \times_{T_i}\mathcal{V}_i, F) \\
& =
\colim_{i \in I} F^+(T_i)
\end{align*}
On peut permuter l'ordre des limites inductives en vertu de
Catégories, lemme \ref{categories-lemma-colimits-commute}.
\end{proof}

\begin{lemma}
\label{lemma-sheaf-finite-presentation}
Soit $S$ un schéma.
Soit $F : (\Sch/S)_{fppf}^{opp} \to \textit{Ens}$ un foncteur.
Supposons que
\begin{enumerate}
\item $F$ soit un faisceau, et
\item il existe un recouvrement fppf $\{U_j \to S\}_{j \in J}$ tel que
$F|_{(\Sch/U_j)_{fppf}}$ soit compatible aux limites.
\end{enumerate}
Alors $F$ est compatible aux limites.
\end{lemma}

\begin{proof}
Soit $T$ un schéma affine sur $S$.
Soit $I$ un ensemble ordonné filtrant et soit
$T_i$ un système projectif de schémas affines sur $S$ tel que
$T = \lim T_i$. Il s'agit de montrer que l'application canonique
$\colim F(T_i) \to F(T)$ est bijective.

\medskip\noindent
Choisissons un élément $0 \in I$ et un recouvrement fppf standard
$\{V_{0, k} \to T_{0}\}_{k = 1, \ldots, m}$ qui raffine
le recouvrement $\{U_j \times_S T_0 \to T_0\}$ obtenu par changement de base à partir du recouvrement donné de $S$.
Pour tout $i \geq 0$, posons $V_{i, k} = T_i \times_{T_0} V_{0, k}$, et
posons $V_k = T \times_{T_0} V_{0, k}$. On a
$V_k = \lim_{i \geq 0} V_{i, k}$, voir
Limites, lemme \ref{limits-lemma-scheme-over-limit}.

\medskip\noindent
Supposons que $x, x' \in \colim F(T_i)$ aient la même image dans
$F(T)$. Représentons $x, x'$ par des éléments $x_i, x'_i \in F(T_i)$
pour un certain $i \in I$ (on peut choisir le même $i$ pour les deux puisque $I$ est filtrant).
L'hypothèse (2) et l'égalité des images de $x_i, x'_i$ dans
$F(T)$ montrent que
$$
x_i|_{V_{i', k}} = x'_i|_{V_{i', k}}
$$
pour un certain $i' \in I$ suffisamment grand. On peut choisir le même $i'$ pour tout
$k$, car $k \in \{1, \ldots, m\}$ parcourt un ensemble fini.
Puisque $\{V_{i', k} \to T_{i'}\}$
est un recouvrement fppf et que $F$ est un faisceau, il vient
$x_i|_{T_{i'}} = x'_i|_{T_{i'}}$. L'application
$\colim F(T_i) \to F(T)$ est injective.

\medskip\noindent
La surjectivité se démontre de manière analogue.
Soit $x \in F(T)$. Pour chaque $k$, l'hypothèse (2) permet de
choisir un indice $i$ tel que $x|_{V_k}$ provienne d'un
élément $x_{i, k} \in F(V_{i, k})$. Comme précédemment, on peut choisir un
même $i$ qui convienne pour tous les $k$. D'après l'injectivité
établie ci-dessus, on a
$$
x_{i, k}|_{V_{i', k} \times_{T_{i'}} V_{i', l}}
=
x_{i, l}|_{V_{i', k} \times_{T_{i'}} V_{i', l}}
$$
pour un $i'$ suffisamment grand. La condition de faisceau pour $F$ implique alors que
les éléments $x_{i, k}|_{V_{i', k}}$ se recollent en un élément $x_{i'} \in F(T_{i'})$,
ce qui conclut.
\end{proof}

\begin{lemma}
\label{lemma-sheafify-finite-presentation-map}
Soit $S$ un schéma contenu dans $\Sch_{fppf}$.
Soient $F, G : (\Sch/S)_{fppf}^{opp} \to \textit{Ens}$ des foncteurs.
Si la transformation $a : F \to G$ est compatible aux limites,
il en est de même de la transformation induite entre les faisceaux associés
$F^\# \to G^\#$.
\end{lemma}

\begin{proof}
Soit $T$ un schéma et soit $y \in G^\#(T)$.
Il s'agit de montrer que le foncteur
$F^\#_y : (\Sch/T)_{fppf}^{opp} \to \textit{Ens}$
construit à partir de $F^\# \to G^\#$ et de $y$ comme dans
la définition \ref{definition-locally-finite-presentation}
est compatible aux limites.
D'après l'équation (\ref{equation-fibre-map-functors}),
$F^\#_y$ est un faisceau. Choisissons un recouvrement fppf
$\{V_j \to T\}_{j \in J}$ tel que $y|_{V_j}$ provienne
d'un élément $y_j \in F(V_j)$.
La restriction de $F^\#$ à $(\Sch/V_j)_{fppf}$
n'est autre que $F^\#_{y_j}$. Si nous montrons que $F^\#_{y_j}$
est compatible aux limites, alors
le lemme \ref{lemma-sheaf-finite-presentation}
garantit que $F^\#_y$ est compatible aux limites,
ce qui conclut. On est donc ramené au cas $y \in G(T)$.

\medskip\noindent
Soit $y \in G(T)$. Dans ce cas, on a $F^\#_y = (F_y)^\#$.
Cette égalité résulte de
l'équation (\ref{equation-fibre-map-functors}).
Le résultat découle donc du
lemme \ref{lemma-sheafify-finite-presentation}.
\end{proof}

\begin{proposition}
\label{proposition-characterize-locally-finite-presentation}
Soit $S$ un schéma. Soit $f : X \to Y$ un morphisme d'espaces algébriques
sur $S$. Les conditions suivantes sont équivalentes :
\begin{enumerate}
\item Le morphisme $f$ est localement de présentation finie en tant que
morphisme d'espaces algébriques, voir
Morphismes d'espaces,
définition \ref{spaces-morphisms-definition-locally-finite-presentation}.
\item Le morphisme $f : X \to Y$ est compatible aux limites en tant que
transformation de foncteurs, voir
définition \ref{definition-locally-finite-presentation}.
\end{enumerate}
\end{proposition}

\begin{proof}
Supposons (1). Soit $T$ un schéma et soit $y \in Y(T)$. Nous devons montrer que
$T \times_Y X$ est compatible aux limites sur $T$ au sens de
la définition \ref{definition-locally-finite-presentation}.
Nous sommes donc réduits à démontrer que si $X$ est un espace algébrique qui
est localement de présentation finie sur $S$ en tant qu'espace algébrique, alors il
est compatible aux limites en tant que foncteur
$X : (\Sch/S)_{fppf}^{opp} \to \textit{Ens}$.
Pour le voir, choisissons une présentation $X = U/R$, voir
Espaces, définition \ref{spaces-definition-presentation}.
Il découle de
Morphismes d'espaces,
définition \ref{spaces-morphisms-definition-locally-finite-presentation}
que $U$ et $R$ sont des schémas localement de présentation finie
sur $S$. Donc, par
Limites, proposition
\ref{limits-proposition-characterize-locally-finite-presentation},
nous avons
$$
U(T) = \colim U(T_i), \quad
R(T) = \colim R(T_i)
$$
chaque fois que $T = \lim_i T_i$ dans $(\Sch/S)_{fppf}$. Il s'ensuit
que le préfaisceau
$$
(\Sch/S)_{fppf}^{opp} \longrightarrow \textit{Ens}, \quad
W \longmapsto U(W)/R(W)
$$
est compatible aux limites. Donc, par le
lemme \ref{lemma-sheafify-finite-presentation},
le faisceau associé $X = U/R$ est lui aussi compatible aux limites.

\medskip\noindent
Supposons (2). Choisissons un schéma $V$ et un morphisme \'etale surjectif
$V \to Y$. Ensuite, choisissons un schéma $U$ et un morphisme \'etale surjectif
$U \to V \times_Y X$. Par le
lemme \ref{lemma-base-change-locally-finite-presentation}, la transformation de foncteurs
$V \times_Y X \to V$ est compatible aux limites. Par
Morphismes d'espaces,
lemme \ref{spaces-morphisms-lemma-etale-locally-finite-presentation},
le morphisme d'espaces algébriques $U \to V \times_Y X$ est localement
de présentation finie, donc est compatible aux limites en tant que
transformation de foncteurs par la première partie de la preuve. Par
lemme \ref{lemma-composition-locally-finite-presentation},
la composition $U \to V \times_Y X \to V$ est compatible aux limites
en tant que transformation de foncteurs. Donc
le morphisme de schémas $U \to V$ est localement de présentation finie par
Limites, proposition
\ref{limits-proposition-characterize-locally-finite-presentation}
(à ceci près d'une remarque ensembliste, voir le dernier paragraphe de la preuve).
Cela signifie, par définition, que la condition (1) est satisfaite.

\medskip\noindent
Remarque ensembliste. Soit $U \to V$ un morphisme de
$(\Sch/S)_{fppf}$. Dans l'énoncé de
Limites, proposition
\ref{limits-proposition-characterize-locally-finite-presentation}
le fait que $U \to V$ soit localement de présentation finie
se caractérise ainsi : pour {\it tout} système projectif filtrant $(T_i, f_{ii'})$ de schémas affines
sur $V$, on a $U(T) = \colim V(T_i)$ ; mais, dans le contexte actuel,
nous ne pouvons considérer que des schémas affines $T_i$ sur $V$ qui sont (isomorphes à)
un objet de $(\Sch/S)_{fppf}$. Nous devons donc nous assurer qu'il y a suffisamment
de schémas affines dans $(\Sch/S)_{fppf}$ pour que la preuve fonctionne.
En examinant la preuve de (2) $\Rightarrow$ (1) de
Limites, proposition
\ref{limits-proposition-characterize-locally-finite-presentation}
nous voyons que la question se réduit au cas où $U$ et $V$ sont affines.
Écrivons $U = \Spec(A)$ et $V = \Spec(B)$. Par construction
de $(\Sch/S)_{fppf}$, le spectre de tout anneau de cardinal
$\leq |B|$ est isomorphe à un objet de $(\Sch/S)_{fppf}$.
Il suffit donc d'observer que, dans la démonstration du sens nécessaire de
Algèbre, lemme \ref{algebra-lemma-characterize-finite-presentation},
seules des $A$-algèbres de cardinal $\leq |B|$ sont utilisées.
\end{proof}

\begin{remark}
\label{remark-limit-preserving}
Voici un cas particulier important de
proposition \ref{proposition-characterize-locally-finite-presentation}.
Soit $S$ un schéma. Soit $X$ un espace algébrique sur $S$.
Alors $X$ est localement de présentation finie sur $S$ si et seulement
si $X$, en tant que foncteur $(\Sch/S)^{opp} \to \textit{Ens}$,
est compatible aux limites. Comparer avec
Limites, remarque \ref{limits-remark-limit-preserving}.
En fait, nous verrons dans le lemme \ref{lemma-surjection-is-enough}
ci-dessous qu'il suffit que l'application
$$
\colim X(T_i) \longrightarrow X(T)
$$
soit surjective lorsque $T = \lim T_i$ est la limite d'un système projectif filtrant de
schémas affines sur $S$.
\end{remark}

\begin{lemma}
\label{lemma-surjection-is-enough}
Soit $S$ un schéma.
Soit $f : X \to Y$ un morphisme d'espaces algébriques sur $S$.
Si, pour toute limite $T = \lim_{i \in I} T_i$ d'un système projectif filtrant
de schémas affines sur $S$, l'application
$$
\colim X(T_i) \longrightarrow X(T) \times_{Y(T)} \colim Y(T_i)
$$
est surjective, alors $f$ est localement de présentation finie.
Autrement dit, dans la partie (2) de la
proposition \ref{proposition-characterize-locally-finite-presentation},
il suffit de vérifier la surjectivité dans le critère du
lemme \ref{lemma-characterize-relative-limit-preserving}.
\end{lemma}

\begin{proof}
Choisissons un schéma $V$ et un morphisme \'etale surjectif $g : V \to Y$.
Ensuite, choisissons un schéma $U$ et un morphisme \'etale surjectif
$h : U \to V \times_Y X$. Il suffit de montrer, pour $T = \lim T_i$
comme dans le lemme, que l'application
$$
\colim U(T_i) \longrightarrow U(T) \times_{V(T)} \colim V(T_i)
$$
est surjective, car alors $U \to V$ sera localement de présentation
finie par Limites, lemme \ref{limits-lemma-surjection-is-enough}
(à ceci près d'une remarque ensembliste exactement comme dans la preuve de
la proposition \ref{proposition-characterize-locally-finite-presentation}).
Ainsi, prenons $a : T \to U$ et $b_i : T_i \to V$ qui déterminent
le même morphisme $T \to V$. Voici le diagramme
$$
\xymatrix{
T \ar[d]_a \ar[rr]_{p_i} & & T_i \ar[d]^{b_i} \ar@{..>}[ld] \\
U \ar[r]^-h & X \times_Y V \ar[d] \ar[r] & V \ar[d]^g \\
& X \ar[r]^f & Y
}
$$
D'après l'hypothèse du lemme, quitte à augmenter $i$, il existe
un morphisme $c_i : T_i \to X$ tel que
$h \circ a = (b_i, c_i) \circ p_i : T_i \to V \times_Y X$
et tel que $f \circ c_i = g \circ b_i$.
Puisque $h$ est un morphisme \'etale d'espaces algébriques
(et donc localement de présentation finie), nous avons la surjectivité de
$$
\colim U(T_i) \longrightarrow U(T) \times_{(X \times_Y V)(T)}
\colim (X \times_Y V)(T_i)
$$
par la proposition \ref{proposition-characterize-locally-finite-presentation}.
Par conséquent, quitte à augmenter encore $i$, il existe un
morphisme $a_i : T_i \to U$ tel que $a = a_i \circ p_i$ et
$b_i = (U \to V) \circ a_i$.
\end{proof}















\section{Limites d'espaces algébriques}
\label{section-limits}

\noindent
Le lemme suivant explique comment nous concevons les limites des espaces
algébriques dans ce chapitre. Nous utiliserons (sans autre mention) que le
changement de base d'un morphisme affine d'espaces algébriques est affine (voir
Morphismes d'espaces, lemme \ref{spaces-morphisms-lemma-base-change-affine}).

\begin{lemma}
\label{lemma-directed-inverse-system-has-limit}
Soit $S$ un schéma. Soit $I$ un ensemble ordonné filtrant.
Soit $(X_i, f_{ii'})$ un système projectif indexé par $I$
dans la catégorie des espaces algébriques sur $S$.
Si les morphismes $f_{ii'} : X_i \to X_{i'}$ sont affines, alors la
limite $X = \lim_i X_i$ (en tant que faisceau fppf) est un espace algébrique.
De plus,
\begin{enumerate}
\item chacun des morphismes $f_i : X \to X_i$ est affine,
\item pour tout $i \in I$ et tout morphisme d'espaces algébriques
$T \to X_i$, nous avons
$$
X \times_{X_i} T = \lim_{i' \geq i} X_{i'} \times_{X_i} T.
$$
Cette égalité étant entendue dans la catégorie des espaces algébriques sur $S$.
\end{enumerate}
\end{lemma}

\begin{proof}
La partie (2) est une conséquence formelle de l'existence de la
limite $X = \lim X_i$ en tant qu'espace algébrique sur $S$.
Choisissons un élément $0 \in I$ (c'est possible puisque tout ensemble filtrant est non vide).
Choisissons un schéma $U_0$ et un morphisme
\'etale surjectif $U_0 \to X_0$. Posons $R_0 = U_0 \times_{X_0} U_0$
de sorte que $X_0 = U_0/R_0$. Pour $i \geq 0$, posons
$U_i = X_i \times_{X_0} U_0$ et
$R_i = X_i \times_{X_0} R_0 = U_i \times_{X_i} U_i$.
D'après Limites, lemme \ref{limits-lemma-directed-inverse-system-has-limit}
on voit que $U = \lim_{i \geq 0} U_i$ et $R = \lim_{i \geq 0} R_i$
sont des schémas. De plus, les deux morphismes $s, t : R \to U$ sont les
changements de base des deux projections $R_0 \to U_0$ par le morphisme
$U \to U_0$, en particulier ils sont \'etales. Le morphisme $R \to U \times_S U$
définit une relation d'équivalence, car une limite projective de relations d'équivalence
indexée par un ensemble filtrant est une relation d'équivalence. Par conséquent, le morphisme
$R \to U \times_S U$ est une relation d'équivalence \'etale. Nous affirmons que
le morphisme naturel
\begin{equation}
\label{equation-isomorphism-sheaves}
U/R \longrightarrow \lim X_i
\end{equation}
est un isomorphisme de faisceaux fppf sur la catégorie des schémas sur $S$.
Cette assertion implique que $X = \lim X_i$ est un espace algébrique
d'après Espaces, théorème \ref{spaces-theorem-presentation}.

\medskip\noindent
Soit $Z$ un schéma et soit $a : Z \to \lim X_i$ un morphisme.
Alors $a = (a_i)$, où $a_i : Z \to X_i$. Posons $W_0 = Z \times_{a_0, X_0} U_0$.
Remarquons que $W_0 = Z \times_{a_i, X_i} U_i$ pour tout $i \geq 0$, par notre
choix de $U_i \to X_i$ ci-dessus. Nous obtenons donc un morphisme
$W_0 \to \lim_{i \geq 0} U_i = U$. Comme $W_0 \to Z$ est surjectif
et \'etale, nous concluons que (\ref{equation-isomorphism-sheaves})
est un morphisme surjectif de faisceaux. Enfin, soit
$Z$ un schéma et soient $a, b : Z \to U/R$ deux morphismes
dont les images par (\ref{equation-isomorphism-sheaves}) sont égales.
Il s'agit de montrer que $a = b$.
Après avoir remplacé $Z$ par les objets d'un recouvrement fppf,
on peut supposer qu'il existe des morphismes $a', b' : Z \to U$ qui
donnent $a$ et $b$. Le fait que les images de $a, b$ par
(\ref{equation-isomorphism-sheaves}) soient égales signifie que
pour chaque $i \geq 0$ les compositions $a_i', b_i' : Z \to U \to U_i$
sont égales comme morphismes à valeurs dans $U_i/R_i = X_i$. Par conséquent,
$(a_i', b_i') : Z \to U_i \times_S U_i$ se factorise par
$R_i$ au moyen d'un morphisme $c_i : Z \to R_i$. Comme
$R = \lim_{i \geq 0} R_i$, on voit que $c = \lim c_i : Z \to R$
est un morphisme qui montre que $a, b$ sont égaux comme morphismes
de $Z$ vers $U/R$.

\medskip\noindent
La partie (1) découle du fait que nous avons vu ci-dessus
$U_i \times_{X_i} X = U$ et que $U \to U_i$ est affine par
construction.
\end{proof}

\begin{lemma}
\label{lemma-space-over-limit}
Soit $S$ un schéma. Soit $I$ un ensemble ordonné filtrant.
Soit $(X_i, f_{ii'})$ un système projectif indexé par $I$ d'espaces algébriques
sur $S$ à morphismes de transition affines.
Soit $X = \lim_i X_i$. Soit $0 \in I$. Supposons que $T \to X_0$ soit un
morphisme d'espaces algébriques. Alors
$$
T \times_{X_0} X = \lim_{i \geq 0} T \times_{X_0} X_i
$$
en tant qu'espaces algébriques sur $S$.
\end{lemma}

\begin{proof}
La limite $X$ est un espace algébrique d'après
le lemme \ref{lemma-directed-inverse-system-has-limit}.
L'égalité est formelle, voir
Catégories, lemme \ref{categories-lemma-colimits-commute}.
\end{proof}

\begin{lemma}
\label{lemma-directed-inverse-system-closed-immersions}
Soit $S$ un schéma. Soit $I$ un ensemble ordonné filtrant.
Soit $(X_i, f_{i'i}) \to (Y_i, g_{i'i})$ un morphisme
de systèmes projectifs indexés par $I$ d'espaces algébriques sur $S$.
Supposons que
\begin{enumerate}
\item les morphismes $f_{i'i} : X_{i'} \to X_i$ soient affines,
\item les morphismes $g_{i'i} : Y_{i'} \to Y_i$ soient affines,
\item les morphismes $X_i \to Y_i$ soient des immersions fermées.
\end{enumerate}
Alors $\lim X_i \to \lim Y_i$ est une immersion fermée.
\end{lemma}

\begin{proof}
Observons que $\lim X_i$ et $\lim Y_i$ existent d'après
le lemme \ref{lemma-directed-inverse-system-has-limit}.
Choisissons $0 \in I$ et un schéma affine $V_0$ ainsi qu'un morphisme \'etale
$V_0 \to Y_0$. Alors les morphismes
$V_i = Y_i \times_{Y_0} V_0 \to U_i = X_i \times_{Y_0} V_0$
sont des immersions fermées de schémas affines.
Par conséquent, le morphisme $V = Y \times_{Y_0} V_0 \to U = X \times_{Y_0} V_0$
est une immersion fermée, car $V = \lim V_i$, $U = \lim U_i$
et parce qu'une limite projective d'immersions fermées de schémas affines est une
immersion fermée : une limite inductive filtrante d'homomorphismes d'anneaux surjectifs
est surjective. Comme les morphismes \'etales $V \to Y$ forment un
recouvrement \'etale de $Y$ lorsque nous faisons varier notre choix de $V_0 \to Y_0$
le lemme s'ensuit.
\end{proof}

\begin{lemma}
\label{lemma-directed-inverse-system-reduced}
Soit $S$ un schéma. Soit $I$ un ensemble ordonné filtrant.
Soit $(X_i, f_{i'i})$ un système projectif indexé par $I$
d'espaces algébriques sur $S$. Si $X_i$ est réduit
pour tout $i$, alors $X$ est réduit.
\end{lemma}

\begin{proof}
Observons que $\lim X_i$ existe d'après
le lemme \ref{lemma-directed-inverse-system-has-limit}.
Choisissons $0 \in I$ et un schéma affine $V_0$ ainsi qu'un morphisme \'etale
$U_0 \to X_0$. Alors les schémas affines
$U_i = X_i \times_{X_0} U_0$ sont réduits.
Par conséquent, $U = X \times_{X_0} U_0$
est un schéma affine réduit, comme limite projective de schémas affines réduits :
une limite inductive filtrante d'anneaux réduits est réduite.
Comme les morphismes \'etales $U \to X$ forment un
recouvrement \'etale de $X$ lorsque nous faisons varier notre choix de $U_0 \to X_0$
le lemme s'ensuit.
\end{proof}

\begin{lemma}
\label{lemma-better-characterize-relative-limit-preserving}
Soit $S$ un schéma. Soit $X \to Y$ un morphisme d'espaces algébriques
sur $S$. Les conditions équivalentes (1) et (2) de la
proposition \ref{proposition-characterize-locally-finite-presentation}
sont également équivalentes à
\begin{enumerate}
\item[(3)] pour toute limite $T = \lim T_i$ d'un système projectif filtrant
d'espaces algébriques quasi-compacts et quasi-séparés $T_i$ sur $S$ à morphismes de transition affines, le diagramme
d'ensembles
$$
\xymatrix{
\colim_i \Mor(T_i, X) \ar[r] \ar[d] & \Mor(T, X) \ar[d] \\
\colim_i \Mor(T_i, Y) \ar[r] & \Mor(T, Y)
}
$$
est un diagramme cartésien.
\end{enumerate}
\end{lemma}

\begin{proof}
Il est clair que (3) implique (2). Nous supposerons (2) et démontrerons (3).
La démonstration est assez formelle et nous laissons au lecteur le soin
d'en trouver une.

\medskip\noindent
Commençons par démontrer que (3) est valable
lorsque $T_i$ est en outre supposé séparé pour tout $i$.
Choisissons $i \in I$ et un morphisme \'etale surjectif $U_i \to T_i$
où $U_i$ est affine. En utilisant le lemme \ref{lemma-space-over-limit},
on voit qu'en posant $U = U_i \times_{T_i} T$ et
$U_{i'} = U_i \times_{T_i} T_{i'}$, on a $U = \lim_{i' \geq i} U_{i'}$.
Bien sûr, $U$ et $U_{i'}$ sont affines (voir
le lemme \ref{lemma-directed-inverse-system-has-limit}).
Comme $T_i$ est séparé, le produit fibré $V_i = U_i \times_{T_i} U_i$
est également un schéma affine et nous obtenons des schémas affines
$V = V_i \times_{T_i} T$ et
$V_{i'} = V_i \times_{T_i} T_{i'}$ tels que $V = \lim_{i' \geq i} V_{i'}$.
Observons que $U \to T$ et $U_i \to T_i$ sont surjectifs et \'etales
et que $V = U \times_T U$ et $V_{i'} = U_{i'} \times_{T_{i'}} U_{i'}$.
Remarquons que $\Mor(T, X)$ est le noyau du couple formé par les deux applications
$\Mor(U, X) \to \Mor(V, X)$ ; cela résulte par exemple du fait que
$X$, comme faisceau sur $(\Sch/S)_{fppf}$, est le conoyau du couple formé par
les deux applications $h_V \to h_u$. De même,
$\Mor(T_{i'}, X)$ est le noyau du couple formé par les
deux applications $\Mor(U_{i'}, X) \to \Mor(V_{i'}, X)$.
Et bien sûr, la même chose vaut en remplaçant $X$ par $Y$.
La condition (2) dit que les diagrammes de (3) sont cartésiens
dans le cas de $U = \lim U_i$ et $V = \lim V_i$.
Il s'ensuit formellement que la même chose vaut pour $T = \lim T_i$.

\medskip\noindent
Dans le cas général, choisissons un schéma affine $U$, un $i \in I$,
et un morphisme \'etale surjectif $U \to T_i$. En répétant
l'argument du paragraphe précédent, on conclut encore :
les schémas $V_{i'}$, $V$ ne sont plus affines, mais ils sont
toujours quasi-compacts et
séparés, et le résultat du paragraphe précédent s'applique.
\end{proof}




\section{Descente des propriétés}
\label{section-descent}

\noindent
Cette section est l'analogue de Limites, section \ref{limits-section-descent}.

\begin{lemma}
\label{lemma-inverse-limit-sets}
Soit $S$ un schéma. Soit $X = \lim_{i \in I} X_i$ la limite d'un système projectif filtrant
d'espaces algébriques sur $S$ à morphismes de transition affines
(lemme \ref{lemma-directed-inverse-system-has-limit}). Si chaque $X_i$
est décent (par exemple quasi-séparé ou localement séparé),
alors $|X| = \lim_i |X_i|$ en tant qu'ensembles.
\end{lemma}

\begin{proof}
Il existe une application canonique $|X| \to \lim |X_i|$. Choisissons $0 \in I$.
Si $W_0 \subset X_0$ est un sous-espace ouvert, alors
$f_0^{-1}W_0 = \lim_{i \geq 0} f_{i0}^{-1}W_0$, voir
le lemme \ref{lemma-directed-inverse-system-has-limit}.
Ainsi, le cas général se ramène à celui des systèmes projectifs pour lesquels $X_0$
est quasi-compact. On peut donc supposer
désormais que $X_0$ est quasi-compact.

\medskip\noindent
Choisissons un schéma affine $U_0$ et un morphisme \'etale surjectif $U_0 \to X_0$.
Posons $U_i = X_i \times_{X_0} U_0$ et $U = X \times_{X_0} U_0$.
Posons $R_i = U_i \times_{X_i} U_i$ et $R = U \times_X U$.
Rappelons que $U = \lim U_i$ et $R = \lim R_i$, voir la démonstration du
lemme \ref{lemma-directed-inverse-system-has-limit}.
Rappelons que $|X| = |U|/|R|$ et $|X_i| = |U_i|/|R_i|$. D'après
Limites, lemme \ref{limits-lemma-topology-limit}, nous avons
$|U| = \lim |U_i|$ et $|R| = \lim |R_i|$.

\medskip\noindent
Surjectivité de $|X| \to \lim |X_i|$. Soit $(x_i) \in \lim |X_i|$. Notons
$S_i \subset |U_i|$ l'image réciproque de $x_i$. Il s'agit d'un ensemble fini non vide
par définition des espaces décents
(Espaces décents, définition \ref{decent-spaces-definition-very-reasonable}).
Donc $\lim S_i$ est non vide, voir
Catégories, lemme \ref{categories-lemma-nonempty-limit}.
Soit $(u_i) \in \lim S_i \subset \lim |U_i|$. D'après ce qui précède, cela détermine
un point $u \in |U|$ ayant pour image un point $x \in |X|$ dont l'image est l'élément donné
$(x_i)$ de $\lim |X_i|$.

\medskip\noindent
Injectivité de $|X| \to \lim |X_i|$. Supposons que $x, x' \in |X|$
aient la même image dans $\lim |X_i|$. Choisissons des relèvements $u, u' \in |U|$
et notons $u_i, u'_i \in |U_i|$ leurs images.
Pour chaque $i$, soit $T_i \subset |R_i|$ l'ensemble des points ayant pour image
sur $(u_i, u'_i) \in |U_i| \times |U_i|$. Il s'agit d'un ensemble fini
par définition des espaces décents
(Espaces décents, définition \ref{decent-spaces-definition-very-reasonable}).
De plus, $T_i$ est non vide puisque nous avons supposé que $x$ et $x'$ avaient
le même point de $X_i$. Donc $\lim T_i$ est non vide, voir
Catégories, lemme \ref{categories-lemma-nonempty-limit}.
Comme précédemment, soit $r \in |R| = \lim |R_i|$ un point correspondant à un
élément de $\lim T_i$. Alors $r$ a pour image $(u, u')$ dans $|U| \times |U|$
par construction et nous voyons que $x = x'$ dans $|X|$, comme voulu.

\medskip\noindent
Pour justifier l'assertion entre parenthèses, un espace algébrique quasi-séparé est décent, voir
Espaces décents, section \ref{decent-spaces-section-reasonable-decent}
(l'observation clé à cet égard est le lemme de Propriétés des espaces,
\ref{spaces-properties-lemma-finite-fibres-presentation}).
Un espace algébrique localement séparé est décent d'après
Espaces décents, lemme \ref{decent-spaces-lemma-locally-separated-decent}.
\end{proof}

\begin{lemma}
\label{lemma-topology-limit}
Avec les mêmes notations et hypothèses que dans le lemme \ref{lemma-inverse-limit-sets},
nous avons $|X| = \lim_i |X_i|$ en tant qu'espaces topologiques.
\end{lemma}

\begin{proof}
Nous utiliserons le critère de
Topologie, lemme \ref{topology-lemma-characterize-limit}.
Nous avons vu que $|X| = \lim_i |X_i|$ en tant qu'ensembles dans
le lemme \ref{lemma-inverse-limit-sets}.
Les $f_i : X \to X_i$ sont des morphismes d'espaces algébriques
et induisent donc des applications continues $|X| \to |X_i|$.
Ainsi $f_i^{-1}(U_i)$ est ouvert pour tout
ouvert $U_i \subset |X_i|$. Enfin,
soit $x \in |X|$ et soit $x \in V \subset |X|$ un voisinage ouvert.
Nous devons trouver un $i$ et un voisinage ouvert
$W_i \subset |X_i|$ de l'image de $x$
tel que $f_i^{-1}(W_i) \subset V$.
Choisissons $0 \in I$. Choisissons un schéma $U_0$ et un morphisme
\'etale surjectif $U_0 \to X_0$. Posons $U = X \times_{X_0} U_0$
et $U_i = X_i \times_{X_0} U_0$ pour $i \geq 0$.
Alors $U = \lim_{i \geq 0} U_i$ dans la catégorie des schémas d'après
le lemme \ref{lemma-directed-inverse-system-has-limit}.
Choisissons $u \in U$ s'envoyant sur $x$. D'après le résultat pour les schémas
(Limites, lemme \ref{limits-lemma-inverse-limit-top})
nous pouvons trouver un $i \geq 0$ et un voisinage ouvert
$E_i \subset U_i$ de l'image de $u$
dont l'image réciproque dans $U$ est contenue dans
l'image réciproque de $V$ dans $U$. Nous pouvons alors poser $W_i \subset |X_i|$
égal à l'image de $E_i$. Cela fonctionne car $|U_i| \to |X_i|$ est ouvert.
\end{proof}

\begin{lemma}
\label{lemma-limit-nonempty}
Soit $S$ un schéma. Soit $X = \lim_{i \in I} X_i$ la limite d'un système
projectif filtrant d'espaces algébriques sur $S$ à morphismes de transition affines
(lemme \ref{lemma-directed-inverse-system-has-limit}). Si chaque $X_i$
est quasi-compact et non vide, alors $|X|$ est non vide.
\end{lemma}

\begin{proof}
Choisissons $0 \in I$.
Choisissons un schéma affine $U_0$ et un morphisme \'etale surjectif $U_0 \to X_0$.
Posons $U_i = X_i \times_{X_0} U_0$ et $U = X \times_{X_0} U_0$.
Alors chaque $U_i$ est un schéma affine non vide. Donc $U = \lim U_i$
est non vide (Limites, lemme \ref{limits-lemma-limit-nonempty}) et ainsi
$X$ est non vide.
\end{proof}

\begin{lemma}
\label{lemma-inverse-limit-irreducibles}
Soit $S$ un schéma. Soit $X = \lim_{i \in I} X_i$ la limite d'un système
projectif filtrant d'espaces algébriques sur $S$ à morphismes de transition affines
(lemme \ref{lemma-directed-inverse-system-has-limit}).
Soit $x \in |X|$ d'images $x_i \in |X_i|$. Si chaque $X_i$ est décent,
alors $\overline{\{x\}} = \lim_i \overline{\{x_i\}}$ en tant qu'ensembles
et comme espaces algébriques lorsqu'ils sont munis de leur structure induite réduite.
\end{lemma}

\begin{proof}
Posons $Z = \overline{\{x\}} \subset |X|$ et
$Z_i = \overline{\{x_i\}} \subset |X_i|$.
Comme l'application $|X| \to |X_i|$ est continue, l'image de $Z$ est contenue dans $Z_i$
pour tout $i$. On obtient donc une application injective $Z \to \lim Z_i$
car $|X| = \lim |X_i|$ en tant qu'ensembles (lemme \ref{lemma-inverse-limit-sets}).
Supposons que $x' \in |X|$ n'appartienne pas à $Z$.
Il existe alors un ouvert $U \subset |X|$ tel que $x' \in U$
et $x \not \in U$. Comme
$|X| = \lim |X_i|$ en tant qu'espaces topologiques (lemme \ref{lemma-topology-limit}),
on peut écrire $U = \bigcup_{j \in J} f_j^{-1}(U_j)$
pour une partie $J \subset I$ et des ouverts $U_j \subset |X_j|$, voir
Topologie, lemme \ref{topology-lemma-describe-limits}.
Il existe donc $j \in J$ tel que $f_j(x') \in U_j$
et $f_j(x) \not \in U_j$. Autrement dit, $f_j(x') \not \in Z_j$.
Ainsi $Z = \lim Z_i$ en tant qu'ensembles.

\medskip\noindent
Munissons ensuite $Z$ et $Z_i$ de leur structure induite réduite, voir
Propriétés des espaces, définition
\ref{spaces-properties-definition-reduced-induced-space}.
Les morphismes de transition $X_{i'} \to X_i$ induisent des morphismes affines
$Z_{i'} \to Z_i$, et les projections $X \to X_i$
induisent des morphismes compatibles $Z \to Z_i$.
On obtient donc des morphismes $Z \to \lim Z_i \to X$ d'espaces algébriques.
D'après le lemme \ref{lemma-directed-inverse-system-closed-immersions},
le morphisme $\lim Z_i \to X$ est une
immersion fermée. D'après le lemme \ref{lemma-directed-inverse-system-reduced},
l'espace algébrique $\lim Z_i$ est réduit.
D'après ce qui précède, $Z \to \lim Z_i$ est bijectif sur les points.
Par unicité de la structure induite réduite de sous-espace fermé,
ce morphisme est un isomorphisme d'espaces algébriques.
\end{proof}

\begin{situation}
\label{situation-descent}
Soit $S$ un schéma. Soit $X = \lim_{i \in I} X_i$ la limite d'un système projectif
filtrant d'espaces algébriques sur $S$ à morphismes de transition affines
(lemme \ref{lemma-directed-inverse-system-has-limit}).
On suppose que $X_i$ est quasi-compact et quasi-séparé pour tout $i \in I$.
On choisit aussi un élément $0 \in I$.
\end{situation}

\begin{lemma}
\label{lemma-descend-section}
Notations et hypothèses comme dans la situation \ref{situation-descent}.
Supposons que $\mathcal{F}_0$ soit un faisceau quasi-cohérent sur $X_0$.
Posons $\mathcal{F}_i = f_{0i}^*\mathcal{F}_0$ pour $i \geq 0$ et posons
$\mathcal{F} = f_0^*\mathcal{F}_0$. Alors
$$
\Gamma(X, \mathcal{F}) = \colim_{i \geq 0} \Gamma(X_i, \mathcal{F}_i)
$$
\end{lemma}

\begin{proof}
Choisissons un morphisme étale surjectif $U_0 \to X_0$, où $U_0$ est un schéma
affine (Propriétés des espaces, lemme
\ref{spaces-properties-lemma-quasi-compact-affine-cover}).
Posons $U_i = X_i \times_{X_0} U_0$.
Posons $R_0 = U_0 \times_{X_0} U_0$ et $R_i = R_0 \times_{X_0} X_i$.
Dans la démonstration du lemme \ref{lemma-directed-inverse-system-has-limit}, nous avons
vu qu'il existe une présentation $X = U/R$ avec
$U = \lim U_i$ et $R = \lim R_i$.
Notons que $U_i$ et $U$ sont affines, et que $R_i$ et $R$ sont
quasi-compacts et séparés (puisque $X_i$ est quasi-séparé). Par conséquent,
Limites, lemme \ref{limits-lemma-descend-section}
donne
$$
\mathcal{F}(U) = \colim \mathcal{F}_i(U_i)
\quad\text{et}\quad
\mathcal{F}(R) = \colim \mathcal{F}_i(R_i).
$$
Le lemme en résulte car
$\Gamma(X, \mathcal{F}) = \Ker(\mathcal{F}(U) \to \mathcal{F}(R))$
et, de même,
$\Gamma(X_i, \mathcal{F}_i) =
\Ker(\mathcal{F}_i(U_i) \to \mathcal{F}_i(R_i))$
\end{proof}

\begin{lemma}
\label{lemma-descend-opens}
Notations et hypothèses comme dans la situation \ref{situation-descent}.
Pour tout sous-espace ouvert quasi-compact $U \subset X$, il existe un $i$
et un ouvert quasi-compact $U_i \subset X_i$ dont l'image inverse dans $X$ est $U$.
\end{lemma}

\begin{proof}
Cela résulte formellement de la construction des limites dans le
lemme \ref{lemma-directed-inverse-system-has-limit}
et du résultat correspondant pour les schémas :
Limites, lemme \ref{limits-lemma-descend-opens}.
\end{proof}

\noindent
Le lemme suivant sera supplanté par le lemme plus fort
\ref{lemma-descend-isomorphism}.

\begin{lemma}
\label{lemma-descend-equality}
Notations et hypothèses comme dans la situation \ref{situation-descent}.
Soit $f_0 : Y_0 \to Z_0$ un morphisme d'espaces algébriques sur $X_0$.
Supposons (a) que $Y_0 \to X_0$ et $Z_0 \to X_0$ soient représentables, (b) que
$Y_0$ et $Z_0$ soient quasi-compacts et quasi-séparés, (c) que
$f_0$ soit localement de présentation finie, et
(d) que $Y_0 \times_{X_0} X \to Z_0 \times_{X_0} X$ soit un isomorphisme.
Il existe alors un $i \geq 0$ tel que
$Y_0 \times_{X_0} X_i \to Z_0 \times_{X_0} X_i$ soit un isomorphisme.
\end{lemma}

\begin{proof}
Choisissons un schéma affine $U_0$ et un morphisme étale surjectif $U_0 \to X_0$.
Posons $U_i = U_0 \times_{X_0} X_i$ et $U = U_0 \times_{X_0} X$.
Appliquons Limites, lemme \ref{limits-lemma-descend-isomorphism} :
on obtient que $Y_0 \times_{X_0} U_i \to Z_0 \times_{X_0} U_i$
est un isomorphisme de schémas pour un certain $i \geq 0$ (détails omis).
Comme $U_i \to X_i$ est étale surjectif, il s'ensuit que
$Y_0 \times_{X_0} X_i \to Z_0 \times_{X_0} X_i$ est un isomorphisme
(détails omis).
\end{proof}

\begin{lemma}
\label{lemma-descend-separated}
Notations et hypothèses comme dans la situation \ref{situation-descent}.
Si $X$ est séparé, alors $X_i$ est séparé pour un certain $i \in I$.
\end{lemma}

\begin{proof}
Choisissons un schéma affine $U_0$ et un morphisme étale surjectif $U_0 \to X_0$.
Pour $i \geq 0$, posons $U_i = U_0 \times_{X_0} X_i$, et posons
$U = U_0 \times_{X_0} X$. Notons que $U_i$ et $U$ sont des schémas affines
munis de morphismes étales surjectifs $U_i \to X_i$
et $U \to X$. Posons $R_i = U_i \times_{X_i} U_i$ et $R = U \times_X U$,
de projections $s_i, t_i : R_i \to U_i$ et $s, t : R \to U$.
Notons que $R_i$ et $R$ sont des schémas quasi-compacts et séparés (puisque les
espaces algébriques $X_i$ et $X$ sont quasi-séparés). Les morphismes
$s_i : R_i \to U_i$ et $s : R \to U$ sont de type fini.
Par définition, $X_i$ est séparé si et seulement si
$(t_i, s_i) : R_i \to U_i \times U_i$
est une immersion fermée; comme $X$ est séparé par hypothèse,
le morphisme $(t, s) : R \to U \times U$ est une immersion fermée. Comme
$R \to U$ est de type fini, il existe un
$i$ tel que le morphisme $R \to U_i \times U$ soit une immersion fermée
(Limites, lemme \ref{limits-lemma-finite-type-eventually-closed}).
Fixons un tel $i \in I$. Appliquons Limites, lemme
\ref{limits-lemma-descend-closed-immersion-finite-presentation}
au système de morphismes $R_{i'} \to U_i \times U_{i'}$ pour $i' \geq i$
(ce qui est permis puisque
$R_{i'} = R_i \times_{U_i \times U_i} U_i \times U_{i'}$).
On voit ainsi que $R_{i'} \to U_i \times U_{i'}$ est une immersion fermée
pour $i'$ assez grand. Il s'ensuit immédiatement
que $R_{i'} \to U_{i'} \times U_{i'}$ est une immersion fermée,
ce qui achève la démonstration du lemme.
\end{proof}

\begin{lemma}
\label{lemma-limit-is-affine}
Notations et hypothèses comme dans la situation \ref{situation-descent}.
Si $X$ est affine, il existe un $i$ tel que $X_i$ soit affine.
\end{lemma}

\begin{proof}
Choisissons $0 \in I$, puis un schéma affine $U_0$ et un morphisme
étale surjectif $U_0 \to X_0$. Posons $U = U_0 \times_{X_0} X$
et $U_i = U_0 \times_{X_0} X_i$ pour $i \geq 0$. Comme les morphismes de transition
sont affines, les espaces algébriques $U_i$ et $U$ sont affines.
Ainsi, $U \to X$ est un morphisme étale de schémas affines. On peut donc
écrire $X = \Spec(A)$, $U = \Spec(B)$ et
$$
B = A[x_1, \ldots, x_n]/(g_1, \ldots, g_n)
$$
où $\Delta = \det(\partial g_\lambda/\partial x_\mu)$ est inversible
dans $B$; voir Algèbre, lemme \ref{algebra-lemma-etale-standard-smooth}.
Posons $A_i = \mathcal{O}_{X_i}(X_i)$. On a $A = \colim A_i$ d'après le
lemme \ref{lemma-descend-section}. Quitte à augmenter $0$, on peut supposer
donnés $g_{1, i}, \ldots, g_{n, i} \in A_i[x_1, \ldots, x_n]$ dont les images sont
$g_1, \ldots, g_n$. Posons
$$
B_i = A_i[x_1, \ldots, x_n]/(g_{1, i}, \ldots, g_{n, i})
$$
pour tout $i \geq 0$. Quitte à augmenter $0$ si nécessaire, on peut supposer que
$\Delta_i = \det(\partial g_{\lambda, i}/\partial x_\mu)$ est inversible
dans $B_i$ pour tout $i \geq 0$. Ainsi, $A_i \to B_i$ est un homomorphisme étale.
Quitte à augmenter $0$, on peut aussi supposer que
$\Spec(B_i) \to \Spec(A_i)$ est surjectif; voir
Limites, lemme \ref{limits-lemma-descend-surjective}. Quitte à augmenter
encore $0$, choisissons des éléments
$h_{1, i}, \ldots, h_{n, i} \in \mathcal{O}_{U_i}(U_i)$ dont les images sont les
classes de $x_1, \ldots, x_n$ dans $B = \mathcal{O}_U(U)$ et tels que
$g_{\lambda, i}(h_{\nu, i}) = 0$ dans $\mathcal{O}_{U_i}(U_i)$. Ainsi,
on obtient un diagramme commutatif
\begin{equation}
\label{equation-to-show-cartesian}
\vcenter{
\xymatrix{
X_i \ar[d] & U_i \ar[l] \ar[d] \\
\Spec(A_i) & \Spec(B_i) \ar[l]
}
}
\end{equation}
Par construction, $B_i = B_0 \otimes_{A_0} A_i$ et
$B = B_0 \otimes_{A_0} A$. Considérons le morphisme
$$
f_0 : U_0 \longrightarrow X_0 \times_{\Spec(A_0)} \Spec(B_0)
$$
C'est un morphisme d'espaces algébriques quasi-compacts et quasi-séparés,
représentable, séparé et étale sur $X_0$. Par nos choix, le changement de base de $f_0$
à $X$ est un isomorphisme. Le lemme
\ref{lemma-descend-equality}
assure donc qu'il existe un $i$ tel que le changement de base de $f_0$
à $X_i$ soit un isomorphisme; autrement dit, le diagramme
(\ref{equation-to-show-cartesian}) est cartésien. Par conséquent,
Descente, lemme \ref{descent-lemma-descent-data-sheaves},
appliqué au recouvrement fppf $\{\Spec(B_i) \to \Spec(A_i)\}$
et combiné avec Descente, lemme \ref{descent-lemma-affine},
montre que $X_i \to \Spec(A_i)$ est représentable par un schéma
affine sur $\Spec(A_i)$, comme voulu. (Il s'ensuit bien sûr également
que $X_i = \Spec(A_i)$, mais nous n'en avons pas besoin.)
\end{proof}

\begin{lemma}
\label{lemma-limit-is-scheme}
Notations et hypothèses comme dans la situation \ref{situation-descent}.
Si $X$ est un schéma, il existe un $i$ tel que $X_i$ soit un schéma.
\end{lemma}

\begin{proof}
Choisissons un recouvrement ouvert affine fini $X = \bigcup W_j$.
D'après le lemme \ref{lemma-descend-opens},
il existe un $i \in I$ et des sous-espaces ouverts $W_{j, i} \subset X_i$
dont le changement de base à $X$ est $W_j \to X$. D'après le
lemme \ref{lemma-limit-is-affine}, on peut supposer que
chaque $W_{j, i}$ est un schéma affine. Il en résulte que $X_i$
est un schéma (voir, par exemple,
Propriétés des espaces, section \ref{spaces-properties-section-schematic}).
\end{proof}

\begin{lemma}
\label{lemma-finite-type-eventually-closed}
Soit $S$ un schéma. Soit $B$ un espace algébrique sur $S$.
Soit $X = \lim X_i$ la limite d'un système projectif filtrant
d'espaces algébriques sur $B$ à morphismes de transition affines.
Soit $Y \to X$ un morphisme d'espaces algébriques sur $B$.
\begin{enumerate}
\item Si $Y \to X$ est une immersion fermée, si $X_i$ est quasi-compact et si
$Y \to B$ est localement de type fini, alors $Y \to X_i$ est une immersion
fermée pour $i$ assez grand.
\item Si $Y \to X$ est une immersion, si $X_i$ est quasi-séparé, si $Y \to B$
est localement de type fini et si $Y$ est quasi-compact, alors $Y \to X_i$ est
une immersion pour $i$ assez grand.
\item Si $Y \to X$ est un isomorphisme, si $X_i$ est quasi-compact,
si $X_i \to B$ est localement de type fini, si les morphismes de transition
$X_{i'} \to X_i$ sont des immersions fermées et si $Y \to B$ est localement
de présentation finie, alors $Y \to X_i$ est un isomorphisme pour $i$
assez grand.
\item Si $Y \to X$ est un monomorphisme, si $X_i$ est quasi-séparé,
si $Y \to B$ est localement de type fini et si $Y$ est quasi-compact, alors
$Y \to X_i$ est un monomorphisme pour $i$ assez grand.
\end{enumerate}
\end{lemma}

\begin{proof}
Démonstration de (1). Choisissons $0 \in I$. Comme $X_0$ est quasi-compact,
on peut choisir un schéma affine $W$ et un morphisme étale $W \to B$ tels que
l'image de $|X_0| \to |B|$ soit contenue dans celle de $|W| \to |B|$.
Choisissons un schéma affine $U_0$ et un morphisme étale
$U_0 \to X_0 \times_B W$ tels que $U_0 \to X_0$ soit surjectif.
(Cela est possible par le choix de $W$ et la quasi-compacité de $X_0$;
les détails sont omis.)
Notons $V \to Y$, resp.\ $U \to X$, resp.\ $U_i \to X_i$, les changements
de base de $U_0 \to X_0$ (pour $i \geq 0$). Il suffit de prouver que
$V \to U_i$ est une immersion fermée pour $i$ assez grand. On est donc ramené
à démontrer le résultat pour $V \to U = \lim U_i$ sur $W$. Le résultat découle
alors du cas des schémas, traité dans Limites, lemme
\ref{limits-lemma-finite-type-eventually-closed}.

\medskip\noindent
Démonstration de (2). Choisissons $0 \in I$. Choisissons un sous-espace ouvert
quasi-compact $X'_0 \subset X_0$ tel que $Y \to X_0$ se factorise par $X'_0$.
Après avoir remplacé $X_i$ par l'image inverse de $X'_0$ pour $i \geq 0$,
on peut supposer que tous les $X_i'$ sont quasi-compacts et quasi-séparés.
Soit $U \subset X$ un ouvert quasi-compact tel que $Y \to X$ se factorise
par une immersion fermée $Y \to U$ (un tel $U$ existe puisque $Y$ est
quasi-compact). D'après le lemme \ref{lemma-descend-opens},
on peut supposer que $U = \lim U_i$, où $U_i \subset X_i$ est un ouvert
quasi-compact. D'après (1), $Y \to U_i$ est une immersion fermée pour un
certain $i$. Ainsi, (2) est établi.

\medskip\noindent
Démonstration de (3). Choisissons $0 \in I$, puis un schéma affine $U_0$
et un morphisme étale surjectif $U_0 \to X_0$.
Posons $U_i = X_i \times_{X_0} U_0$ et
$U = X \times_{X_0} U_0 = Y \times_{X_0} U_0$. Alors $U = \lim U_i$ est une
limite de schémas affines, les morphismes de transition du système sont des
immersions fermées, et $U \to U_0$ est de présentation finie (car
$U \to B$ est localement de présentation finie, $U_0 \to B$ est localement
de type fini, et d'après
Morphismes d'espaces, lemme
\ref{spaces-morphisms-lemma-finite-presentation-permanence}).
On est ainsi ramené au fait algébrique suivant : si $A = \lim A_i$
est une limite inductive filtrante de $R$-algèbres à morphismes de transition
surjectifs et si $A$ est de présentation finie sur $A_0$, alors $A = A_i$
pour un certain $i$. En effet, écrivons $A = A_0/(f_1, \ldots, f_n)$.
Choisissons $i$ tel que $f_1, \ldots, f_n$ aient une image nulle par le morphisme surjectif $A_0 \to A_i$.

\medskip\noindent
Démonstration de (4). Posons $Z_i = Y \times_{X_i} Y$.
Comme les morphismes de transition $X_{i'} \to X_i$ sont affines, donc
séparés, les morphismes de transition $Z_{i'} \to Z_i$ sont des immersions
fermées; voir Morphismes d'espaces, lemme
\ref{spaces-morphisms-lemma-fibre-product-after-map}.
On a $\lim Z_i = Y \times_X Y = Y$, puisque $Y \to X$ est un monomorphisme.
Choisissons $0 \in I$. Comme $Y \to X_0$ est localement de type fini
(Morphismes d'espaces, lemme
\ref{spaces-morphisms-lemma-permanence-finite-type})
le morphisme $Y \to Z_0$ est localement de présentation finie
(Morphismes d'espaces, lemme
\ref{spaces-morphisms-lemma-diagonal-morphism-finite-type}).
Les morphismes $Z_i \to Z_0$ sont localement de type fini
(ce sont des immersions fermées).
Enfin, $Z_i = Y \times_{X_i} Y$ est quasi-compact puisque
$X_i$ est quasi-séparé et $Y$ est quasi-compact.
La partie (3) s'applique donc à $Y = \lim_{i \geq 0} Z_i$ sur $Z_0$,
et l'on conclut que $Y = Z_i$ pour un certain $i$. Cela démontre (4) et le lemme.
\end{proof}

\begin{lemma}
\label{lemma-eventually-separated}
Soit $S$ un schéma. Soit $Y$ un espace algébrique sur $S$.
Soit $X = \lim X_i$ la limite d'un système projectif filtrant d'espaces
algébriques sur $Y$ à morphismes de transition affines. Supposons que
\begin{enumerate}
\item $Y$ soit quasi-séparé,
\item $X_i$ soit quasi-compact et quasi-séparé,
\item le morphisme $X \to Y$ soit séparé.
\end{enumerate}
Alors $X_i \to Y$ est séparé pour tout $i$ assez grand.
\end{lemma}

\begin{proof}
Soit $0 \in I$. Choisissons un schéma affine $W$ et un morphisme étale
$W \to Y$ tels que l'image de $|W| \to |Y|$ contienne celle de
$|X_0| \to |Y|$. Cela est possible puisque $X_0$ est quasi-compact.
Il suffit de vérifier que $W \times_Y X_i \to W$ est séparé
pour un certain $i \geq 0$, car la diagonale de $W \times_Y X_i$
sur $W$ est le changement de base de $X_i \to X_i \times_Y X_i$ par
le morphisme étale surjectif $(X_i \times_Y X_i) \times_Y W \to
X_i \times_Y X_i$. Comme $Y$ est quasi-séparé, les espaces algébriques
$W \times_Y X_i$ sont quasi-compacts (et aussi quasi-séparés).
On peut donc effectuer le changement de base à $W$ et supposer que $Y$ est un schéma affine.
Lorsque $Y$ est un schéma affine, il faut montrer que $X_i$ est un espace
algébrique séparé pour $i$ assez grand, sachant que $X$ est un espace
algébrique séparé. Ce cas résulte donc du
lemme \ref{lemma-descend-separated}.
\end{proof}

\begin{lemma}
\label{lemma-eventually-affine}
Soit $S$ un schéma. Soit $Y$ un espace algébrique sur $S$.
Soit $X = \lim X_i$ la limite d'un système projectif filtrant d'espaces
algébriques sur $Y$ à morphismes de transition affines. Supposons que
\begin{enumerate}
\item $Y$ soit quasi-compact et quasi-séparé,
\item $X_i$ soit quasi-compact et quasi-séparé,
\item $X \to Y$ soit affine.
\end{enumerate}
Alors $X_i \to Y$ est affine pour $i$ assez grand.
\end{lemma}

\begin{proof}
Choisissons un schéma affine $W$ et un morphisme étale surjectif $W \to Y$.
Alors $X \times_Y W$ est affine, et il suffit de vérifier que
$X_i \times_Y W$ est affine pour un certain $i$ (Morphismes d'espaces,
lemme \ref{spaces-morphisms-lemma-affine-local}).
Cela résulte du lemme \ref{lemma-limit-is-affine}.
\end{proof}

\begin{lemma}
\label{lemma-eventually-finite}
Soit $S$ un schéma. Soit $Y$ un espace algébrique sur $S$.
Soit $X = \lim X_i$ la limite d'un système projectif filtrant d'espaces
algébriques sur $Y$ à morphismes de transition affines. Supposons que
\begin{enumerate}
\item $Y$ soit quasi-compact et quasi-séparé,
\item $X_i$ soit quasi-compact et quasi-séparé,
\item les morphismes de transition $X_{i'} \to X_i$ soient finis,
\item $X_i \to Y$ soit localement de type fini,
\item $X \to Y$ soit intégral.
\end{enumerate}
Alors $X_i \to Y$ est fini pour $i$ assez grand.
\end{lemma}

\begin{proof}
Choisissons un schéma affine $W$ et un morphisme étale surjectif $W \to Y$.
Alors $X \times_Y W$ est fini sur $W$, et il suffit de vérifier que
$X_i \times_Y W$ est fini sur $W$ pour un certain $i$ (Morphismes d'espaces,
lemme \ref{spaces-morphisms-lemma-integral-local}). D'après le
lemme \ref{lemma-limit-is-scheme}, on est ramené au cas des schémas.
Dans ce cas, le résultat découle de
Limites, lemme \ref{limits-lemma-eventually-finite}.
\end{proof}

\begin{lemma}
\label{lemma-eventually-closed-immersion}
Soit $S$ un schéma. Soit $Y$ un espace algébrique sur $S$.
Soit $X = \lim X_i$ la limite d'un système projectif filtrant d'espaces
algébriques sur $Y$ à morphismes de transition affines. Supposons que
\begin{enumerate}
\item $Y$ soit quasi-compact et quasi-séparé,
\item $X_i$ soit quasi-compact et quasi-séparé,
\item les morphismes de transition $X_{i'} \to X_i$ soient des immersions fermées,
\item $X_i \to Y$ soit localement de type fini,
\item $X \to Y$ soit une immersion fermée.
\end{enumerate}
Alors $X_i \to Y$ est une immersion fermée pour $i$ assez grand.
\end{lemma}

\begin{proof}
Choisissons un schéma affine $W$ et un morphisme étale surjectif $W \to Y$.
Alors $X \times_Y W$ est un sous-espace fermé de $W$, et il suffit de vérifier
que $X_i \times_Y W$ est un sous-espace fermé de $W$ pour un certain $i$
(Morphismes d'espaces,
lemme \ref{spaces-morphisms-lemma-closed-immersion-local}). D'après le
lemme \ref{lemma-limit-is-scheme}, on est ramené au cas des schémas.
Dans ce cas, le résultat découle de
Limites, lemme \ref{limits-lemma-eventually-closed-immersion}.
\end{proof}

















\section{Descente de propriétés des morphismes}
\label{section-descent-of-properties}

\noindent
Cette section est l'analogue de la section \ref{section-descent}
pour les propriétés des morphismes. Nous nous placerons dans la situation suivante.

\begin{situation}
\label{situation-descent-property}
Soit $S$ un schéma. Soit $B = \lim B_i$ la limite d'un système projectif filtrant
d'espaces algébriques sur $S$ à morphismes de transition affines
(lemme \ref{lemma-directed-inverse-system-has-limit}).
Soit $0 \in I$ et soit $f_0 : X_0 \to Y_0$ un morphisme d'espaces algébriques
sur $B_0$. Supposons $B_0$, $X_0$, $Y_0$ quasi-compacts et quasi-séparés.
Soit $f_i : X_i \to Y_i$ le changement de base de $f_0$ à $B_i$ et
soit $f : X \to Y$ le changement de base de $f_0$ à $B$.
\end{situation}

\begin{lemma}
\label{lemma-descend-etale}
Notations et hypothèses comme dans la
situation \ref{situation-descent-property}. Si
\begin{enumerate}
\item $f$ est étale,
\item $f_0$ est localement de présentation finie,
\end{enumerate}
alors $f_i$ est étale pour un certain $i \geq 0$.
\end{lemma}

\begin{proof}
Choisissons un schéma affine $V_0$ et un morphisme étale surjectif
$V_0 \to Y_0$. Choisissons un schéma affine $U_0$ et un morphisme étale
surjectif $U_0 \to V_0 \times_{Y_0} X_0$. On a le diagramme
$$
\xymatrix{
U_0 \ar[d] \ar[r] & V_0 \ar[d] \\
X_0 \ar[r] & Y_0
}
$$
Les flèches verticales sont étales et surjectives par construction.
Après changement de base de ce diagramme à $B_i$ ou à $B$, on obtient
$$
\vcenter{
\xymatrix{
U_i \ar[d] \ar[r] & V_i \ar[d] \\
X_i \ar[r] & Y_i
}
}
\quad\text{et}\quad
\vcenter{
\xymatrix{
U \ar[d] \ar[r] & V \ar[d] \\
X \ar[r] & Y
}
}
$$
Remarquons que $U_i, V_i, U, V$ sont des schémas affines,
que les morphismes verticaux sont étales et surjectifs, et que la limite des
morphismes $U_i \to V_i$ est $U \to V$. Rappelons que $X_i \to Y_i$ est étale
si et seulement si $U_i \to V_i$ est
étale et, de même, que $X \to Y$ est étale si et seulement si
$U \to V$ est étale
(Morphismes des espaces, lemme \ref{spaces-morphisms-lemma-etale-local}).
Comme $f_0$ est localement de présentation
finie, le morphisme $U_0 \to V_0$ l'est aussi. Le lemme résulte donc
de Limites, lemme \ref{limits-lemma-descend-etale}.
\end{proof}

\begin{lemma}
\label{lemma-descend-smooth}
Notations et hypothèses comme dans la
situation \ref{situation-descent-property}. Si
\begin{enumerate}
\item $f$ est lisse,
\item $f_0$ est localement de présentation finie,
\end{enumerate}
alors $f_i$ est lisse pour un certain $i \geq 0$.
\end{lemma}

\begin{proof}
Choisissons un schéma affine $V_0$ et un morphisme étale surjectif
$V_0 \to Y_0$. Choisissons un schéma affine $U_0$ et un morphisme étale
surjectif $U_0 \to V_0 \times_{Y_0} X_0$. On a le diagramme
$$
\xymatrix{
U_0 \ar[d] \ar[r] & V_0 \ar[d] \\
X_0 \ar[r] & Y_0
}
$$
Les flèches verticales sont étales et surjectives par construction.
Après changement de base de ce diagramme à $B_i$ ou à $B$, on obtient
$$
\vcenter{
\xymatrix{
U_i \ar[d] \ar[r] & V_i \ar[d] \\
X_i \ar[r] & Y_i
}
}
\quad\text{et}\quad
\vcenter{
\xymatrix{
U \ar[d] \ar[r] & V \ar[d] \\
X \ar[r] & Y
}
}
$$
Remarquons que $U_i, V_i, U, V$ sont des schémas affines,
que les morphismes verticaux sont étales et surjectifs, et que la limite des
morphismes $U_i \to V_i$ est $U \to V$. Rappelons que $X_i \to Y_i$ est lisse
si et seulement si $U_i \to V_i$ est lisse et, de même,
$X \to Y$ est lisse si et seulement si $U \to V$ est lisse
(Morphismes des espaces, définition \ref{spaces-morphisms-definition-smooth}).
Comme $f_0$ est localement de présentation
finie, le morphisme $U_0 \to V_0$ l'est aussi. Le lemme résulte donc
de Limites, lemme \ref{limits-lemma-descend-smooth}.
\end{proof}

\begin{lemma}
\label{lemma-descend-surjective}
Notations et hypothèses comme dans la
situation \ref{situation-descent-property}. Si
\begin{enumerate}
\item $f$ est surjectif,
\item $f_0$ est localement de présentation finie,
\end{enumerate}
alors $f_i$ est surjectif pour un certain $i \geq 0$.
\end{lemma}

\begin{proof}
Choisissons un schéma affine $V_0$ et un morphisme étale surjectif
$V_0 \to Y_0$. Choisissons un schéma affine $U_0$ et un morphisme étale
surjectif $U_0 \to V_0 \times_{Y_0} X_0$. On a le diagramme
$$
\xymatrix{
U_0 \ar[d] \ar[r] & V_0 \ar[d] \\
X_0 \ar[r] & Y_0
}
$$
Les flèches verticales sont étales et surjectives par construction.
Après changement de base de ce diagramme à $B_i$ ou à $B$, on obtient
$$
\vcenter{
\xymatrix{
U_i \ar[d] \ar[r] & V_i \ar[d] \\
X_i \ar[r] & Y_i
}
}
\quad\text{et}\quad
\vcenter{
\xymatrix{
U \ar[d] \ar[r] & V \ar[d] \\
X \ar[r] & Y
}
}
$$
Remarquons que $U_i, V_i, U, V$ sont des schémas affines, que les morphismes verticaux sont
étales et surjectifs, que la limite des morphismes $U_i \to V_i$ est $U \to V$,
et que les morphismes $U_i \to X_i \times_{Y_i} V_i$ et
$U \to X \times_Y V$ sont surjectifs (car ce sont les changements de base de
$U_0 \to X_0 \times_{Y_0} V_0$). En particulier, on voit que
$X_i \to Y_i$ est surjectif si et seulement si $U_i \to V_i$ est surjectif
et, de même, que $X \to Y$ est surjectif si et seulement si $U \to V$ est surjectif.
Comme $f_0$ est localement de présentation
finie, le morphisme $U_0 \to V_0$ l'est aussi. Le lemme résulte donc
du cas des schémas (Limites, lemme \ref{limits-lemma-descend-surjective}).
\end{proof}

\begin{lemma}
\label{lemma-descend-universally-injective}
Notations et hypothèses comme dans la situation \ref{situation-descent-property}. Si
\begin{enumerate}
\item $f$ est universellement injectif,
\item $f_0$ est localement de type fini,
\end{enumerate}
alors $f_i$ est universellement injectif pour un certain $i \geq 0$.
\end{lemma}

\begin{proof}
Rappelons qu'un morphisme $X \to Y$ est universellement injectif si et
seulement si la diagonale $X \to X \times_Y X$ est surjective
(Morphismes des espaces, définition
\ref{spaces-morphisms-definition-universally-injective} et
lemme \ref{spaces-morphisms-lemma-universally-injective}).
Remarquons que $X_0 \to X_0 \times_{Y_0} X_0$ est localement de
présentation finie (Morphismes des espaces, lemme
\ref{spaces-morphisms-lemma-diagonal-morphism-finite-type}).
Le lemme résulte donc du lemme \ref{lemma-descend-surjective}
appliqué au morphisme $X_0 \to X_0 \times_{Y_0} X_0$.
\end{proof}

\begin{lemma}
\label{lemma-descend-affine}
Notations et hypothèses comme dans la situation \ref{situation-descent-property}. Si
$f$ est affine, alors $f_i$ est affine pour un certain $i \geq 0$.
\end{lemma}

\begin{proof}
Choisissons un schéma affine $V_0$ et un morphisme étale surjectif $V_0 \to Y_0$.
Posons $V_i = V_0 \times_{Y_0} Y_i$ et $V = V_0 \times_{Y_0} Y$.
Comme $f$ est affine, on voit que $V \times_Y X = \lim V_i \times_{Y_i} X_i$
est affine. D'après le lemme \ref{lemma-limit-is-affine},
$V_i \times_{Y_i} X_i$ est affine pour un certain $i \geq 0$. Pour cet $i$, le morphisme
$f_i$ est affine
(Morphismes des espaces, lemme \ref{spaces-morphisms-lemma-affine-local}).
\end{proof}

\begin{lemma}
\label{lemma-descend-finite}
Notations et hypothèses comme dans la situation \ref{situation-descent-property}. Si
\begin{enumerate}
\item $f$ est fini,
\item $f_0$ est localement de type fini,
\end{enumerate}
alors $f_i$ est fini pour un certain $i \geq 0$.
\end{lemma}

\begin{proof}
Choisissons un schéma affine $V_0$ et un morphisme étale surjectif $V_0 \to Y_0$.
Posons $V_i = V_0 \times_{Y_0} Y_i$ et $V = V_0 \times_{Y_0} Y$.
Comme $f$ est fini, on voit que $V \times_Y X = \lim V_i \times_{Y_i} X_i$
est un schéma fini sur $V$. D'après le lemme \ref{lemma-limit-is-affine},
$V_i \times_{Y_i} X_i$ est affine pour un certain $i \geq 0$. Quitte à augmenter $i$ si
nécessaire, on trouve que $V_i \times_{Y_i} X_i \to V_i$ est fini d'après
Limites, lemme \ref{limits-lemma-descend-finite-finite-presentation}.
Pour cet $i$, le morphisme $f_i$ est fini
(Morphismes des espaces, lemme \ref{spaces-morphisms-lemma-integral-local}).
\end{proof}

\begin{lemma}
\label{lemma-descend-closed-immersion}
Notations et hypothèses comme dans la situation \ref{situation-descent-property}. Si
\begin{enumerate}
\item $f$ est une immersion fermée,
\item $f_0$ est localement de type fini,
\end{enumerate}
alors $f_i$ est une immersion fermée pour un certain $i \geq 0$.
\end{lemma}

\begin{proof}
Choisissons un schéma affine $V_0$ et un morphisme étale surjectif $V_0 \to Y_0$.
Posons $V_i = V_0 \times_{Y_0} Y_i$ et $V = V_0 \times_{Y_0} Y$.
Comme $f$ est une immersion fermée, on voit que
$V \times_Y X = \lim V_i \times_{Y_i} X_i$
est un sous-schéma fermé du schéma affine $V$. D'après le
lemme \ref{lemma-limit-is-affine},
$V_i \times_{Y_i} X_i$ est affine pour un certain $i \geq 0$. Quitte à augmenter $i$ si
nécessaire, on trouve que $V_i \times_{Y_i} X_i \to V_i$ est une immersion fermée d'après
Limites, lemme \ref{limits-lemma-descend-closed-immersion-finite-presentation}.
Pour cet $i$, le morphisme $f_i$ est une immersion fermée
(Morphismes des espaces, lemme \ref{spaces-morphisms-lemma-integral-local}).
\end{proof}

\begin{lemma}
\label{lemma-descend-separated-morphism}
Notations et hypothèses comme dans la situation \ref{situation-descent-property}.
Si $f$ est séparé, alors $f_i$ est séparé pour un certain $i \geq 0$.
\end{lemma}

\begin{proof}
Appliquons le lemme \ref{lemma-descend-closed-immersion}
au morphisme diagonal $\Delta_{X_0/Y_0} : X_0 \to X_0 \times_{Y_0} X_0$.
(Les morphismes diagonaux sont localement de type fini
et le produit fibré $X_0 \times_{Y_0} X_0$ est quasi-compact et
quasi-séparé. Certains détails sont omis.)
\end{proof}

\begin{lemma}
\label{lemma-descend-isomorphism}
Notations et hypothèses comme dans la situation \ref{situation-descent-property}. Si
\begin{enumerate}
\item $f$ est un isomorphisme,
\item $f_0$ est localement de présentation finie,
\end{enumerate}
alors $f_i$ est un isomorphisme pour un certain $i \geq 0$.
\end{lemma}

\begin{proof}
Être un isomorphisme équivaut à être étale, universellement injectif
et surjectif, voir
Morphismes des espaces, lemme
\ref{spaces-morphisms-lemma-etale-universally-injective-open}.
Le lemme résulte donc des
lemmes \ref{lemma-descend-etale},
\ref{lemma-descend-surjective} et
\ref{lemma-descend-universally-injective}.
\end{proof}

\begin{lemma}
\label{lemma-descend-monomorphism}
Notations et hypothèses comme dans la situation \ref{situation-descent-property}. Si
\begin{enumerate}
\item $f$ est un monomorphisme,
\item $f_0$ est localement de type fini,
\end{enumerate}
alors $f_i$ est un monomorphisme pour un certain $i \geq 0$.
\end{lemma}

\begin{proof}
Rappelons qu'un morphisme est un monomorphisme si et seulement si sa diagonale est
un isomorphisme. Le morphisme $X_0 \to X_0 \times_{Y_0} X_0$ est localement de
présentation finie d'après
Morphismes des espaces, lemme
\ref{spaces-morphisms-lemma-diagonal-morphism-finite-type}.
Comme $X_0 \times_{Y_0} X_0$ est quasi-compact et quasi-séparé,
on déduit du
lemme \ref{lemma-descend-isomorphism}
que $\Delta_i : X_i \to X_i \times_{Y_i} X_i$ est un isomorphisme pour
un certain $i \geq 0$. Pour cet $i$, le morphisme $f_i$ est un monomorphisme.
\end{proof}

\begin{lemma}
\label{lemma-descend-flat}
Notations et hypothèses comme dans la situation \ref{situation-descent-property}.
Soit $\mathcal{F}_0$ un $\mathcal{O}_{X_0}$-module quasi-cohérent;
notons $\mathcal{F}_i$ son image inverse sur $X_i$ et $\mathcal{F}$
son image inverse sur $X$. Si
\begin{enumerate}
\item $\mathcal{F}$ est plat sur $Y$,
\item $\mathcal{F}_0$ est de présentation finie, et
\item $f_0$ est localement de présentation finie,
\end{enumerate}
alors $\mathcal{F}_i$ est plat sur $Y_i$ pour un certain $i \geq 0$.
En particulier, si $f_0$ est localement de présentation finie et si
$f$ est plat, alors $f_i$ est plat pour un certain $i \geq 0$.
\end{lemma}

\begin{proof}
Choisissons un schéma affine $V_0$ et un morphisme étale surjectif
$V_0 \to Y_0$. Choisissons un schéma affine $U_0$ et un morphisme étale
surjectif $U_0 \to V_0 \times_{Y_0} X_0$. On a le diagramme
$$
\xymatrix{
U_0 \ar[d] \ar[r] & V_0 \ar[d] \\
X_0 \ar[r] & Y_0
}
$$
Les flèches verticales sont étales et surjectives par construction.
Après changement de base de ce diagramme à $B_i$ ou à $B$, on obtient
$$
\vcenter{
\xymatrix{
U_i \ar[d] \ar[r] & V_i \ar[d] \\
X_i \ar[r] & Y_i
}
}
\quad\text{et}\quad
\vcenter{
\xymatrix{
U \ar[d] \ar[r] & V \ar[d] \\
X \ar[r] & Y
}
}
$$
Remarquons que $U_i, V_i, U, V$ sont des schémas affines, que les morphismes verticaux sont
étales et surjectifs, et que la limite des morphismes $U_i \to V_i$ est
$U \to V$. Rappelons que $\mathcal{F}_i$ est plat sur $Y_i$ si et seulement si
$\mathcal{F}_i|_{U_i}$ est plat sur $V_i$ et, de même, que $\mathcal{F}$ est plat
sur $Y$ si et seulement si $\mathcal{F}|_U$ est plat sur $V$
(Morphismes des espaces, définition \ref{spaces-morphisms-definition-flat}).
Comme $f_0$ est localement de présentation finie, le morphisme
$U_0 \to V_0$ l'est aussi. Le lemme résulte donc
de Limites, lemme \ref{limits-lemma-descend-module-flat-finite-presentation}.
\end{proof}

\begin{lemma}
\label{lemma-eventually-proper}
Hypothèses et notations comme dans la situation \ref{situation-descent-property}.
Si
\begin{enumerate}
\item $f$ est propre, et
\item $f_0$ est localement de type fini,
\end{enumerate}
alors il existe un $i$ tel que $f_i$ soit propre.
\end{lemma}

\begin{proof}
Choisissons un schéma affine $V_0$ et un morphisme étale surjectif $V_0 \to Y_0$.
Posons $V_i = Y_i \times_{Y_0} V_0$ et $V = Y \times_{Y_0} V_0$.
Il suffit de montrer que le changement de base de $f_i$ à $V_i$ est
propre, voir Morphismes des espaces, lemme
\ref{spaces-morphisms-lemma-proper-local}.
On peut donc supposer $Y_0$ affine.

\medskip\noindent
D'après le lemme \ref{lemma-descend-separated-morphism}, on voit que
$f_i$ est séparé pour un certain $i \geq 0$. En remplaçant
$0$ par $i$, on peut supposer que $f_0$ est séparé.
Remarquons que $f_0$ est quasi-compact. Ainsi, $f_0$ est séparé et
de type fini. D'après
Cohomologie des espaces, lemme \ref{spaces-cohomology-lemma-weak-chow},
on peut choisir un diagramme
$$
\xymatrix{
X_0 \ar[rd] & X_0' \ar[d] \ar[l]^\pi \ar[r] & \mathbf{P}^n_{Y_0} \ar[dl] \\
& Y_0 &
}
$$
où $X_0' \to \mathbf{P}^n_{Y_0}$ est une immersion, et
$\pi : X_0' \to X_0$ est propre et surjectif. Introduisons
$X' = X_0' \times_{Y_0} Y$ et $X_i' = X_0' \times_{Y_0} Y_i$.
D'après Morphismes des espaces, lemmes
\ref{spaces-morphisms-lemma-composition-proper} et
\ref{spaces-morphisms-lemma-base-change-proper}
on voit que $X' \to Y$ est propre. Ainsi, $X' \to \mathbf{P}^n_Y$ est
une immersion fermée (Morphismes des espaces, lemme
\ref{spaces-morphisms-lemma-universally-closed-permanence}). D'après
Morphismes des espaces, lemme \ref{spaces-morphisms-lemma-image-proper-is-proper},
il suffit de montrer que $X'_i \to Y_i$ est propre pour un certain $i$.
D'après le lemme \ref{lemma-descend-closed-immersion},
on trouve que $X'_i \to \mathbf{P}^n_{Y_i}$ est
une immersion fermée pour $i$ assez grand. Alors $X'_i \to Y_i$
est propre, ce qui conclut la démonstration.
\end{proof}

\begin{lemma}
\label{lemma-eventually-relative-dimension}
Hypothèses et notations comme dans la situation \ref{situation-descent-property}.
Soit $d \geq 0$. Si
\begin{enumerate}
\item $f$ est de dimension relative $\leq d$
(Morphismes des espaces, définition
\ref{spaces-morphisms-definition-relative-dimension}), et
\item $f_0$ est localement de type fini,
\end{enumerate}
alors il existe un $i$ tel que $f_i$ soit de dimension relative $\leq d$.
\end{lemma}

\begin{proof}
Choisissons un schéma affine $V_0$ et un morphisme étale surjectif
$V_0 \to Y_0$. Choisissons un schéma affine $U_0$ et un morphisme étale
surjectif $U_0 \to V_0 \times_{Y_0} X_0$. On a le diagramme
$$
\xymatrix{
U_0 \ar[d] \ar[r] & V_0 \ar[d] \\
X_0 \ar[r] & Y_0
}
$$
Les flèches verticales sont étales et surjectives par construction.
Après changement de base de ce diagramme à $B_i$ ou à $B$, on obtient
$$
\vcenter{
\xymatrix{
U_i \ar[d] \ar[r] & V_i \ar[d] \\
X_i \ar[r] & Y_i
}
}
\quad\text{et}\quad
\vcenter{
\xymatrix{
U \ar[d] \ar[r] & V \ar[d] \\
X \ar[r] & Y
}
}
$$
Remarquons que $U_i, V_i, U, V$ sont des schémas affines,
que les morphismes verticaux sont étales et surjectifs, et que la limite des
morphismes $U_i \to V_i$ est $U \to V$.
Dans cette situation, $X_i \to Y_i$ est de dimension relative $\leq d$
si et seulement si $U_i \to V_i$ est de dimension relative $\leq d$
(au sens de Morphismes, définition
\ref{morphisms-definition-relative-dimension-d}).
Pour voir l'équivalence, utilisons le fait que la définition pour les morphismes
d'espaces algébriques fait intervenir Morphismes des espaces, définition
\ref{spaces-morphisms-definition-dimension-fibre}
qui utilise la localisation étale. Il en va de même pour $X \to Y$ et $U \to V$.
Comme $f_0$ est localement de type fini, le morphisme $U_0 \to V_0$ l'est aussi.
Le lemme résulte donc du résultat plus général
Limites, lemme \ref{limits-lemma-limit-dimension}.
\end{proof}














\section{Descente des objets relatifs}
\label{section-descending-relative}

\noindent
Le lemme suivant est représentatif du type de résultats de cette section.

\begin{lemma}
\label{lemma-descend-finite-presentation}
Soit $S$ un schéma. Soit $I$ un ensemble préordonné filtrant.
Soit $(X_i, f_{ii'})$ un système projectif indexé par $I$ d'espaces algébriques
sur $S$. Supposons que
\begin{enumerate}
\item les morphismes $f_{ii'} : X_i \to X_{i'}$ soient affines,
\item les espaces $X_i$ soient quasi-compacts et quasi-séparés.
\end{enumerate}
Soit $X = \lim_i X_i$. Alors la catégorie des espaces algébriques
de présentation finie sur $X$ est la limite inductive sur $I$ des
catégories des espaces algébriques de présentation finie sur $X_i$.
\end{lemma}

\begin{proof}
Fixons $0 \in I$. Choisissons un morphisme étale surjectif $U_0 \to X_0$, où
$U_0$ est un schéma affine (Propriétés des espaces, lemme
\ref{spaces-properties-lemma-quasi-compact-affine-cover}).
Posons $U_i = X_i \times_{X_0} U_0$. Posons $R_0 = U_0 \times_{X_0} U_0$ et
$R_i = R_0 \times_{X_0} X_i$. Notons $s_i, t_i : R_i \to U_i$ et
$s, t : R \to U$ les deux projections. Dans la démonstration du
lemme \ref{lemma-directed-inverse-system-has-limit}, on a
vu qu'il existe une présentation $X = U/R$ avec
$U = \lim U_i$ et $R = \lim R_i$. Remarquons que $U_i$ et $U$ sont affines et
que $R_i$ et $R$ sont quasi-compacts et séparés (puisque $X_i$ est
quasi-séparé). Soit $Y$ un espace algébrique sur $S$ et soit
$Y \to X$ un morphisme de présentation finie. Posons $V = U \times_X Y$.
C'est un espace algébrique de présentation finie sur $U$.
Choisissons un schéma affine $W$ et un morphisme étale surjectif $W \to V$.
Alors $W \to Y$ est également étale et surjectif. Posons $R' = W \times_Y W$,
de sorte que $Y = W/R'$ (voir Espaces, section \ref{spaces-section-presentations}).
Remarquons que $W$ est un schéma de présentation finie sur $U$ et que $R'$
est un schéma de présentation finie sur $R$ (détails omis).
D'après Limites, lemme \ref{limits-lemma-descend-finite-presentation},
on peut trouver un indice $i$ et un morphisme de schémas $W_i \to U_i$ de
présentation finie dont le changement de base à $U$ donne $W \to U$. De même,
on peut trouver, quitte à augmenter $i$, un schéma $R'_i$ de présentation
finie sur $R_i$ dont le changement de base à $R$ est $R'$.
Les morphismes de projection $s', t' : R' \to W$ sont des morphismes au-dessus des
morphismes de projection $s, t : R \to U$. On peut donc considérer $s'$,
resp.\ $t'$, comme un morphisme entre schémas de présentation finie sur
$U$ (le morphisme structural $R' \to U$ étant donné par $R' \to R$ suivi
de $s$, resp.\ $t$). On peut donc appliquer
de nouveau le lemme \ref{limits-lemma-descend-finite-presentation} de Limites,
pour voir que, quitte à augmenter $i$, il existe des
morphismes $s'_i, t'_i : R'_i \to W_i$ qui, après changement de base à $U$,
sont $S', t'$. D'après Limites, lemmes \ref{limits-lemma-descend-etale} et
\ref{limits-lemma-descend-monomorphism},
on peut supposer que $s'_i, t'_i$ sont étales et que
$j'_i : R'_i \to W_i \times_{X_i} W_i$ est un monomorphisme (ici, on
considère $j'_i$ comme un morphisme de schémas de présentation finie sur $U_i$ via
l'une des projections -- peu importe laquelle). En posant
$Y_i = W_i/R'_i$ (voir Espaces, théorème \ref{spaces-theorem-presentation}),
on obtient un espace algébrique de présentation finie
sur $X_i$ dont le changement de base à $X$ est isomorphe à $Y$.

\medskip\noindent
Cela montre que tout espace algébrique de présentation finie sur $X$ provient
d'un espace algébrique de présentation finie sur un certain $X_i$, c'est-à-dire
que le foncteur du lemme est essentiellement surjectif. Pour
montrer qu'il est pleinement fidèle, considérons un indice $0 \in I$ et deux
espaces algébriques $Y_0, Z_0$ de présentation finie sur $X_0$.
Posons $Y_i = X_i \times_{X_0} Y_0$, $Y = X \times_{X_0} Y_0$,
$Z_i = X_i \times_{X_0} Z_0$ et $Z = X \times_{X_0} Z_0$. Soit
$\alpha : Y \to Z$ un morphisme d'espaces algébriques sur $X$.
Choisissons un morphisme étale surjectif $V_0 \to Y_0$, où $V_0$ est
un schéma affine. Posons $V_i = V_0 \times_{Y_0} Y_i$ et
$V = V_0 \times_{Y_0} Y$, qui sont des schémas affines munis de
morphismes étales surjectifs vers $Y_i$ et $Y$. Le composé
$V \to Y \to Z \to Z_0$ provient d'un morphisme (essentiellement unique)
$V_i \to Z_0$ pour un certain $i \geq 0$ d'après la
proposition \ref{proposition-characterize-locally-finite-presentation}
(appliquée à $Z_0 \to X_0$, qui est de présentation finie par hypothèse).
Quitte à augmenter $i$, les deux composés
$$
V_i \times_{Y_i} V_i \to V_i \to Z_0
$$
sont égaux, puisque c'est vrai à la limite. On obtient donc un morphisme
(essentiellement unique) $Y_i \to Z_0$. Comme c'est un morphisme sur $X_0$,
il induit un morphisme vers $Z_i = Z_0 \times_{X_0} X_i$, comme souhaité.
\end{proof}

\begin{lemma}
\label{lemma-descend-modules-finite-presentation}
Avec les notations et hypothèses du
lemme \ref{lemma-descend-finite-presentation}.
La catégorie des $\mathcal{O}_X$-modules de présentation finie est la
limite inductive sur $I$ des catégories des $\mathcal{O}_{X_i}$-modules de
présentation finie.
\end{lemma}

\begin{proof}
Choisissons $0 \in I$. Choisissons un schéma affine $U_0$ et un morphisme
étale surjectif $U_0 \to X_0$. Posons $U_i = X_i \times_{X_0} U_0$.
Posons $R_0 = U_0 \times_{X_0} U_0$ et $R_i = R_0 \times_{X_0} X_i$.
Notons $s_i, t_i : R_i \to U_i$ et $s, t : R \to U$ les deux
projections. Dans la démonstration du
lemme \ref{lemma-directed-inverse-system-has-limit}, on a
vu qu'il existe une présentation $X = U/R$ avec
$U = \lim U_i$ et $R = \lim R_i$. Remarquons que $U_i$ et $U$ sont affines et
que $R_i$ et $R$ sont quasi-compacts et séparés (puisque $X_i$ est
quasi-séparé). En outre, on a également
$R \times_{s, U, t} R = \colim R_i \times_{s_i, U_i, t_i} R_i$.
On sait donc que
$\QCoh(\mathcal{O}_U) = \colim \QCoh(\mathcal{O}_{U_i})$,
$\QCoh(\mathcal{O}_R) = \colim \QCoh(\mathcal{O}_{R_i})$,
et
$\QCoh(\mathcal{O}_{R \times_{s, U, t} R}) =
\colim \QCoh(\mathcal{O}_{R_i \times_{s_i, U_i, t_i} R_i})$
d'après Limites, lemme \ref{limits-lemma-descend-modules-finite-presentation}.
On a $\QCoh(\mathcal{O}_X) = \QCoh(U, R, s, t, c)$ et
$\QCoh(\mathcal{O}_{X_i}) = \QCoh(U_i, R_i, s_i, t_i, c_i)$,
voir Propriétés des espaces, proposition
\ref{spaces-properties-proposition-quasi-coherent}.
Le résultat en découle donc formellement.
\end{proof}

\begin{lemma}
\label{lemma-descend-invertible-modules}
Avec les notations et hypothèses du
lemme \ref{lemma-descend-finite-presentation}. Alors
\begin{enumerate}
\item tout $\mathcal{O}_X$-module localement libre de type fini est l'image inverse
d'un $\mathcal{O}_{X_i}$-module localement libre de type fini pour un certain $i$,
\item tout $\mathcal{O}_X$-module inversible est l'image inverse d'un
$\mathcal{O}_{X_i}$-module inversible pour un certain $i$.
\end{enumerate}
\end{lemma}

\begin{proof}
Démonstration de (2).
Soit $\mathcal{L}$ un $\mathcal{O}_X$-module inversible. Comme les
modules inversibles sont de présentation finie, on peut trouver un $i$
et des modules $\mathcal{L}_i$ et $\mathcal{N}_i$ de présentation finie
sur $X_i$ tels que $f_i^*\mathcal{L}_i \cong \mathcal{L}$ et
$f_i^*\mathcal{N}_i \cong \mathcal{L}^{\otimes -1}$, voir le
lemme \ref{lemma-descend-modules-finite-presentation}.
Comme l'image inverse commute au produit tensoriel, on voit que
$f_i^*(\mathcal{L}_i \otimes_{\mathcal{O}_{X_i}} \mathcal{N}_i)$
est isomorphe à $\mathcal{O}_X$. Comme le produit tensoriel de
modules de présentation finie est de présentation finie, le même
lemme montre que
$f_{i'i}^*\mathcal{L}_i
\otimes_{\mathcal{O}_{X_{i'}}} f_{i'i}^*\mathcal{N}_i$
est isomorphe à $\mathcal{O}_{X_{i'}}$ pour un certain $i' \geq i$.
Il s'ensuit que $f_{i'i}^*\mathcal{L}_i$ est inversible
(Modules sur les sites, lemme \ref{sites-modules-lemma-invertible})
et la démonstration est terminée.

\medskip\noindent
Démonstration de (1). Omise. Indication : raisonner comme dans la démonstration de (2)
en utilisant qu'un module (sur un site localement annelé) est localement libre
de type fini si et seulement s'il admet un dual, voir
Modules sur les sites, section \ref{sites-modules-section-duals}.
On peut aussi raisonner comme dans la démonstration pour les schémas, voir
Limites, lemme \ref{limits-lemma-descend-invertible-modules}.
\end{proof}














\section{Approximation noethérienne absolue}
\label{section-approximation}

\noindent
Le résultat suivant est \cite[théorème 1.2.2]{CLO}.
Un ingrédient essentiel de la démonstration est
Espaces décents, lemme
\ref{decent-spaces-lemma-filter-quasi-compact-quasi-separated}.

\begin{proposition}
\label{proposition-approximate}
\begin{reference}
Notre démonstration suit de près celle donnée dans \cite[théorème 1.2.2]{CLO}.
\end{reference}
Soit $X$ un espace algébrique quasi-compact et quasi-séparé sur
$\Spec(\mathbf{Z})$. Il existe un ensemble préordonné filtrant $I$
et un système projectif d'espaces algébriques $(X_i, f_{ii'})$ indexé par $I$
tels que
\begin{enumerate}
\item les morphismes de transition $f_{ii'}$ soient affines,
\item chaque $X_i$ soit quasi-séparé et de type fini sur
$\mathbf{Z}$, et
\item $X = \lim X_i$.
\end{enumerate}
\end{proposition}

\begin{proof}
Appliquons Espaces décents, lemme
\ref{decent-spaces-lemma-filter-quasi-compact-quasi-separated}
pour obtenir des sous-espaces ouverts $U_p \subset X$, des schémas $V_p$ et des morphismes
$f_p : V_p \to U_p$ possédant les propriétés énoncées. Remarquons que
$f_n : V_n \to U_n$ est un morphisme étale d'espaces algébriques
dont la restriction à l'image inverse de $T_n = (V_n)_{red}$ est un
isomorphisme. Ainsi $f_n$ est un isomorphisme, par exemple d'après
Morphismes des espaces, lemme
\ref{spaces-morphisms-lemma-etale-universally-injective-open}.
En particulier, $U_n$ est un schéma quasi-compact et séparé.
On peut donc écrire $U_n = \lim U_{n, i}$ comme limite projective filtrante
de schémas de type fini sur $\mathbf{Z}$ dont les morphismes de transition sont
affines, voir Limites, proposition \ref{limits-proposition-approximate}.
En procédant par récurrence descendante sur $p$, on voit donc qu'on a ramené
la question au problème posé dans le paragraphe suivant.

\medskip\noindent
On dispose ici de $U \subset X$, $U = \lim U_i$, $Z \subset X$ et
$f : V \to X$ vérifiant les propriétés suivantes
\begin{enumerate}
\item $X$ est un espace algébrique quasi-compact et quasi-séparé,
\item $V$ est un schéma quasi-compact et séparé,
\item $U \subset X$ est un sous-espace ouvert quasi-compact,
\item $(U_i, g_{ii'})$ est un système projectif filtrant d'espaces algébriques
quasi-séparés
de type fini sur $\mathbf{Z}$, à morphismes de transition affines,
dont la limite est $U$,
\item $Z \subset X$ est un sous-espace fermé tel que $|X| = |U| \amalg |Z|$,
\item $f : V \to X$ est un morphisme étale surjectif tel que
$f^{-1}(Z) \to Z$ soit un isomorphisme.
\end{enumerate}
Problème : montrer que la conclusion de la proposition vaut pour $X$.

\medskip\noindent
Remarquons que $W = f^{-1}(U) \subset V$ est un sous-schéma ouvert quasi-compact
étale sur $U$. On peut donc appliquer les
lemmes \ref{lemma-descend-finite-presentation} et \ref{lemma-descend-etale}
pour trouver un indice $0 \in I$ et un morphisme étale $W_0 \to U_0$
de présentation finie dont le changement de base à $U$ donne $W$. En posant
$W_i = W_0 \times_{U_0} U_i$, on voit que $W = \lim_{i \geq 0} W_i$. Quitte à
augmenter $0$, on peut supposer que les $W_i$ sont des schémas, voir le
lemme \ref{lemma-limit-is-scheme}.
En outre, $W_i$ est de type fini sur $\mathbf{Z}$.

\medskip\noindent
Appliquons Limites, lemme \ref{limits-lemma-approximate}, à
$W = \lim_{i \geq 0} W_i$ et à l'inclusion $W \subset V$. Remplaçons $I$
par l'ensemble préordonné filtrant $J$ fourni par ce lemme. Cela permet
d'écrire $V$ comme limite projective filtrante $V = \lim V_i$ de schémas de type fini sur
$\mathbf{Z}$ à morphismes de transition affines, tels que chaque $V_i$ contienne
$W_i$ comme sous-schéma ouvert (de façon compatible avec les morphismes de transition).
Pour chaque $i$, on peut former la somme amalgamée
$$
\xymatrix{
W_i \ar[r] \ar[d]_\Delta & V_i \ar[d] \\
W_i \times_{U_i} W_i \ar[r] & R_i
}
$$
dans la catégorie des schémas. En effet, la flèche verticale de gauche et la flèche horizontale supérieure
sont des immersions ouvertes de schémas. Autrement dit, on peut construire
$R_i$ par recollement de $V_i$ et de $W_i \times_{U_i} W_i$ le long de l'ouvert commun
$W_i$ (voir Schémas, section \ref{schemes-section-glueing-schemes}). Remarquons que
les projections étales $W_i \times_{U_i} W_i \to W_i$ se prolongent
en morphismes étales $s_i, t_i : R_i \to V_i$. Il est clair que le
morphisme $j_i = (t_i, s_i) : R_i \to V_i \times V_i$ est une relation d'équivalence étale
sur $V_i$. Remarquons que $W_i \times_{U_i} W_i$ est
quasi-compact (puisque $U_i$ est quasi-séparé et $W_i$ quasi-compact)
et que $V_i$ est quasi-compact; ainsi $R_i$ est quasi-compact. Pour
$i \geq i'$, le diagramme
\begin{equation}
\label{equation-cartesian}
\vcenter{
\xymatrix{
R_i \ar[r] \ar[d]_{s_i} & R_{i'} \ar[d]^{s_{i'}} \\
V_i \ar[r] & V_{i'}
}
}
\end{equation}
est cartésien car
$$
(W_{i'} \times_{U_{i'}} W_{i'}) \times_{U_{i'}} U_i =
W_{i'} \times_{U_{i'}} U_i \times_{U_i} U_i \times_{U_{i'}} W_{i'} =
W_i \times_{U_i} W_i.
$$
Considérons l'espace algébrique $X_i = V_i/R_i$ (voir
Espaces, théorème \ref{spaces-theorem-presentation}).
Comme $V_i$ est de type fini sur $\mathbf{Z}$ et que $R_i$ est quasi-compact,
on voit que $X_i$ est quasi-séparé et de type fini sur $\mathbf{Z}$
(voir
Propriétés des espaces, lemme \ref{spaces-properties-lemma-quasi-separated}
et
Morphismes des espaces, lemmes
\ref{spaces-morphisms-lemma-surjection-from-quasi-compact} et
\ref{spaces-morphisms-lemma-finite-type-local}).
Comme la construction de $R_i$ ci-dessus est compatible
avec les morphismes de transition, on obtient des morphismes d'espaces algébriques
$X_i \to X_{i'}$ pour $i \geq i'$. Les diagrammes commutatifs
$$
\xymatrix{
V_i \ar[r] \ar[d] & V_{i'} \ar[d] \\
X_i \ar[r] & X_{i'}
}
$$
sont cartésiens puisque (\ref{equation-cartesian}) l'est, voir
Groupoïdes, lemme \ref{groupoids-lemma-criterion-fibre-product}.
Comme $V_i \to V_{i'}$ est affine, cela implique que $X_i \to X_{i'}$
est affine, voir
Morphismes des espaces, lemme \ref{spaces-morphisms-lemma-affine-local}.
On peut donc former la limite $X' = \lim X_i$ d'après le
lemme \ref{lemma-directed-inverse-system-has-limit}.
Nous affirmons que $X \cong X'$, ce qui achève la démonstration de la proposition.

\medskip\noindent
Démonstration de l'affirmation. Posons $R = \lim R_i$.
Par construction, l'espace algébrique $X'$ est
muni d'un morphisme étale surjectif $V \to X'$ tel que
$$
V \times_{X'} V \cong R
$$
(utiliser le lemme \ref{lemma-directed-inverse-system-has-limit}).
Par construction, $\lim W_i \times_{U_i} W_i = W \times_U W$ et $V = \lim V_i$,
de sorte que $R$ est la réunion de $W \times_U W$ et de $V$, recollés le long de $W$.
La propriété (6) implique que les projections $V \times_X V \to V$ sont des isomorphismes
au-dessus de $f^{-1}(Z) \subset V$. Ainsi, le schéma $V \times_X V$ est la réunion
des ouverts $\Delta_{V/X}(V)$ et $W \times_U W$, dont l'intersection
est $\Delta_{W/X}(W)$. On conclut qu'il existe un unique isomorphisme
$R \cong V \times_X V$ compatible avec les projections vers $V$.
Comme $V \to X$ et $V \to X'$ sont étales et surjectifs, on voit que
$$
X = V/ V \times_X V = V/R = V/V \times_{X'} V = X'
$$
d'après Espaces, lemme \ref{spaces-lemma-space-presentation}, ce qui conclut.
\end{proof}




\section{Applications}
\label{section-applications}

\noindent
Le lemme suivant se déduit aussi directement de
Espaces décents, lemme
\ref{decent-spaces-lemma-filter-quasi-compact-quasi-separated}
sans passer par l'approximation noethérienne absolue.

\begin{lemma}
\label{lemma-colimit-finitely-presented}
Soient $S$ un schéma et $X$ un espace algébrique quasi-compact et quasi-séparé
sur $S$. Tout $\mathcal{O}_X$-module quasi-cohérent est une
limite inductive filtrante de $\mathcal{O}_X$-modules de présentation finie.
\end{lemma}

\begin{proof}
On peut considérer $X$ comme un espace algébrique sur $\Spec(\mathbf{Z})$, voir
Espaces, définition \ref{spaces-definition-base-change}, et
Propriétés des espaces, définition \ref{spaces-properties-definition-separated}.
On peut donc appliquer la proposition \ref{proposition-approximate}
et écrire $X = \lim X_i$, où $X_i$ est de présentation finie sur $\mathbf{Z}$.
Ainsi, $X_i$ est un espace algébrique noethérien, voir
Morphismes des espaces, lemme
\ref{spaces-morphisms-lemma-finite-presentation-noetherian}.
Le morphisme $X \to X_i$ est affine, voir le
lemme \ref{lemma-directed-inverse-system-has-limit}.
On conclut par
Cohomologie des espaces, lemme
\ref{spaces-cohomology-lemma-direct-colimit-finite-presentation}.
\end{proof}

\noindent
Le reste de cette section consiste en des applications immédiates
du lemme \ref{lemma-colimit-finitely-presented}.

\begin{lemma}
\label{lemma-directed-colimit-finite-type}
Soient $S$ un schéma et $X$ un espace algébrique quasi-compact et quasi-séparé
sur $S$.
Soit $\mathcal{F}$ un $\mathcal{O}_X$-module quasi-cohérent.
Alors $\mathcal{F}$ est la limite inductive filtrante de ses sous-modules quasi-cohérents
de type fini.
\end{lemma}

\begin{proof}
Si $\mathcal{G}, \mathcal{H} \subset \mathcal{F}$ sont des
$\mathcal{O}_X$-sous-modules quasi-cohérents de type fini, alors l'image
de $\mathcal{G} \oplus \mathcal{H} \to \mathcal{F}$ est encore un
$\mathcal{O}_X$-sous-module quasi-cohérent de type fini qui les contient
tous deux. On voit ainsi que le système est filtrant.
Pour montrer que $\mathcal{F}$ est la limite inductive de ce système, écrivons
$\mathcal{F} = \colim_i \mathcal{F}_i$ comme
limite inductive filtrante de modules quasi-cohérents de présentation finie, comme dans le
lemme \ref{lemma-colimit-finitely-presented}.
Alors les images $\mathcal{G}_i = \Im(\mathcal{F}_i \to \mathcal{F})$ sont des
sous-modules quasi-cohérents de type fini de $\mathcal{F}$. Comme
$\mathcal{F}$ est leur limite inductive, le résultat en découle.
\end{proof}

\begin{lemma}
\label{lemma-finite-directed-colimit-surjective-maps}
Soient $S$ un schéma et $X$ un espace algébrique quasi-compact et quasi-séparé
sur $S$. Soit $\mathcal{F}$ un $\mathcal{O}_X$-module quasi-cohérent
de type fini. On peut alors écrire
$\mathcal{F} = \lim \mathcal{F}_i$, où chaque $\mathcal{F}_i$ est un
$\mathcal{O}_X$-module de présentation finie et tous les morphismes de transition
$\mathcal{F}_i \to \mathcal{F}_{i'}$ sont surjectifs.
\end{lemma}

\begin{proof}
Écrivons $\mathcal{F} = \colim \mathcal{G}_i$ comme limite inductive filtrante de
$\mathcal{O}_X$-modules de présentation finie
(lemme \ref{lemma-colimit-finitely-presented}).
Nous affirmons que $\mathcal{G}_i \to \mathcal{F}$ est surjectif pour un certain $i$.
En effet, choisissons une surjection étale $U \to X$, où $U$ est un schéma affine.
Choisissons un nombre fini de sections $s_k \in \mathcal{F}(U)$ qui engendrent
$\mathcal{F}|_U$. Comme $U$ est affine, on voit que chaque $s_k$ appartient à l'image
de $\mathcal{G}_i \to \mathcal{F}$ pour $i$ assez grand. Ainsi,
$\mathcal{G}_i \to \mathcal{F}$ est surjectif pour $i$ assez grand.
Fixons un tel $i$ et soit $\mathcal{K} \subset \mathcal{G}_i$ le
noyau de l'application $\mathcal{G}_i \to \mathcal{F}$. Écrivons
$\mathcal{K} = \colim \mathcal{K}_a$
comme limite inductive filtrante de ses sous-modules quasi-cohérents de type fini
(lemme \ref{lemma-directed-colimit-finite-type}). Alors
$\mathcal{F} = \colim \mathcal{G}_i/\mathcal{K}_a$ répond à la question
posée par le lemme.
\end{proof}

\noindent
Soit $X$ un espace algébrique. Dans le lemme suivant, on utilise la notion
d'une {\it $\mathcal{O}_X$-algèbre quasi-cohérente de présentation finie
$\mathcal{A}$}. Cela signifie que, pour tout schéma affine
$U = \Spec(R)$ étale sur $X$, on a $\mathcal{A}|_U = \widetilde{A}$,
où $A$ est une $R$-algèbre (commutative) de présentation finie
comme $R$-algèbre.

\begin{lemma}
\label{lemma-algebra-directed-colimit-finite-presentation}
Soient $S$ un schéma et $X$ un espace algébrique quasi-compact et quasi-séparé
sur $S$.
Soit $\mathcal{A}$ une $\mathcal{O}_X$-algèbre quasi-cohérente.
Alors $\mathcal{A}$ est une limite inductive filtrante d'$\mathcal{O}_X$-algèbres
quasi-cohérentes de présentation finie.
\end{lemma}

\begin{proof}
Écrivons d'abord $\mathcal{A} = \colim_i \mathcal{F}_i$ comme limite inductive filtrante
de modules quasi-cohérents de présentation finie, comme dans le
lemme \ref{lemma-colimit-finitely-presented}.
Pour chaque $i$, soit $\mathcal{B}_i = \text{Sym}(\mathcal{F}_i)$ l'algèbre
symétrique de $\mathcal{F}_i$ sur $\mathcal{O}_X$. Écrivons
$\mathcal{I}_i = \Ker(\mathcal{B}_i \to \mathcal{A})$. Écrivons
$\mathcal{I}_i = \colim_j \mathcal{F}_{i, j}$, où
$\mathcal{F}_{i, j}$ est un sous-module quasi-cohérent de type fini de
$\mathcal{I}_i$, voir le
lemme \ref{lemma-directed-colimit-finite-type}.
Posons $\mathcal{I}_{i, j} \subset \mathcal{I}_i$
égal à l'idéal de $\mathcal{B}_i$ engendré par $\mathcal{F}_{i, j}$.
Posons $\mathcal{A}_{i, j} = \mathcal{B}_i/\mathcal{I}_{i, j}$.
Alors $\mathcal{A}_{i, j}$ est une $\mathcal{O}_X$-algèbre quasi-cohérente
de présentation finie. Définissons $(i, j) \leq (i', j')$ si
$i \leq i'$ et si l'application $\mathcal{B}_i \to \mathcal{B}_{i'}$
envoie l'idéal $\mathcal{I}_{i, j}$ dans l'idéal $\mathcal{I}_{i', j'}$.
Alors il est clair que $\mathcal{A} = \colim_{i, j} \mathcal{A}_{i, j}$.
\end{proof}

\noindent
Soit $X$ un espace algébrique. Dans le lemme suivant, on utilise la notion
d'une {\it $\mathcal{O}_X$-algèbre quasi-cohérente $\mathcal{A}$
de type fini}. Cela signifie que, pour tout schéma affine
$U = \Spec(R)$ étale sur $X$, on a $\mathcal{A}|_U = \widetilde{A}$,
où $A$ est une $R$-algèbre (commutative) de type fini
comme $R$-algèbre.

\begin{lemma}
\label{lemma-algebra-directed-colimit-finite-type}
Soient $S$ un schéma et $X$ un espace algébrique quasi-compact et quasi-séparé
sur $S$. Soit $\mathcal{A}$ une $\mathcal{O}_X$-algèbre quasi-cohérente.
Alors $\mathcal{A}$ est la limite inductive filtrante de ses
$\mathcal{O}_X$-sous-algèbres quasi-cohérentes de type fini.
\end{lemma}

\begin{proof}
Omise. Indication : comparer avec la démonstration du
lemme \ref{lemma-directed-colimit-finite-type}.
\end{proof}

\noindent
Soit $X$ un espace algébrique. Dans le lemme suivant, on utilise la notion
d'une {\it $\mathcal{O}_X$-algèbre quasi-cohérente finie (resp.\ entière)
$\mathcal{A}$}. Cela signifie que, pour tout schéma
affine $U = \Spec(R)$ étale sur $X$, on a
$\mathcal{A}|_U = \widetilde{A}$, où $A$ est une $R$-algèbre (commutative)
finie (resp.\ entière) comme $R$-algèbre.

\begin{lemma}
\label{lemma-finite-algebra-directed-colimit-finite-finitely-presented}
Soient $S$ un schéma et $X$ un espace algébrique quasi-compact et quasi-séparé
sur $S$. Soit $\mathcal{A}$ une $\mathcal{O}_X$-algèbre quasi-cohérente
finie. Alors $\mathcal{A} = \colim \mathcal{A}_i$
est une limite inductive filtrante d'$\mathcal{O}_X$-algèbres quasi-cohérentes
finies et de présentation finie, à morphismes de transition surjectifs.
\end{lemma}

\begin{proof}
D'après le lemme \ref{lemma-finite-directed-colimit-surjective-maps},
il existe un $\mathcal{O}_X$-module de présentation finie
$\mathcal{F}$ et une surjection $\mathcal{F} \to \mathcal{A}$.
La structure d'algèbre fournit une surjection
$$
\text{Sym}^*_{\mathcal{O}_X}(\mathcal{F}) \longrightarrow \mathcal{A}
$$
Notons $\mathcal{J}$ le noyau. Écrivons $\mathcal{J} = \colim \mathcal{E}_i$
comme limite inductive filtrante de $\mathcal{O}_X$-sous-modules de type fini
$\mathcal{E}_i$ (lemme \ref{lemma-directed-colimit-finite-type}). Posons
$$
\mathcal{A}_i = \text{Sym}^*_{\mathcal{O}_X}(\mathcal{F})/(\mathcal{E}_i)
$$
où $(\mathcal{E}_i)$ désigne le faisceau d'idéaux engendré par
l'image de $\mathcal{E}_i \to \text{Sym}^*_{\mathcal{O}_X}(\mathcal{F})$.
Alors chaque $\mathcal{A}_i$ est une $\mathcal{O}_X$-algèbre de présentation finie,
les morphismes de transition sont surjectifs et $\mathcal{A} = \colim \mathcal{A}_i$.
Pour achever la démonstration, il reste
à montrer que $\mathcal{A}_i$ est une $\mathcal{O}_X$-algèbre finie
pour $i$ assez grand. Pour cela, choisissons un morphisme étale surjectif
$U \to X$, où $U$ est un schéma affine. Prenons des générateurs
$f_1, \ldots, f_m \in \Gamma(U, \mathcal{F})$.
Comme $\mathcal{A}(U)$ est une $\mathcal{O}_X(U)$-algèbre finie,
on voit que, pour chaque $j$, il existe un polynôme unitaire
$P_j \in \mathcal{O}(U)[T]$ tel que $P_j(f_j)$ soit nul dans $\mathcal{A}(U)$.
Comme $\mathcal{A} = \colim \mathcal{A}_i$ par construction,
on a $P_j(f_j) = 0$ dans $\mathcal{A}_i(U)$ pour tout $i$ assez grand.
Pour un tel $i$, les algèbres $\mathcal{A}_i$ sont finies.
\end{proof}

\begin{lemma}
\label{lemma-integral-algebra-directed-colimit-finite}
Soient $S$ un schéma et $X$ un espace algébrique quasi-compact et quasi-séparé
sur $S$. Soit $\mathcal{A}$ une $\mathcal{O}_X$-algèbre quasi-cohérente
entière. Alors
\begin{enumerate}
\item $\mathcal{A}$ est la limite inductive filtrante de ses
$\mathcal{O}_X$-sous-algèbres quasi-cohérentes finies, et
\item $\mathcal{A}$ est une limite inductive filtrante d'$\mathcal{O}_X$-algèbres
finies et de présentation finie.
\end{enumerate}
\end{lemma}

\begin{proof}
D'après le lemme \ref{lemma-algebra-directed-colimit-finite-type}, on a
$\mathcal{A} = \colim \mathcal{A}_i$, où
$\mathcal{A}_i \subset \mathcal{A}$ parcourt les
$\mathcal{O}_X$-sous-algèbres quasi-cohérentes de type fini.
Toute $\mathcal{O}_X$-sous-algèbre quasi-cohérente de type fini
de $\mathcal{A}$ est finie (utiliser Algèbre, lemme
\ref{algebra-lemma-characterize-finite-in-terms-of-integral},
sur les schémas affines étales sur $X$). Cela prouve (1).

\medskip\noindent
Pour prouver (2), écrivons $\mathcal{A} = \colim \mathcal{F}_i$
comme limite inductive de $\mathcal{O}_X$-modules de présentation finie au moyen du
lemme \ref{lemma-colimit-finitely-presented}.
Pour chaque $i$, soit $\mathcal{J}_i$ le noyau de l'application
$$
\text{Sym}^*_{\mathcal{O}_X}(\mathcal{F}_i) \longrightarrow \mathcal{A}
$$
Pour $i' \geq i$, il existe une application induite $\mathcal{J}_i \to \mathcal{J}_{i'}$,
et l'on a $\mathcal{A} =
\colim \text{Sym}^*_{\mathcal{O}_X}(\mathcal{F}_i)/\mathcal{J}_i$.
En outre, les $\mathcal{O}_X$-algèbres quasi-cohérentes
$\text{Sym}^*_{\mathcal{O}_X}(\mathcal{F}_i)/\mathcal{J}_i$
sont finies (voir ci-dessus). Écrivons $\mathcal{J}_i = \colim \mathcal{E}_{ik}$
comme limite inductive de $\mathcal{O}_X$-modules de présentation finie.
Étant donnés $i' \geq i$ et $k$, il existe un $k'$ tel qu'on
dispose d'une application $\mathcal{E}_{ik} \to \mathcal{E}_{i'k'}$
rendant
$$
\xymatrix{
\mathcal{J}_i \ar[r] & \mathcal{J}_{i'} \\
\mathcal{E}_{ik} \ar[u] \ar[r] & \mathcal{E}_{i'k'} \ar[u]
}
$$
le diagramme commutatif. Cela résulte de
Cohomologie des espaces, lemme
\ref{spaces-cohomology-lemma-finite-presentation-quasi-compact-colimit}.
Cela induit une application
$$
\mathcal{A}_{ik} =
\text{Sym}^*_{\mathcal{O}_X}(\mathcal{F}_i)/(\mathcal{E}_{ik})
\longrightarrow
\text{Sym}^*_{\mathcal{O}_X}(\mathcal{F}_{i'})/(\mathcal{E}_{i'k'}) =
\mathcal{A}_{i'k'}
$$
où $(\mathcal{E}_{ik})$ désigne l'idéal engendré par $\mathcal{E}_{ik}$.
Les $\mathcal{O}_X$-algèbres quasi-cohérentes $\mathcal{A}_{ki}$
sont de présentation finie et finies pour $k$ assez grand
(voir la démonstration du
lemme \ref{lemma-finite-algebra-directed-colimit-finite-finitely-presented}).
Enfin, on a
$$
\colim \mathcal{A}_{ik} = \colim \mathcal{A}_i = \mathcal{A}
$$
En effet, la première égalité a été démontrée dans la démonstration du
lemme \ref{lemma-finite-algebra-directed-colimit-finite-finitely-presented},
et la seconde parce que $\mathcal{A}$ est la limite inductive des
modules $\mathcal{F}_i$.
\end{proof}

\begin{lemma}
\label{lemma-extend}
Soit $S$ un schéma.
Soit $X$ un espace algébrique quasi-compact et quasi-séparé sur $S$.
Soit $U \subset X$ un sous-espace ouvert quasi-compact.
Soit $\mathcal{F}$ un $\mathcal{O}_X$-module quasi-cohérent.
Soit $\mathcal{G} \subset \mathcal{F}|_U$ un
$\mathcal{O}_U$-sous-module quasi-cohérent de type fini. Alors
il existe un sous-module quasi-cohérent $\mathcal{G}' \subset \mathcal{F}$
de type fini tel que $\mathcal{G}'|_U = \mathcal{G}$.
\end{lemma}

\begin{proof}
Notons $j : U \to X$ le morphisme d'inclusion. Comme $X$ est quasi-séparé
et $U$ quasi-compact, le morphisme $j$ est quasi-compact. Ainsi,
$j_*\mathcal{G} \subset j_*\mathcal{F}|_U$ sont des modules quasi-cohérents
sur $X$ (Morphismes des espaces, lemme
\ref{spaces-morphisms-lemma-pushforward}).
Posons $\mathcal{H} =
\Ker(j_*\mathcal{G} \oplus \mathcal{F} \to j_*\mathcal{F}|_U)$.
Alors $\mathcal{H}|_U = \mathcal{G}$. D'après le
lemme \ref{lemma-directed-colimit-finite-type},
on peut trouver un sous-module quasi-cohérent de type fini
$\mathcal{H}' \subset \mathcal{H}$ tel que
$\mathcal{H}'|_U = \mathcal{H}|_U = \mathcal{G}$.
Posons $\mathcal{G}' = \Im(\mathcal{H}' \to \mathcal{F})$
pour conclure.
\end{proof}









\section{Approximation relative}
\label{section-relative-approximation}

\noindent
On examine des variantes de la proposition \ref{proposition-approximate}
sur une base.

\begin{lemma}
\label{lemma-approximate-morphism}
Soit $f : X \to Y$ un morphisme d'espaces algébriques quasi-compacts et quasi-séparés
sur $\mathbf{Z}$.
Il existe un ensemble préordonné filtrant $I$ et un système projectif $(f_i : X_i \to Y_i)$
de morphismes d'espaces algébriques indexé par $I$, tel que les morphismes de transition
$X_i \to X_{i'}$ et $Y_i \to Y_{i'}$ soient affines, que $X_i$
et $Y_i$ soient quasi-séparés et de type fini sur $\mathbf{Z}$, et
que $(X \to Y) = \lim (X_i \to Y_i)$.
\end{lemma}

\begin{proof}
Écrivons $X = \lim_{a \in A} X_a$ et $Y = \lim_{b \in B} Y_b$ comme dans la
proposition \ref{proposition-approximate}, c'est-à-dire avec $X_a$ et $Y_b$
quasi-séparés et de type fini sur $\mathbf{Z}$ et
des morphismes de transition affines.

\medskip\noindent
Fixons $b \in B$.
D'après le lemme \ref{lemma-better-characterize-relative-limit-preserving}
appliqué à $Y_b$ et à $X = \lim X_a$ sur $\mathbf{Z}$,
il existe un $a \in A$ et un morphisme
$f_{a, b} : X_a \to Y_b$ rendant le diagramme
$$
\xymatrix{
X \ar[d] \ar[r] & Y \ar[d] \\
X_a \ar[r] & Y_b
}
$$
commutatif. Soit $I$ l'ensemble des triplets $(a, b, f_{a, b})$ qu'on
obtient ainsi.

\medskip\noindent
Soient $(a, b, f_{a, b})$ et $(a', b', f_{a', b'})$ dans $I$.
Soit $b'' \leq \min(b, b')$. D'après le
lemme \ref{lemma-better-characterize-relative-limit-preserving},
il existe encore un $a'' \geq \max(a, a')$ tel que
les composés $X_{a''} \to X_a \to Y_b \to Y_{b''}$ et
$X_{a''} \to X_{a'} \to Y_{b'} \to Y_{b''}$ soient égaux.
Munissons $I$ du préordre
$$
(a, b, f_{a, b}) \geq (a', b', f_{a', b'})
\Leftrightarrow
a \geq a',\ b \geq b',\text{ et }
g_{b, b'} \circ f_{a, b} = f_{a', b'} \circ h_{a, a'}
$$
où $h_{a, a'} : X_a \to X_{a'}$ et $g_{b, b'} : Y_b \to Y_{b'}$
sont les morphismes de transition. Les remarques précédentes montrent que $I$
est filtrant et que les applications $I \to A$, $(a, b, f_{a, b}) \mapsto a$,
et $I \to B$, $(a, b, f_{a, b})$, sont cofinales. Si, pour $i = (a, b, f_{a, b})$,
on pose $X_i = X_a$, $Y_i = Y_b$ et $f_i = f_{a, b}$, on obtient
un système projectif de morphismes indexé par $I$ et l'on a
$$
\lim_{i \in I} X_i = \lim_{a \in A} X_a = X
\quad\text{et}\quad
\lim_{i \in I} S_i = \lim_{b \in B} Y_b = Y
$$
d'après Catégories, lemme \ref{categories-lemma-initial} (rappelons que
les limites sur $I$ sont en réalité les limites sur la catégorie opposée
associée à $I$, de sorte que cofinal devient initial).
Cela achève la démonstration.
\end{proof}

\begin{lemma}
\label{lemma-relative-approximation}
Soient $S$ un schéma et $f : X \to Y$ un morphisme d'espaces algébriques
sur $S$. Supposons que
\begin{enumerate}
\item $X$ soit quasi-compact et quasi-séparé, et
\item $Y$ soit quasi-séparé.
\end{enumerate}
Alors $X = \lim X_i$ est la limite d'un système projectif filtrant d'espaces algébriques
$X_i$ de présentation finie sur $Y$, à morphismes de transition affines
sur $Y$.
\end{lemma}

\begin{proof}
Comme $|f|(|X|)$ est quasi-compact, on peut remplacer $Y$ par un sous-espace
ouvert quasi-compact dont l'ensemble des points contient $|f|(|X|)$. On peut donc supposer
$Y$ également quasi-compact.
D'après le lemme \ref{lemma-approximate-morphism}, on peut écrire
$(X \to Y) = \lim (X_i \to Y_i)$ pour un certain système projectif filtrant
de morphismes d'espaces algébriques de type fini sur $\mathbf{Z}$, à
morphismes de transition affines. Comme les limites commutent aux limites
(Catégories, lemme \ref{categories-lemma-colimits-commute}),
on a $X = \lim X_i \times_{Y_i} Y$. Pour $i \geq i'$,
le morphisme de transition
$X_i \times_{Y_i} Y \to X_{i'} \times_{Y_{i'}} Y$ est affine, car c'est
le composé
$$
X_i \times_{Y_i} Y \to X_i \times_{Y_{i'}} Y \to X_{i'} \times_{Y_{i'}} Y
$$
où le premier morphisme est une immersion fermée (d'après
Morphismes des espaces, lemme
\ref{spaces-morphisms-lemma-fibre-product-after-map})
et le second est un changement de base d'un morphisme affine
(Morphismes des espaces, lemme \ref{spaces-morphisms-lemma-base-change-affine}),
et où le composé de morphismes affines est affine
(Morphismes des espaces, lemme \ref{spaces-morphisms-lemma-composition-affine}).
Les morphismes $f_i$ sont de présentation finie
(Morphismes des espaces, lemmes
\ref{spaces-morphisms-lemma-noetherian-finite-type-finite-presentation} et
\ref{spaces-morphisms-lemma-finite-presentation-permanence})
et leurs changements de base $X_i \times_{f_i, Y_i} Y \to Y$
sont donc de présentation finie
(Morphismes des espaces, lemme
\ref{spaces-morphisms-lemma-base-change-finite-presentation}).
\end{proof}






\section{Réalisation d'un morphisme de type fini comme fermé dans un morphisme de présentation finie}
\sectionmark{Type fini comme fermé dans une présentation finie}
\label{section-finite-type-closed-in-finite-presentation}

\noindent
Cette section est l'analogue de la
section \ref{limits-section-finite-type-closed-in-finite-presentation} de Limites.

\begin{lemma}
\label{lemma-affine-morphism-is-limit}
Soient $S$ un schéma et $f : X \to Y$ un morphisme affine d'espaces
algébriques sur $S$. Si $Y$ est quasi-compact et
quasi-séparé, alors $X$ est une limite projective filtrante $X = \lim X_i$,
où chaque $X_i$ est affine et de présentation finie sur $Y$.
\end{lemma}

\begin{proof}
Considérons le $\mathcal{O}_Y$-module quasi-cohérent
$\mathcal{A} = f_*\mathcal{O}_X$. D'après le
lemme \ref{lemma-algebra-directed-colimit-finite-presentation},
on peut écrire $\mathcal{A} = \colim \mathcal{A}_i$ comme une limite
inductive filtrante d'algèbres de présentation finie
sur $\mathcal{O}_Y$, notées $\mathcal{A}_i$.
Posons $X_i = \underline{\Spec}_Y(\mathcal{A}_i)$, voir
Morphismes des espaces, définition
\ref{spaces-morphisms-definition-relative-spec}.
Par construction, $X_i \to Y$ est affine et de présentation finie
et $X = \lim X_i$.
\end{proof}

\begin{lemma}
\label{lemma-integral-limit-finite-and-finite-presentation}
Soient $S$ un schéma et $f : X \to Y$ un morphisme entier d'espaces
algébriques sur $S$. Supposons $Y$ quasi-compact et quasi-séparé.
Alors $X$ s'écrit comme une limite projective filtrante $X = \lim X_i$,
où les $X_i$ sont finis et de présentation finie sur $Y$.
\end{lemma}

\begin{proof}
Considérons le $\mathcal{O}_Y$-module quasi-cohérent
$\mathcal{A} = f_*\mathcal{O}_X$. D'après le
lemme \ref{lemma-integral-algebra-directed-colimit-finite},
on peut écrire $\mathcal{A} = \colim \mathcal{A}_i$ comme une limite
inductive filtrante de $\mathcal{O}_Y$-algèbres
$\mathcal{A}_i$ finies et de présentation finie.
Posons $X_i = \underline{\Spec}_Y(\mathcal{A}_i)$, voir
Morphismes des espaces, définition
\ref{spaces-morphisms-definition-relative-spec}.
Par construction, $X_i \to Y$ est fini et de présentation finie et
$X = \lim X_i$.
\end{proof}

\begin{lemma}
\label{lemma-finite-in-finite-and-finite-presentation}
Soient $S$ un schéma et $f : X \to Y$ un morphisme fini d'espaces
algébriques sur $S$. Supposons $Y$ quasi-compact et quasi-séparé.
Alors $X$ s'écrit comme une limite projective filtrante $X = \lim X_i$,
où les morphismes de transition sont des immersions fermées et les objets
$X_i$ sont finis et de présentation finie sur $Y$.
\end{lemma}

\begin{proof}
Considérons le $\mathcal{O}_Y$-module quasi-cohérent de type fini
$\mathcal{A} = f_*\mathcal{O}_X$. D'après le
lemme \ref{lemma-finite-algebra-directed-colimit-finite-finitely-presented},
on peut écrire $\mathcal{A} = \colim \mathcal{A}_i$ comme une limite
inductive filtrante de $\mathcal{O}_Y$-algèbres
$\mathcal{A}_i$ finies et de présentation finie, à morphismes de transition surjectifs.
Posons $X_i = \underline{\Spec}_Y(\mathcal{A}_i)$, voir
Morphismes des espaces, définition
\ref{spaces-morphisms-definition-relative-spec}.
Par construction, $X_i \to Y$ est fini et de présentation finie,
les morphismes de transition sont des immersions fermées et $X = \lim X_i$.
\end{proof}

\begin{lemma}
\label{lemma-closed-is-limit-closed-and-finite-presentation}
\begin{slogan}
Les immersions fermées d'espaces algébriques qcqs peuvent être approchées
par des immersions fermées de présentation finie.
\end{slogan}
Soient $S$ un schéma et $f : X \to Y$ une immersion fermée d'espaces
algébriques sur $S$. Supposons $Y$ quasi-compact et quasi-séparé.
Alors $X$ s'écrit comme une limite projective filtrante $X = \lim X_i$,
où les morphismes de transition sont des immersions fermées et les morphismes
$X_i \to Y$ sont des immersions fermées de présentation finie.
\end{lemma}

\begin{proof}
Soit $\mathcal{I} \subset \mathcal{O}_Y$ le faisceau quasi-cohérent
d'idéaux définissant $X$ comme sous-espace fermé de $Y$. D'après le
lemme \ref{lemma-directed-colimit-finite-type},
on peut écrire $\mathcal{I} = \colim \mathcal{I}_i$ comme la
limite inductive filtrante de ses sous-modules quasi-cohérents de type fini.
Soit $X_i$ le sous-espace fermé de $X$ défini par $\mathcal{I}_i$.
Alors $X_i \to Y$ est une immersion fermée de présentation finie
et $X = \lim X_i$. Certains détails sont omis.
\end{proof}

\begin{lemma}
\label{lemma-quasi-affine-closed-in-finite-presentation}
Soient $S$ un schéma et $f : X \to Y$ un morphisme d'espaces algébriques
sur $S$. Supposons que
\begin{enumerate}
\item $f$ soit localement de type fini et quasi-affine, et
\item $Y$ soit quasi-compact et quasi-séparé.
\end{enumerate}
Alors il existe un morphisme de présentation finie
$f' : X' \to Y$ et une immersion fermée $X \to X'$ sur $Y$.
\end{lemma}

\begin{proof}
D'après Morphismes des espaces, lemme
\ref{spaces-morphisms-lemma-characterize-quasi-affine},
il existe une factorisation $X \to Z \to Y$ où
$X \to Z$ est une immersion ouverte quasi-compacte et
$Z \to Y$ est affine. Écrivons $Z = \lim Z_i$, où $Z_i$ est affine et
de présentation finie sur $Y$ (lemme \ref{lemma-affine-morphism-is-limit}).
Pour un certain $0 \in I$, il existe un ouvert quasi-compact $U_0 \subset Z_0$
tel que $X$ soit isomorphe à l'image réciproque de $U_0$ dans $Z$
(lemme \ref{lemma-descend-opens}). Soit $U_i$ l'image réciproque de
$U_0$ dans $Z_i$, de sorte que $U = \lim U_i$. D'après le
lemme \ref{lemma-finite-type-eventually-closed},
le morphisme $X \to U_i$ est une immersion fermée pour un certain $i$ assez grand.
En posant $X' = U_i$, on conclut.
\end{proof}

\begin{lemma}
\label{lemma-finite-type-closed-in-finite-presentation}
Soient $S$ un schéma et $f : X \to Y$ un morphisme d'espaces algébriques
sur $S$. Supposons :
\begin{enumerate}
\item $f$ soit localement de type fini,
\item $X$ soit quasi-compact et quasi-séparé, et
\item $Y$ soit quasi-compact et quasi-séparé.
\end{enumerate}
Alors il existe un morphisme de présentation finie
$f' : X' \to Y$ et une immersion fermée $X \to X'$ d'espaces
algébriques sur $Y$.
\end{lemma}

\begin{proof}
D'après la proposition \ref{proposition-approximate},
on peut écrire $X = \lim_i X_i$, où $X_i$ est quasi-séparé et de type fini sur
$\mathbf{Z}$ et les morphismes de transition $f_{ii'} : X_i \to X_{i'}$ sont affines.
Considérons le diagramme commutatif
$$
\xymatrix{
X \ar[r] \ar[rd] & X_{i, Y} \ar[r] \ar[d] & X_i \ar[d] \\
& Y \ar[r] & \Spec(\mathbf{Z})
}
$$
Notons que $X_i$ est de présentation finie sur $\Spec(\mathbf{Z})$, voir
Morphismes des espaces, lemme
\ref{spaces-morphisms-lemma-noetherian-finite-type-finite-presentation}.
Le changement de base $X_{i, Y} \to Y$ est donc de présentation finie d'après
Morphismes des espaces, lemme
\ref{spaces-morphisms-lemma-base-change-finite-presentation}.
Observons que $\lim X_{i, Y} = X \times Y$ et que $X \to X \times Y$ est un
monomorphisme. D'après le lemme \ref{lemma-finite-type-eventually-closed},
le morphisme $X \to X_{i, Y}$ est un monomorphisme pour $i$ assez grand.
Fixons un tel $i$. Notons que $X \to X_{i, Y}$ est localement de type fini
(Morphismes des espaces, lemme
\ref{spaces-morphisms-lemma-permanence-finite-type})
et est un monomorphisme, donc est séparé et localement quasi-fini
(Morphismes des espaces, lemme
\ref{spaces-morphisms-lemma-monomorphism-loc-finite-type-loc-quasi-finite}).
Le morphisme $X \to X_{i, Y}$ est donc représentable.
Le morphisme $X \to X_{i, Y}$ est alors quasi-affine, car on peut appliquer le
principe énoncé dans Espaces, lemme
\ref{spaces-lemma-representable-transformations-property-implication},
et le résultat pour les morphismes de schémas de Compléments sur les morphismes, lemme
\ref{more-morphisms-lemma-quasi-finite-separated-quasi-affine}.
Ainsi, le lemme \ref{lemma-quasi-affine-closed-in-finite-presentation}
fournit une factorisation $X \to X' \to X_{i, Y}$,
où $X \to X'$ est une immersion fermée et $X' \to X_{i, Y}$ est de
présentation finie. Enfin, $X' \to Y$ est de présentation finie comme
composé de morphismes de présentation finie
(Morphismes des espaces, lemme
\ref{spaces-morphisms-lemma-composition-finite-presentation}).
\end{proof}

\begin{proposition}
\label{proposition-separated-closed-in-finite-presentation}
Soient $S$ un schéma et $f : X \to Y$ un morphisme d'espaces algébriques
sur $S$. Supposons que
\begin{enumerate}
\item $f$ soit de type fini et séparé, et
\item $Y$ soit quasi-compact et quasi-séparé.
\end{enumerate}
Alors il existe un morphisme séparé de présentation finie
$f' : X' \to Y$ et une immersion fermée $X \to X'$ sur $Y$.
\end{proposition}

\begin{proof}
D'après le lemme \ref{lemma-finite-type-closed-in-finite-presentation},
il existe une immersion fermée $X \to Z$, où $Z/Y$ est de
présentation finie. Soit $\mathcal{I} \subset \mathcal{O}_Z$
le faisceau quasi-cohérent d'idéaux définissant $X$ comme
sous-espace fermé de $Y$. D'après le
lemme \ref{lemma-directed-colimit-finite-type},
on peut écrire $\mathcal{I}$ comme une limite inductive filtrante
$\mathcal{I} = \colim_{a \in A} \mathcal{I}_a$ de ses
faisceaux quasi-cohérents d'idéaux de type fini.
Soit $X_a \subset Z$ le sous-espace fermé défini par $\mathcal{I}_a$.
Ceux-ci forment un système projectif indexé par $A$.
Les morphismes de transition $X_a \to X_{a'}$ sont affines, car
ce sont des immersions fermées. Chaque $X_a$ est quasi-compact et quasi-séparé,
puisqu'il est un sous-espace fermé de $Z$ et que $Z$ est quasi-compact et
quasi-séparé d'après nos hypothèses.
On a $X = \lim_a X_a$, comme il résulte directement de
l'égalité $\mathcal{I} = \colim_{a \in A} \mathcal{I}_a$.
Chacun des morphismes $X_a \to Z$ est de présentation finie, voir
Morphismes, lemme \ref{morphisms-lemma-closed-immersion-finite-presentation}.
Par conséquent, les morphismes $X_a \to Y$ sont de présentation finie.
Il suffit donc de montrer que $X_a \to Y$ est séparé pour un certain
$a \in A$. Cela résulte du lemme \ref{lemma-eventually-separated},
car on a supposé que $X \to Y$ était séparé.
\end{proof}






\section{Approximation des morphismes propres}
\label{section-approximate-proper}

\begin{lemma}
\label{lemma-proper-limit-of-proper-finite-presentation}
Soient $S$ un schéma et $f : X \to Y$ un morphisme propre d'espaces
algébriques sur $S$, avec $Y$ quasi-compact et quasi-séparé. Alors
$X = \lim X_i$ est une limite projective filtrante d'espaces algébriques $X_i$
propres et de présentation finie sur $Y$, dont les morphismes de transition
et les morphismes $X \to X_i$ sont des immersions fermées.
\end{lemma}

\begin{proof}
D'après la proposition \ref{proposition-separated-closed-in-finite-presentation},
il existe une immersion fermée $X \to X'$, où $X'$ est séparé et de
présentation finie sur $Y$. D'après le
lemme \ref{lemma-closed-is-limit-closed-and-finite-presentation},
on peut écrire $X = \lim X_i$, où $X_i \to X'$ est une immersion fermée de
présentation finie. Nous affirmons que, pour tout $i$ assez grand,
le morphisme $X_i \to Y$ est propre, ce qui achèvera la démonstration.

\medskip\noindent
Pour le montrer, on peut supposer que $Y$ est un schéma affine, voir
Morphismes des espaces, lemme \ref{spaces-morphisms-lemma-proper-local}.
Ensuite, on applique la version faible du lemme de Chow, voir
Cohomologie des espaces, lemme \ref{spaces-cohomology-lemma-weak-chow},
pour obtenir un diagramme
$$
\xymatrix{
X' \ar[rd] & X'' \ar[d] \ar[l]^\pi \ar[r] & \mathbf{P}^n_Y \ar[dl] \\
& Y &
}
$$
où $X'' \to \mathbf{P}^n_Y$ est une immersion et
$\pi : X'' \to X'$ est propre et surjectif. Notons
$X'_i \subset X''$, resp.\ $\pi^{-1}(X)$, l'image réciproque schématique de
$X_i \subset X'$, resp.\ $X \subset X'$.
Alors $\lim X'_i = \pi^{-1}(X)$. Comme $\pi^{-1}(X) \to Y$ est propre
(Morphismes des espaces, lemme \ref{spaces-morphisms-lemma-composition-proper}),
on voit que $\pi^{-1}(X) \to \mathbf{P}^n_Y$ est une immersion fermée
(Morphismes des espaces, lemmes
\ref{spaces-morphisms-lemma-universally-closed-permanence} et
\ref{spaces-morphisms-lemma-immersion-when-closed}).
Ainsi, pour $i$ assez grand,
$X'_i \to \mathbf{P}^n_Y$ est une immersion fermée d'après le
lemme \ref{lemma-eventually-closed-immersion}.
Par conséquent, $X'_i$ est propre sur $Y$.
Pour un tel $i$, le morphisme $X_i \to Y$ est propre d'après
Morphismes des espaces, lemme \ref{spaces-morphisms-lemma-image-proper-is-proper}.
\end{proof}

\begin{lemma}
\label{lemma-proper-limit-of-proper-finite-presentation-noetherian}
Soit $f : X \to Y$ un morphisme propre d'espaces algébriques sur $\mathbf{Z}$,
avec $Y$ quasi-compact et quasi-séparé. Il existe alors un ensemble préordonné
filtrant $I$, un système projectif $(f_i : X_i \to Y_i)$ de morphismes d'espaces
algébriques indexé par $I$, tel que les morphismes de transition $X_i \to X_{i'}$
et $Y_i \to Y_{i'}$ soient affines, que $f_i$ soit propre et de
présentation finie, que $Y_i$ soit de présentation finie sur
$\mathbf{Z}$ et que $(X \to Y) = \lim (X_i \to Y_i)$.
\end{lemma}

\begin{proof}
D'après le lemme \ref{lemma-proper-limit-of-proper-finite-presentation},
on peut écrire $X = \lim_{k \in K} X_k$, avec $X_k \to Y$ propre et
de présentation finie. Ensuite, par approximation noethérienne absolue
(proposition \ref{proposition-approximate}), on peut
écrire $Y = \lim_{j \in J} Y_j$, avec $Y_j$ de présentation finie
sur $\mathbf{Z}$.
Pour chaque $k$, il existe un $j$ et un morphisme $X_{k, j} \to Y_j$
de présentation finie tels que $X_k \cong Y \times_{Y_j} X_{k, j}$
comme espaces algébriques sur $Y$, voir le
lemme \ref{lemma-descend-finite-presentation}.
Quitte à augmenter $j$, on peut supposer $X_{k, j} \to Y_j$ propre, voir le
lemme \ref{lemma-eventually-proper}. L'ensemble $I$ sera constitué
de ces couples $(k, j)$ et le morphisme correspondant sera $X_{k, j} \to Y_j$.
Pour tout $k' \geq k$, on peut trouver un $j' \geq j$ et un morphisme
$X_{j', k'} \to X_{j, k}$ sur $Y_{j'} \to Y_j$ dont le changement de base à $Y$
donne le morphisme $X_{k'} \to X_k$ (cela résulte encore du
lemme \ref{lemma-descend-finite-presentation}).
Ces morphismes constituent les morphismes de transition du système. Certains détails
sont omis.
\end{proof}

\noindent
Rappelons le support schématique d'un
module quasi-cohérent de type fini, voir
Morphismes des espaces, définition
\ref{spaces-morphisms-definition-scheme-theoretic-support}.

\begin{lemma}
\label{lemma-eventually-proper-support}
Hypothèses et notation comme dans la situation \ref{situation-descent-property}.
Soit $\mathcal{F}_0$ un $\mathcal{O}_{X_0}$-module quasi-cohérent.
Notons $\mathcal{F}$ et $\mathcal{F}_i$ les images réciproques de
$\mathcal{F}_0$ sur $X$ et $X_i$. Supposons que
\begin{enumerate}
\item $f_0$ soit localement de type fini,
\item $\mathcal{F}_0$ soit de type fini,
\item le support schématique de $\mathcal{F}$ soit propre sur $Y$.
\end{enumerate}
Alors le support schématique de $\mathcal{F}_i$ est propre sur $Y_i$
pour un certain $i$.
\end{lemma}

\begin{proof}
On peut remplacer $X_0$ par le support schématique de $\mathcal{F}_0$.
D'après Morphismes des espaces, lemme \ref{spaces-morphisms-lemma-support-finite-type},
cela garantit que $X_i$ est le support de $\mathcal{F}_i$ et que $X$ est le
support de $\mathcal{F}$. Si $Z \subset X$ désigne alors le support
schématique de $\mathcal{F}$, on voit que $Z \to X$ est un homéomorphisme
universel. On en conclut que $X \to Y$ est propre, puisque c'est le cas de
$Z \to Y$ par hypothèse, voir
Morphismes, lemme \ref{morphisms-lemma-image-proper-is-proper}.
D'après le lemme \ref{lemma-eventually-proper}, le morphisme $X_i \to Y$ est propre
pour un certain $i$. Il s'ensuit que le support schématique $Z_i$ de
$\mathcal{F}_i$ est propre sur $Y$ d'après
Morphismes des espaces, lemmes
\ref{spaces-morphisms-lemma-closed-immersion-proper} et
\ref{spaces-morphisms-lemma-composition-proper}.
\end{proof}








\section{Immersion dans un espace affine}
\label{section-embedding}

\noindent
Quelques lemmes techniques qui serviront plus loin à démontrer le lemme de Chow.

\begin{lemma}
\label{lemma-embedding-into-affine-over-ls-qs}
Soient $S$ un schéma et $f : U \to X$ un morphisme d'espaces
algébriques sur $S$. Supposons que $U$ soit un schéma affine, que $f$ soit
localement de type fini, et que $X$ soit quasi-séparé et localement séparé.
Alors il existe une immersion $U \to \mathbf{A}^n_X$ sur $X$.
\end{lemma}

\begin{proof}
Écrivons $U = \Spec(A)$. Écrivons $A = \colim A_i$ comme limite inductive filtrante
de sous-algèbres de type fini sur $\mathbf{Z}$. Pour tout $i$, le morphisme
$U \to U_i = \Spec(A_i)$ induit un morphisme
$$
U \longrightarrow X \times U_i
$$
sur $X$. À la limite, le morphisme $U \to X \times U$ est une immersion
puisque $X$ est localement séparé, voir
Morphismes des espaces, lemme
\ref{spaces-morphisms-lemma-semi-diagonal}.
D'après le lemme \ref{lemma-finite-type-eventually-closed},
le morphisme $U \to X \times U_i$ est une immersion pour un certain $i$.
Comme $U_i$ est isomorphe à un sous-schéma fermé de
$\mathbf{A}^n_{\mathbf{Z}}$, le lemme en résulte.
\end{proof}

\begin{remark}
\label{remark-cannot-embed-in-general}
Nous avons vu dans Exemples, section \ref{examples-section-embedding-affines},
que le lemme \ref{lemma-embedding-into-affine-over-ls-qs}
est faux si l'on omet l'hypothèse que $X$ soit localement séparé.
Cela soulève la question suivante : le
lemme \ref{lemma-embedding-into-affine-over-ls-qs}
reste-t-il vrai si l'on omet l'hypothèse que $X$ soit quasi-séparé ?
Si vous connaissez la réponse, merci d'écrire à
\href{mailto:stacks.project@gmail.com}{stacks.project@gmail.com}.
\end{remark}

\begin{lemma}
\label{lemma-embedding-into-affine-over-qs}
Soient $S$ un schéma et $f : Y \to X$ un morphisme d'espaces
algébriques sur $S$. Supposons que $X$ soit noethérien et que $f$ soit de présentation finie.
Alors il existe un ouvert dense $V \subset Y$ et une immersion
$V \to \mathbf{A}^n_X$.
\end{lemma}

\begin{proof}
Les hypothèses impliquent que $Y$ est noethérien
(Morphismes des espaces, lemme
\ref{spaces-morphisms-lemma-finite-presentation-noetherian}).
Alors $Y$ est quasi-séparé et possède donc un sous-schéma ouvert dense
(Propriétés des espaces, proposition
\ref{spaces-properties-proposition-locally-quasi-separated-open-dense-scheme}).
On peut donc supposer que $Y$ est un schéma noethérien.
En retirant les intersections des composantes irréductibles de $Y$
(utiliser Topologie, lemme \ref{topology-lemma-Noetherian} et
Propriétés, lemme \ref{properties-lemma-Noetherian-topology}),
on peut supposer que $Y$ est une réunion disjointe de schémas
noethériens irréductibles. Comme il existe une immersion
$$
\mathbf{A}^n_X \amalg \mathbf{A}^m_X
\longrightarrow
\mathbf{A}^{\max(n, m) + 1}_X
$$
(détails omis), il suffit de démontrer le résultat lorsque
$Y$ est irréductible.

\medskip\noindent
Supposons que $Y$ soit un schéma irréductible. Soit $T \subset |X|$ l'adhérence de
l'image de $f : Y \to X$. Comme $|Y|$ et $|X|$ sont des espaces topologiques sobres
(Propriétés des espaces, lemme
\ref{spaces-properties-lemma-quasi-separated-sober})
$T$ est irréductible et possède un unique point générique $\xi$, lequel est
l'image du point générique $\eta$ de $Y$.
Soit $\mathcal{I} \subset X$ un faisceau quasi-cohérent d'idéaux
définissant sur $T$ la structure d'espace réduit induite
(Propriétés des espaces, définition
\ref{spaces-properties-definition-reduced-induced-space}).
Comme $\mathcal{O}_{Y, \eta}$ est un anneau local artinien, il existe
$n > 0$ tel que $f^{-1}\mathcal{I}^n \mathcal{O}_{Y, \eta} = 0$.
Comme $f^{-1}\mathcal{I}\mathcal{O}_Y$ est un idéal quasi-cohérent de type fini,
on en conclut que $f^{-1}\mathcal{I}^n\mathcal{O}_V = 0$ pour
un ouvert non vide $V \subset Y$. Soit $Z \subset X$ le sous-espace fermé
défini par $\mathcal{I}^n$. Par construction, $V \to Y \to X$ se factorise par
$Z$. Comme $\mathbf{A}^n_Z \to \mathbf{A}^n_X$ est une immersion,
on peut remplacer $X$ par $Z$ et $Y$ par $V$.
Nous sommes ainsi ramenés au cas où $Y$ et $X$ sont irréductibles et où
$Y \to X$ envoie le point générique de $Y$ sur le point générique de $X$.

\medskip\noindent
Supposons que $Y$ et $X$ soient irréductibles, que $Y$ soit un schéma,
et que $Y \to X$ envoie le point générique de
$Y$ sur le point générique de $X$. D'après Propriétés des espaces, proposition
\ref{spaces-properties-proposition-locally-quasi-separated-open-dense-scheme}
$X$ possède un sous-schéma ouvert dense $U \subset X$. Choisissons un ouvert affine
non vide $V \subset Y$ dont l'image dans $X$ est contenue dans $U$. D'après
Morphismes, lemme \ref{morphisms-lemma-quasi-affine-finite-type-over-S},
le morphisme $V \to U$ se factorise en $V \to \mathbf{A}^n_U \to U$. En composant
avec $\mathbf{A}^n_U \to \mathbf{A}^n_X$, on obtient l'immersion voulue.
\end{proof}








\section{sections à support dans une partie fermée}
\label{section-sections-with-support-in-closed}

\noindent
Cette section est l'analogue de
Propriétés, section \ref{properties-section-sections-with-support-in-closed}.

\begin{lemma}
\label{lemma-quasi-coherent-finite-type-ideals}
Soit $S$ un schéma.
Soit $X$ un espace algébrique quasi-compact et quasi-séparé.
Soit $U \subset X$ un sous-espace ouvert. Les conditions suivantes sont équivalentes :
\begin{enumerate}
\item $U \to X$ est quasi-compact,
\item $U$ est quasi-compact, et
\item il existe un faisceau quasi-cohérent d'idéaux de type fini
$\mathcal{I} \subset \mathcal{O}_X$ tel que
$|X| \setminus |U| = |V(\mathcal{I})|$.
\end{enumerate}
\end{lemma}

\begin{proof}
Soient $W$ un schéma affine et $\varphi : W \to X$ un morphisme étale
surjectif, voir Propriétés des espaces, lemme
\ref{spaces-properties-lemma-quasi-compact-affine-cover}.
Si (1) est vérifiée, alors $\varphi^{-1}(U) \to W$ est quasi-compact, donc
$\varphi^{-1}(U)$ est quasi-compact, et par conséquent $U$ est quasi-compact
(car $|\varphi^{-1}(U)| \to |U|$ est surjectif). Si (2) est vérifiée, alors
$\varphi^{-1}(U)$ est quasi-compact, car $\varphi$ est quasi-compact
puisque $X$ est quasi-séparé (Morphismes des espaces,
lemme \ref{spaces-morphisms-lemma-quasi-compact-quasi-separated-permanence}).
Ainsi, $\varphi^{-1}(U) \to W$ est un morphisme quasi-compact de schémas d'après
Propriétés, lemme \ref{properties-lemma-quasi-coherent-finite-type-ideals}.
Il en résulte que $U \to X$ est quasi-compact d'après
Morphismes des espaces, lemme \ref{spaces-morphisms-lemma-quasi-compact-local}.
Les conditions (1) et (2) sont donc équivalentes.

\medskip\noindent
Supposons (1) et (2). D'après
Propriétés des espaces, lemme
\ref{spaces-properties-lemma-reduced-closed-subspace},
il existe un unique faisceau quasi-cohérent d'idéaux $\mathcal{J}$ définissant
la structure de sous-espace fermé réduit induite sur $|X| \setminus |U|$.
Remarquons que $\mathcal{J}|_U = \mathcal{O}_U$, qui est un
$\mathcal{O}_U$-module de type fini.
Comme $U$ est quasi-compact, il résulte du
lemme \ref{lemma-directed-colimit-finite-type}
qu'il existe un sous-faisceau quasi-cohérent
$\mathcal{I} \subset \mathcal{J}$ qui est de type fini
et tel que $\mathcal{I}|_U = \mathcal{J}|_U$.
Alors $|X| \setminus |U| = |V(\mathcal{I})|$, ce qui donne (3). Inversement,
si $\mathcal{I}$ est comme en (3), alors $\varphi^{-1}(U) \subset W$
est un ouvert quasi-compact d'après le lemme pour les schémas
(Propriétés, lemme \ref{properties-lemma-quasi-coherent-finite-type-ideals})
appliqué à $\varphi^{-1}\mathcal{I}$ sur $W$.
La condition (2) est donc vérifiée.
\end{proof}

\begin{lemma}
\label{lemma-sections-annihilated-by-ideal}
Soient $S$ un schéma et $X$ un espace algébrique sur $S$.
Soit $\mathcal{I} \subset \mathcal{O}_X$ un faisceau quasi-cohérent d'idéaux.
Soit $\mathcal{F}$ un $\mathcal{O}_X$-module quasi-cohérent.
Considérons le faisceau de $\mathcal{O}_X$-modules $\mathcal{F}'$
qui associe à tout objet $U$ de $X_\etale$ le module
$$
\mathcal{F}'(U)
=
\{s \in \mathcal{F}(U) \mid
\mathcal{I}s = 0\}
$$
Supposons que $\mathcal{I}$ soit de type fini. Alors
\begin{enumerate}
\item $\mathcal{F}'$ est un faisceau quasi-cohérent de $\mathcal{O}_X$-modules,
\item pour tout $U$ affine dans $X_\etale$, on a
$\mathcal{F}'(U) = \{s \in \mathcal{F}(U) \mid \mathcal{I}(U)s = 0\}$, et
\item $\mathcal{F}'_x = \{s \in \mathcal{F}_x \mid \mathcal{I}_x s = 0\}$.
\end{enumerate}
\end{lemma}

\begin{proof}
Il est clair que la règle définissant $\mathcal{F}'$ donne un sous-faisceau
de $\mathcal{F}$. On peut donc travailler localement pour la topologie étale sur $X$
afin de vérifier les autres assertions. Le lemme se ramène ainsi au cas des schémas,
qui est donné par
Propriétés, lemme \ref{properties-lemma-sections-annihilated-by-ideal}.
\end{proof}

\begin{definition}
\label{definition-subsheaf-sections-annihilated-by-ideal}
Soient $S$ un schéma et $X$ un espace algébrique sur $S$.
Soit $\mathcal{I} \subset \mathcal{O}_X$ un faisceau quasi-cohérent
d'idéaux de type fini.
Soit $\mathcal{F}$ un $\mathcal{O}_X$-module quasi-cohérent.
Le sous-faisceau $\mathcal{F}' \subset \mathcal{F}$ défini au
lemme \ref{lemma-sections-annihilated-by-ideal} ci-dessus est appelé
le {\it sous-faisceau des sections annulées par $\mathcal{I}$}.
\end{definition}

\begin{lemma}
\label{lemma-push-sections-annihilated-by-ideal}
Soit $S$ un schéma.
Soit $f : X \to Y$ un morphisme quasi-compact et quasi-séparé
d'espaces algébriques sur $S$.
Soit $\mathcal{I} \subset \mathcal{O}_Y$ un faisceau quasi-cohérent
d'idéaux de type fini. Soit $\mathcal{F}$ un
$\mathcal{O}_X$-module quasi-cohérent. Soit $\mathcal{F}' \subset \mathcal{F}$
le sous-faisceau des sections annulées par $f^{-1}\mathcal{I}\mathcal{O}_X$.
Alors $f_*\mathcal{F}' \subset f_*\mathcal{F}$ est le sous-faisceau
des sections annulées par $\mathcal{I}$.
\end{lemma}

\begin{proof}
Démonstration omise. Indication : l'hypothèse que $f$ est quasi-compact et
quasi-séparé implique que $f_*\mathcal{F}$ est quasi-cohérent
(Morphismes des espaces, lemme \ref{spaces-morphisms-lemma-pushforward}),
de sorte que le lemme \ref{lemma-sections-annihilated-by-ideal} s'applique
à $\mathcal{I}$ et à $f_*\mathcal{F}$.
\end{proof}

\noindent
Passons maintenant au faisceau des sections à support dans une partie fermée.
Là encore, ce n'est pas toujours un faisceau quasi-cohérent ; mais si le complémentaire
de la partie fermée est « rétrocompact » dans l'espace algébrique donné, alors
il l'est.

\begin{lemma}
\label{lemma-sections-supported-on-closed-subset}
Soient $S$ un schéma et $X$ un espace algébrique sur $S$.
Soient $T \subset |X|$ une partie fermée et $U \subset X$
le sous-espace ouvert tel que $T \amalg |U| = |X|$.
Soit $\mathcal{F}$ un $\mathcal{O}_X$-module quasi-cohérent.
Considérons le faisceau de $\mathcal{O}_X$-modules $\mathcal{F}'$
qui associe à tout objet $\varphi : W \to X$ de
$X_\etale$ le module
$$
\mathcal{F}'(W)
=
\{s \in \mathcal{F}(W) \mid
\text{le support de }s\text{ est contenu dans }|\varphi|^{-1}(T)\}
$$
Si $U \to X$ est quasi-compact, alors
\begin{enumerate}
\item pour $W$ affine, il existe un idéal de type fini
$I \subset \mathcal{O}_X(W)$ tel que $|\varphi|^{-1}(T) = V(I)$,
\item pour $W$ et $I$ comme en (1), on a
$\mathcal{F}'(W) = \{x \in \mathcal{F}(W) \mid
I^nx = 0 \text{ pour un certain } n\}$,
\item $\mathcal{F}'$ est un faisceau quasi-cohérent de $\mathcal{O}_X$-modules.
\end{enumerate}
\end{lemma}

\begin{proof}
Il est clair que la règle définissant $\mathcal{F}'$ donne un sous-faisceau
de $\mathcal{F}$. On peut donc travailler localement pour la topologie étale sur $X$
afin de vérifier les autres assertions. Le lemme se ramène ainsi au cas des schémas,
qui est donné par
Propriétés, lemme \ref{properties-lemma-sections-supported-on-closed-subset}.
\end{proof}

\begin{definition}
\label{definition-subsheaf-sections-supported-on-closed}
Soient $S$ un schéma et $X$ un espace algébrique sur $S$.
Soit $T \subset |X|$ une partie fermée dont le complémentaire
correspond à un sous-espace ouvert $U \subset X$
dont le morphisme d'inclusion $U \to X$ est quasi-compact.
Soit $\mathcal{F}$ un $\mathcal{O}_X$-module quasi-cohérent.
Le sous-faisceau quasi-cohérent $\mathcal{F}' \subset \mathcal{F}$ défini au
lemme \ref{lemma-sections-supported-on-closed-subset} ci-dessus est appelé
le {\it sous-faisceau des sections à support dans $T$}.
\end{definition}

\begin{lemma}
\label{lemma-push-sections-supported-on-closed-subset}
Soit $S$ un schéma.
Soit $f : X \to Y$ un morphisme quasi-compact et quasi-séparé
d'espaces algébriques sur $S$. Soit $T \subset |Y|$ une partie fermée.
Supposons que $|Y| \setminus T$ corresponde à un sous-espace ouvert $V \subset Y$
tel que $V \to Y$ soit quasi-compact. Soit $\mathcal{F}$ un
$\mathcal{O}_X$-module quasi-cohérent. Soit $\mathcal{F}' \subset \mathcal{F}$
le sous-faisceau des sections à support dans $|f|^{-1}T$.
Alors $f_*\mathcal{F}' \subset f_*\mathcal{F}$ est le sous-faisceau
des sections à support dans $T$.
\end{lemma}

\begin{proof}
Démonstration omise. Indications : $|X| \setminus |f|^{-1}T$ est le support du sous-espace ouvert
$U = f^{-1}V \subset X$. Comme $V \to Y$ est quasi-compact, il en est de même de
$U \to X$ (par changement de base). L'hypothèse que $f$ est quasi-compact et
quasi-séparé implique que $f_*\mathcal{F}$ est quasi-cohérent.
Le lemme \ref{lemma-sections-supported-on-closed-subset}
s'applique donc à $T$ et à $f_*\mathcal{F}$, ainsi qu'à
$|f|^{-1}T$ et à $\mathcal{F}$. L'égalité des modules quasi-cohérents considérés
résulte immédiatement des définitions.
\end{proof}















\section{Caractérisation des espaces affines}
\label{section-affine}

\noindent
Cette section est l'analogue de Limites, section \ref{limits-section-affine}.

\begin{lemma}
\label{lemma-affine}
Soient $S$ un schéma et $f : X \to Y$ un morphisme d'espaces algébriques
sur $S$. Supposons que $f$ soit surjectif et fini, et que $X$
soit affine. Alors $Y$ est affine.
\end{lemma}

\begin{proof}
Nous pouvons, et allons, considérer $f : X \to Y$ comme un morphisme d'espaces algébriques sur
$\Spec(\mathbf{Z})$ (voir
Espaces, définition \ref{spaces-definition-base-change}).
Remarquons qu'un morphisme fini est affine et universellement fermé, voir
Morphismes des espaces, lemme
\ref{spaces-morphisms-lemma-integral-universally-closed}.
D'après Morphismes des espaces, lemme
\ref{spaces-morphisms-lemma-image-universally-closed-separated}
on voit que $Y$ est un espace algébrique séparé.
Comme $f$ est surjectif et $X$ quasi-compact, on voit que $Y$ est
quasi-compact.

\medskip\noindent
D'après le lemme \ref{lemma-finite-in-finite-and-finite-presentation},
on peut écrire $X = \lim X_a$, où chaque $X_a \to Y$ est fini et de
présentation finie. D'après le
lemme \ref{lemma-limit-is-affine},
on voit que $X_a$ est affine pour $a$ assez grand.
On peut donc supposer, et l'on suppose, que $f : X \to Y$ est fini, surjectif et
de présentation finie.

\medskip\noindent
D'après la proposition \ref{proposition-approximate}, on peut écrire
$Y = \lim Y_i$ comme limite projective filtrante d'espaces
algébriques de présentation finie sur $\mathbf{Z}$.
D'après le lemme \ref{lemma-descend-finite-presentation}, on peut
trouver $0 \in I$ et un morphisme $X_0 \to Y_0$ de présentation finie
tels que $X_i = X_0 \times_{Y_0} Y_i$ pour $i \geq 0$
et que $X = \lim_i X_i$. D'après le
lemme \ref{lemma-descend-finite},
on voit que $X_i \to Y_i$ est fini pour $i$ assez grand.
D'après le lemme \ref{lemma-descend-surjective},
on voit que $X_i \to Y_i$ est surjectif pour $i$ assez grand.
D'après le lemme \ref{lemma-limit-is-affine}, on voit que $X_i$ est
affine pour $i$ assez grand. Ainsi, pour $i$ assez grand, on peut appliquer
Cohomologie des espaces, lemme
\ref{spaces-cohomology-lemma-image-affine-finite-morphism-affine-Noetherian}
pour conclure que $Y_i$ est affine. Il s'ensuit que $Y$ est affine, ce qui
achève la démonstration.
\end{proof}

\begin{proposition}
\label{proposition-affine}
Soient $S$ un schéma et $f : X \to Y$ un morphisme d'espaces algébriques
sur $S$. Supposons que $X$ soit affine et que $f$ soit surjectif et universellement
fermé\footnote{Un morphisme entier est universellement fermé, voir
Morphismes des espaces, lemme
\ref{spaces-morphisms-lemma-integral-universally-closed}.}. Alors $Y$ est affine.
\end{proposition}

\begin{proof}
Nous pouvons, et allons, considérer $f : X \to Y$ comme un morphisme d'espaces algébriques sur
$\Spec(\mathbf{Z})$ (voir
Espaces, définition \ref{spaces-definition-base-change}).
D'après Morphismes des espaces, lemme
\ref{spaces-morphisms-lemma-image-universally-closed-separated}
on voit que $Y$ est un espace algébrique séparé. Puis, d'après
Morphismes des espaces, lemme \ref{spaces-morphisms-lemma-affine-permanence},
on obtient que $f$ est affine. Dès lors, d'après
Morphismes des espaces, lemme
\ref{spaces-morphisms-lemma-integral-universally-closed}
on voit que $f$ est entier.

\medskip\noindent
D'après le paragraphe précédent, on peut supposer que $f : X \to Y$
est surjectif et entier, que $X$ est affine et que $Y$ est séparé.
Comme $f$ est surjectif et $X$ quasi-compact, on en déduit aussi que $Y$ est
quasi-compact.

\medskip\noindent
Considérons le faisceau $\mathcal{A} = f_*\mathcal{O}_X$.
C'est un faisceau quasi-cohérent de $\mathcal{O}_Y$-algèbres, voir
Morphismes des espaces, lemme \ref{spaces-morphisms-lemma-pushforward}.
D'après le lemme \ref{lemma-colimit-finitely-presented},
on peut écrire $\mathcal{A} = \colim_i \mathcal{F}_i$ comme limite inductive
filtrante de $\mathcal{O}_Y$-modules de type fini. Soit
$\mathcal{A}_i \subset \mathcal{A}$ la sous-$\mathcal{O}_Y$-algèbre
engendrée par $\mathcal{F}_i$. Comme le morphisme d'algèbres
$\mathcal{O}_Y \to \mathcal{A}$ est entier, on voit que chaque $\mathcal{A}_i$
est une $\mathcal{O}_Y$-algèbre quasi-cohérente finie. Ainsi,
$$
X_i = \underline{\Spec}_Y(\mathcal{A}_i) \longrightarrow Y
$$
est un morphisme fini d'espaces algébriques. Ici,
$\underline{\Spec}$ est la construction de Morphismes des espaces, lemme
\ref{spaces-morphisms-lemma-affine-equivalence-algebras}.
Il est clair
que $X = \lim_i X_i$. Par conséquent, d'après le
lemme \ref{lemma-limit-is-affine},
on voit que, pour $i$ assez grand, le schéma $X_i$ est affine.
De plus, puisque $X \to Y$ se factorise par chaque $X_i$, on voit que
$X_i \to Y$ est surjectif. On conclut donc que $Y$ est affine d'après le
lemme \ref{lemma-affine}.
\end{proof}

\noindent
Le corollaire suivant du résultat précédent se trouve dans
\cite{CLO}.

\begin{lemma}
\label{lemma-reduction-scheme}
\begin{reference}
\cite[3.1.12]{CLO}
\end{reference}
Soient $S$ un schéma et $X$ un espace algébrique sur $S$.
Si $X_{red}$ est un schéma, alors $X$ est un schéma.
\end{lemma}

\begin{proof}
Soit $U' \subset X_{red}$ un sous-schéma ouvert affine.
Soit $U \subset X$ le sous-espace ouvert correspondant à l'ouvert
$|U'| \subset |X_{red}| = |X|$. Alors $U' \to U$ est surjectif et
entier. Par conséquent, $U$ est affine d'après la
proposition \ref{proposition-affine}.
Ainsi, tout point est contenu dans un sous-schéma ouvert de $X$, c'est-à-dire que
$X$ est un schéma.
\end{proof}

\begin{lemma}
\label{lemma-integral-universally-bijective-scheme}
Soient $S$ un schéma et $f : X \to Y$ un morphisme d'espaces algébriques
sur $S$. Supposons que $f$ soit entier et induise une bijection $|X| \to |Y|$.
Alors $X$ est un schéma si et seulement si $Y$ est un schéma.
\end{lemma}

\begin{proof}
Un morphisme entier est représentable par définition ; si $Y$
est un schéma, il en est donc de même de $X$. Inversement, supposons que $X$ soit un schéma.
Soit $U \subset X$ un ouvert affine. Un morphisme entier est
fermé et $|f|$ est bijectif ; ainsi $|f|(|U|) \subset |Y|$
est ouvert comme complémentaire de $|f|(|X| \setminus |U|)$. Soit
$V \subset Y$ le sous-espace ouvert tel que $|V| = |f|(|U|)$, voir
Propriétés des espaces, lemme \ref{spaces-properties-lemma-open-subspaces}.
Alors $U \to V$ est entier et surjectif ; par conséquent,
$V$ est un schéma affine d'après la proposition \ref{proposition-affine}.
Cela achève la démonstration.
\end{proof}

\begin{lemma}
\label{lemma-check-closed-infinitesimally}
Soit $S$ un schéma.
Soient $f : X \to B$ et $B' \to B$ des morphismes d'espaces algébriques sur $S$.
Supposons que
\begin{enumerate}
\item $B' \to B$ soit une immersion fermée,
\item $|B'| \to |B|$ soit bijectif,
\item $X \times_B B' \to B'$ soit une immersion fermée, et
\item $X \to B$ soit de type fini ou $B' \to B$ de présentation finie.
\end{enumerate}
Alors $f : X \to B$ est une immersion fermée.
\end{lemma}

\begin{proof}
Les hypothèses (1) et (2) impliquent que $B_{red} = B'_{red}$.
Posons $X' = X \times_B B'$. Alors $X' \to X$ est une immersion fermée
et $X'_{red} = X_{red}$. Soit $U \to B$ un morphisme étale,
avec $U$ affine. Alors $X' \times_B U \to X \times_B U$ est une
immersion fermée d'espaces algébriques induisant un isomorphisme
sur les espaces réduits sous-jacents. Comme $X' \times_B U$ est un schéma
(puisque $B' \to B$ et $X' \to B'$ sont représentables), il en est de même de
$X \times_B U$ d'après le lemme \ref{lemma-reduction-scheme}.
Ainsi, $X \to B$ est également représentable. On se ramène donc au
cas des schémas, voir
Morphismes, lemme \ref{morphisms-lemma-check-closed-infinitesimally}.
\end{proof}










\section{Revêtement fini par un schéma}
\label{section-finite-cover}

\noindent
Comme application des résultats de ce chapitre sur les limites, nous démontrons que, pour
tout espace algébrique quasi-compact et quasi-séparé $X$, il existe un schéma
$Y$ et un morphisme fini surjectif $Y \to X$. Nous utiliserons le résultat déjà
démontré selon lequel on peut trouver un revêtement entier par un schéma,
établi dans
Espaces décents, section \ref{decent-spaces-section-integral-cover}.

\begin{proposition}
\label{proposition-there-is-a-scheme-finite-over}
Soient $S$ un schéma et $X$ un espace algébrique quasi-compact et quasi-séparé
sur $S$.
\begin{enumerate}
\item Il existe un morphisme fini surjectif $Y \to X$
de présentation finie, où $Y$ est un schéma,
\item étant donné un morphisme étale surjectif $U \to X$, on peut choisir
$Y \to X$ de telle sorte que, pour tout $y \in Y$, il existe un voisinage ouvert
$V \subset Y$ tel que $V \to X$ se factorise par $U$.
\end{enumerate}
\end{proposition}

\begin{proof}
La partie (1) est le cas particulier de (2) où $U = X$.
Soit $Y \to X$ comme dans
Espaces décents, lemme \ref{decent-spaces-lemma-there-is-a-scheme-integral-over}.
Choisissons un recouvrement ouvert affine fini $Y = \bigcup V_j$ tel que $V_j \to X$
se factorise par $U$. On peut écrire $Y = \lim Y_i$, où
$Y_i \to X$ est fini et de présentation finie, voir le
lemme \ref{lemma-integral-limit-finite-and-finite-presentation}.
Pour $i$ assez grand, l'espace algébrique $Y_i$ est un schéma, voir le
lemme \ref{lemma-limit-is-scheme}.
Pour $i$ assez grand, on peut trouver des ouverts affines $V_{i, j} \subset Y_i$
dont l'image réciproque dans $Y$ redonne $V_j$, voir le
lemme \ref{lemma-descend-opens}.
Pour $i$ encore plus grand, les morphismes $V_j \to U$ sur $X$ proviennent
de morphismes $V_{i, j} \to U$ sur $X$, voir la proposition
\ref{proposition-characterize-locally-finite-presentation}.
Cela achève la démonstration.
\end{proof}

\begin{lemma}
\label{lemma-integral-limit-finite-and-finite-presentation-refined}
Soient $S$ un schéma et $f : X \to Y$ un morphisme entier d'espaces
algébriques sur $S$. Supposons $Y$ quasi-compact et quasi-séparé.
Soit $V \subset Y$ un sous-espace ouvert quasi-compact tel que
$f^{-1}(V) \to V$ soit fini et de présentation finie.
Alors $X$ peut s'écrire comme limite projective filtrante $X = \lim X_i$,
où les $f_i : X_i \to Y$ sont finis et de présentation finie,
et $f^{-1}(V) \to f_i^{-1}(V)$ est un isomorphisme pour tout $i$.
\end{lemma}

\begin{proof}
Ce lemme est un léger raffinement de la proposition
\ref{proposition-there-is-a-scheme-finite-over}.
Considérons la $\mathcal{O}_Y$-algèbre quasi-cohérente entière
$\mathcal{A} = f_*\mathcal{O}_X$. Au paragraphe suivant, nous allons
écrire $\mathcal{A} = \colim \mathcal{A}_i$ comme limite inductive
filtrante de $\mathcal{O}_Y$-algèbres finies de présentation finie
$\mathcal{A}_i$ telles que $\mathcal{A}_i|_V = \mathcal{A}|_V$.
Cela fait, posons $X_i = \underline{\Spec}_Y(\mathcal{A}_i)$, voir
Morphismes des espaces, définition
\ref{spaces-morphisms-definition-relative-spec}.
Par construction, $X_i \to Y$ est fini et de présentation finie,
$X = \lim X_i$ et $f_i^{-1}(V) = f^{-1}(V)$.

\medskip\noindent
La démonstration de l'assertion sur les algèbres est analogue à
celle de la partie (2) du
lemme \ref{lemma-integral-algebra-directed-colimit-finite}.
Écrivons d'abord $\mathcal{A} = \colim \mathcal{F}_i$
comme limite inductive de $\mathcal{O}_Y$-modules de présentation finie, à l'aide du
lemme \ref{lemma-colimit-finitely-presented}.
Comme $\mathcal{A}|_V$ est un $\mathcal{O}_V$-module de type fini,
on peut supposer, et l'on suppose, que $\mathcal{F}_i|_V \to \mathcal{A}|_V$
est surjectif pour tout $i$.
Pour tout $i$, soit $\mathcal{J}_i$ le noyau du morphisme
$$
\text{Sym}^*_{\mathcal{O}_X}(\mathcal{F}_i) \longrightarrow \mathcal{A}
$$
Pour $i' \geq i$, on obtient un morphisme $\mathcal{J}_i \to \mathcal{J}_{i'}$.
On a $\mathcal{A} =
\colim \text{Sym}^*_{\mathcal{O}_X}(\mathcal{F}_i)/\mathcal{J}_i$.
De plus, les $\mathcal{O}_X$-algèbres quasi-cohérentes
$\text{Sym}^*_{\mathcal{O}_X}(\mathcal{F}_i)/\mathcal{J}_i$
sont finies (comme sous-algèbres quasi-cohérentes de type fini de la
$\mathcal{O}_Y$-algèbre quasi-cohérente entière $\mathcal{A}$ sur $\mathcal{O}_X$).
La restriction de $\text{Sym}^*_{\mathcal{O}_X}(\mathcal{F}_i)/\mathcal{J}_i$
à $V$ est $\mathcal{A}|_V$ par la surjectivité ci-dessus.
Ainsi, $\mathcal{J}_i|_V$ est de type fini comme faisceau d'idéaux de
$\text{Sym}^*_{\mathcal{O}_X}(\mathcal{F}_i)|_V$, du fait que
$\mathcal{A}|_V$ est de présentation finie comme
$\mathcal{O}_Y$-algèbre.
Écrivons $\mathcal{J}_i = \colim \mathcal{E}_{ik}$
comme limite inductive de $\mathcal{O}_X$-modules de présentation finie.
On peut supposer, et l'on suppose, que $\mathcal{E}_{ik}|_V$
engendre $\mathcal{J}_i|_V$ comme faisceau d'idéaux de
$\text{Sym}^*_{\mathcal{O}_X}(\mathcal{F}_i)|_V$, d'après
l'assertion de type fini ci-dessus.
Étant donnés $i' \geq i$ et $k$, il existe $k'$ tel que l'on
dispose d'un morphisme $\mathcal{E}_{ik} \to \mathcal{E}_{i'k'}$
rendant le diagramme
$$
\xymatrix{
\mathcal{J}_i \ar[r] & \mathcal{J}_{i'} \\
\mathcal{E}_{ik} \ar[u] \ar[r] & \mathcal{E}_{i'k'} \ar[u]
}
$$
commutatif. Cela résulte de Cohomologie des espaces, lemme
\ref{spaces-cohomology-lemma-finite-presentation-quasi-compact-colimit}.
On en déduit un morphisme
$$
\mathcal{A}_{ik} =
\text{Sym}^*_{\mathcal{O}_X}(\mathcal{F}_i)/(\mathcal{E}_{ik})
\longrightarrow
\text{Sym}^*_{\mathcal{O}_X}(\mathcal{F}_{i'})/(\mathcal{E}_{i'k'}) =
\mathcal{A}_{i'k'}
$$
où $(\mathcal{E}_{ik})$ désigne l'idéal engendré par $\mathcal{E}_{ik}$.
Les $\mathcal{O}_X$-algèbres quasi-cohérentes $\mathcal{A}_{ki}$
sont de présentation finie et finies pour $k$ assez grand
(voir la démonstration du
lemme \ref{lemma-finite-algebra-directed-colimit-finite-finitely-presented}).
De plus, on a $\mathcal{A}_{ik}|_V = \mathcal{A}|_V$ par construction.
Enfin, on a
$$
\colim \mathcal{A}_{ik} = \colim \mathcal{A}_i = \mathcal{A}
$$
En effet, la première égalité a été démontrée dans la démonstration du
lemme \ref{lemma-finite-algebra-directed-colimit-finite-finitely-presented},
et la seconde vaut parce que $\mathcal{A}$ est la limite inductive des
modules $\mathcal{F}_i$.
\end{proof}

\begin{lemma}
\label{lemma-there-is-a-scheme-finite-over-refined}
Soient $S$ un schéma et $X$ un espace algébrique quasi-compact et quasi-séparé
sur $S$, tel que $|X|$ ait un nombre fini de composantes
irréductibles.
\begin{enumerate}
\item Il existe un morphisme fini surjectif $f : Y \to X$
de présentation finie, où $Y$ est un schéma, tel que $f$
soit fini étale au-dessus d'un ouvert dense quasi-compact $U \subset X$,
\item étant donné un morphisme étale surjectif $V \to X$, on peut choisir
$Y \to X$ de telle sorte que, pour tout $y \in Y$, il existe un voisinage ouvert
$W \subset Y$ tel que $W \to X$ se factorise par $V$.
\end{enumerate}
\end{lemma}

\begin{proof}
La partie (1) est le cas particulier de (2) où $V = X$.

\medskip\noindent
Preuve de (2). Soit $\pi : Y \to X$ comme dans Espaces décents, lemme
\ref{decent-spaces-lemma-there-is-a-scheme-integral-over-refined}
et soit $U \subset X$ un ouvert dense quasi-compact tel que
$\pi^{-1}(U) \to U$ soit fini étale.
Choisissons un recouvrement ouvert affine fini $Y = \bigcup W_j$ tel que
$W_j \to X$ se factorise par $V$. On peut écrire $Y = \lim Y_i$, où
$\pi_i : Y_i \to X$ est fini et de présentation finie et où
$\pi^{-1}(U) \to \pi_i^{-1}(U)$ est un isomorphisme, voir le
lemme \ref{lemma-integral-limit-finite-and-finite-presentation-refined}.
Pour $i$ assez grand, l'espace algébrique $Y_i$ est un schéma, voir le
lemme \ref{lemma-limit-is-scheme}.
Pour $i$ assez grand, on peut trouver des ouverts affines $W_{i, j} \subset Y_i$
dont l'image réciproque dans $Y$ redonne $W_j$, voir le
lemme \ref{lemma-descend-opens}.
Pour $i$ encore plus grand, les morphismes $W_j \to V$ sur $X$ proviennent
de morphismes $W_{i, j} \to U$ sur $X$, voir la proposition
\ref{proposition-characterize-locally-finite-presentation}.
Cela achève la démonstration.
\end{proof}

\begin{lemma}
\label{lemma-there-is-a-scheme-finite-over-filtered}
Soient $S$ un schéma et $X$ un espace algébrique quasi-compact et quasi-séparé
sur $S$. Il existe un entier $t \geq 0$ et des
sous-espaces fermés
$$
X \supset Z_0 \supset Z_1 \supset \ldots \supset Z_t = \emptyset
$$
tels que $Z_i \to X$ soit de présentation finie,
que $Z_0 \subset X$ soit un épaississement et que, pour tout $i = 0, \ldots t - 1$,
il existe un schéma $Y_i$ et un morphisme fini surjectif de
présentation finie $Y_i \to Z_i$, qui soit fini étale au-dessus de
$Z_i \setminus Z_{i + 1}$.
\end{lemma}

\begin{proof}
On peut considérer $X$ comme un espace algébrique sur $\Spec(\mathbf{Z})$, voir
Espaces, définition \ref{spaces-definition-base-change}, et
Propriétés des espaces, définition \ref{spaces-properties-definition-separated}.
On peut donc appliquer la proposition \ref{proposition-approximate}.
Il en résulte que l'on peut trouver un morphisme affine $X \to X_0$,
où $X_0$ est de présentation finie sur $\mathbf{Z}$.
Si l'on peut démontrer le lemme pour $X_0$, on peut alors ramener par changement de base
la stratification et les morphismes à $X$ et obtenir le résultat pour $X$ ;
certains détails sont omis. On se ramène ainsi au cas examiné au
paragraphe suivant.

\medskip\noindent
Supposons $X$ de présentation finie sur $\mathbf{Z}$.
Alors $X$ est noethérien et $|X|$ est un espace topologique
noethérien (ayant un nombre fini de composantes irréductibles), de dimension finie.
On peut donc raisonner par récurrence sur $\dim(|X|)$.
Tout morphisme fini vers $X$ est de présentation finie ;
on peut donc négliger cette condition dans la suite de la démonstration.
D'après le lemme \ref{lemma-there-is-a-scheme-finite-over-refined},
il existe un morphisme fini surjectif $Y \to X$ qui est
fini étale au-dessus d'un ouvert dense $U \subset X$.
Posons $Z_0 = X$ et soit $Z_1 \subset X$ le sous-espace fermé réduit
tel que $|Z_1| = |X| \setminus |U|$.
Par récurrence, on trouve un entier $t \geq 0$ et une filtration
$$
Z_1 \supset Z_{1, 0} \supset Z_{1, 1} \supset \ldots
\supset Z_{1, t} = \emptyset
$$
par sous-espaces fermés, où $Z_{1, 0} \to Z_1$ est un épaississement
et où il existe des morphismes finis surjectifs $Y_{1, i} \to Z_{1, i}$
qui sont finis étales au-dessus de $Z_{1, i} \setminus Z_{1, i + 1}$.
Comme $Z_1$ est réduit, on a $Z_1 = Z_{1, 0}$.
On peut donc poser $Z_i = Z_{1, i - 1}$ et $Y_i = Y_{1, i - 1}$
pour $i \geq 1$, ce qui démontre le lemme.
\end{proof}











\section{Obtention de schémas}
\label{section-representable}

\noindent
Voici quelques techniques supplémentaires pour montrer qu'un espace algébrique est un schéma.
La première consiste à montrer qu'il existe un sous-espace fermé minimal
qui n'est pas un schéma.

\begin{lemma}
\label{lemma-minimal-closed-subspace}
Soient $S$ un schéma et $X$ un espace algébrique quasi-compact et quasi-séparé
sur $S$. Si $X$ n'est pas un schéma, alors il existe
un sous-espace fermé $Z \subset X$ tel que $Z$ ne soit pas un schéma, mais que
tout sous-espace fermé propre $Z' \subset Z$ soit un schéma.
\end{lemma}

\begin{proof}
Nous démontrons cela par le lemme de Zorn. Soit $\mathcal{Z}$ l'ensemble
des sous-espaces fermés $Z$ qui ne sont pas des schémas, ordonné par inclusion.
Par hypothèse, $\mathcal{Z}$ contient $X$ et n'est donc pas vide.
Si les $Z_\alpha$ forment une partie totalement ordonnée de $\mathcal{Z}$, alors
$Z = \bigcap Z_\alpha$ appartient à $\mathcal{Z}$. En effet,
$$
Z = \lim Z_\alpha
$$
et les morphismes de transition sont affines.
On peut donc appliquer le lemme \ref{lemma-limit-is-scheme} : si $Z$
était un schéma, l'un des $Z_\alpha$ le serait aussi.
(Cela fonctionne même si $Z = \emptyset$, mais remarquons que, d'après le
lemme \ref{lemma-limit-nonempty}, ce cas ne peut pas se produire.)
Ainsi, $\mathcal{Z}$ possède des éléments minimaux d'après le lemme de Zorn.
\end{proof}

\noindent
Nous pouvons maintenant établir quelques propriétés de ces non-schémas minimaux.

\begin{lemma}
\label{lemma-minimal-nonscheme}
Soient $S$ un schéma et $X$ un espace algébrique quasi-compact et quasi-séparé
sur $S$. Supposons que tout sous-espace fermé propre
$Z \subset X$ soit un schéma, mais que $X$ ne soit pas un schéma. Alors $X$ est réduit
et irréductible.
\end{lemma}

\begin{proof}
On voit que $X$ est réduit d'après le lemme \ref{lemma-reduction-scheme}.
Choisissons des parties fermées $T_1 \subset |X|$ et $T_2 \subset |X|$ telles que
$|X| = T_1 \cup T_2$. Si $T_1$ et $T_2$ sont des parties fermées propres,
alors les sous-espaces fermés réduits induits correspondants $Z_1, Z_2 \subset X$
(Propriétés des espaces, définition
\ref{spaces-properties-definition-reduced-induced-space})
sont des schémas, de même que $Z = Z_1 \times_X Z_2 = Z_1 \cap Z_2$, comme sous-schéma
fermé de $Z_1$ ou de $Z_2$. Remarquons que la somme amalgamée
$Z_1 \amalg_Z Z_2$ existe dans la catégorie des schémas, voir
Compléments sur les morphismes, lemme
\ref{more-morphisms-lemma-pushout-along-closed-immersions}.
Une manière de procéder consiste à montrer que $Z_1 \amalg_Z Z_2$ est isomorphe à $X$,
mais nous ne pouvons pas l'utiliser ici, car les sommes amalgamées d'espaces
algébriques ne sont étudiées que plus loin. Nous utiliserons plutôt le
lemme \ref{lemma-affine} pour trouver un voisinage affine de tout point.
Soit en effet $x \in |X|$. Si $x \not \in Z_1$, alors $x$ possède un voisinage
qui est un schéma, à savoir $X \setminus Z_1$. De même si $x \not \in Z_2$.
Si $x \in Z = Z_1 \cap Z_2$, choisissons un ouvert affine
$U \subset Z_1 \amalg_Z Z_2$ contenant $z$. Alors $U_1 = Z_1 \cap U$
et $U_2 = Z_2 \cap U$ sont des ouverts affines dont les intersections avec
$Z$ coïncident. Comme $|Z_1| = T_1$ et $|Z_2| = T_2$ sont des parties fermées de
$|X|$ dont l'intersection est $|Z|$, on trouve un ouvert $W \subset |X|$
tel que $W \cap T_1 = |U_1|$ et $W \cap T_2 = |U_2|$. Notons encore $W$ le
sous-espace ouvert correspondant de $X$. Alors $x \in |W|$ et le morphisme
$U_1 \amalg U_2 \to W$ est un morphisme fini surjectif dont la source
est un schéma affine. Ainsi, $W$ est un schéma affine d'après le
lemme \ref{lemma-affine}.
\end{proof}

\noindent
Un point essentiel du lemme suivant est qu'il suffit de vérifier
la condition aux images des points de $X$.

\begin{lemma}
\label{lemma-enough-local}
Soit $f: X \to S$ un morphisme quasi-compact et quasi-séparé d'un
espace algébrique vers un schéma $S$. Si, pour tout $x \in |X|$ d'image
$s = f(x) \in S$, l'espace algébrique $X \times_S \Spec(\mathcal{O}_{S,s})$
est un schéma, alors $X$ est un schéma.
\end{lemma}

\begin{proof}
Soit $x \in |X|$. Il suffit de trouver un voisinage ouvert $U$ de
$s = f(x)$ tel que $X \times_S U$ soit un schéma.
Comme $X \times_S \Spec(\mathcal{O}_{S, s})$ est un schéma et que
$\mathcal{O}_{S, s} = \colim \mathcal{O}_S(U)$, où la limite inductive porte
sur les voisinages ouverts affines de $s$ dans $S$, on voit que
$$
X \times_S \Spec(\mathcal{O}_{S, s}) = \lim X \times_S U
$$
D'après le lemme \ref{lemma-limit-is-scheme}, on voit que $X \times_S U$
est un schéma pour un certain $U$.
\end{proof}

\noindent
Au lieu de se restreindre aux anneaux locaux comme dans le lemme \ref{lemma-enough-local},
on peut se restreindre aux sous-schémas fermés de la base.

\begin{lemma}
\label{lemma-maximal-ideal}
Soit $\varphi : X \to \Spec(A)$ un morphisme quasi-compact et quasi-séparé
d'un espace algébrique vers un schéma affine.
Si $X$ n'est pas un schéma, alors il existe un idéal $I \subset A$
tel que le changement de base $X_{A/I}$ ne soit pas un schéma, mais que,
pour tout $I \subset I'$ avec $I \not = I'$, le changement de base
$X_{A/I'}$ soit un schéma.
\end{lemma}

\begin{proof}
Nous démontrons cela par le lemme de Zorn. Soit $\mathcal{I}$ l'ensemble
des idéaux $I$ tels que $X_{A/I}$ ne soit pas un schéma. Par
hypothèse, $\mathcal{I}$ contient $(0)$. Si les $I_\alpha$ forment
une chaîne d'idéaux de $\mathcal{I}$, alors
$I = \bigcup I_\alpha$ appartient à $\mathcal{I}$. En effet,
$A/I = \colim A/I_\alpha$, d'où
$$
X_{A/I} = \lim X_{A/I_\alpha}
$$
On peut donc appliquer le lemme \ref{lemma-limit-is-scheme} : si $X_{A/I}$
était un schéma, l'un des $X_{A/I_\alpha}$ le serait aussi.
Ainsi, $\mathcal{I}$ possède des éléments maximaux d'après le lemme de Zorn.
\end{proof}






\section{Recollement dans les fibres fermées}
\label{section-change-over-closed-points}

\noindent
En appliquant la théorie précédente au spectre d'un anneau local, on obtient
quelques agréables résultats de recollement pour les espaces algébriques relatifs.
Nous démontrons d'abord un lemme auxiliaire (qui sera considérablement généralisé
dans Amorçage, section \ref{bootstrap-section-applications}).

\begin{lemma}
\label{lemma-relative-glueing}
Soit $S = U \cup W$ un recouvrement ouvert d'un schéma. Alors le foncteur
$$
FP_S \longrightarrow FP_U \times_{FP_{U \cap W}} FP_W
$$
donné par changement de base est une équivalence, où $FP_T$
est la catégorie des espaces algébriques de présentation finie sur
le schéma $T$.
\end{lemma}

\begin{proof}
Tout d'abord, puisque $S = U \cup W$ est un recouvrement de Zariski, on voit que la
catégorie des faisceaux sur $(\Sch/S)_{fppf}$ est équivalente à la catégorie
des triplets $(\mathcal{F}_U, \mathcal{F}_W, \varphi)$, où
$\mathcal{F}_U$ est un faisceau sur $(\Sch/U)_{fppf}$,
$\mathcal{F}_W$ est un faisceau sur $(\Sch/W)_{fppf}$, et où
$$
\varphi :
\mathcal{F}_U|_{(\Sch/U \cap W)_{fppf}}
\longrightarrow
\mathcal{F}_W|_{(\Sch/U \cap W)_{fppf}}
$$
est un isomorphisme. Voir Sites, lemme \ref{sites-lemma-mapping-property-glue}
(remarquons qu'aucune autre donnée de recollement n'est nécessaire, car
$U \times_S U = U$, $W \times_S W = W$, et que la condition de cocycle
est automatique pour la même raison).
Maintenant, si le faisceau $\mathcal{F}$ sur $(\Sch/S)_{fppf}$
correspond à $(\mathcal{F}_U, \mathcal{F}_W, \varphi)$
par cette équivalence, alors $\mathcal{F}$ est un espace algébrique
si et seulement si $\mathcal{F}_U$ et $\mathcal{F}_W$ sont des espaces algébriques.
Cela résulte immédiatement de
Espaces algébriques, lemme \ref{spaces-lemma-glueing-algebraic-spaces},
car $\mathcal{F}_U \to \mathcal{F}$ et $\mathcal{F}_W \to \mathcal{F}$
sont représentables par des immersions ouvertes et recouvrent $\mathcal{F}$.
Enfin, dans ce cas, l'espace algébrique $\mathcal{F}$ est de
présentation finie sur $S$ si et seulement si $\mathcal{F}_U$ est de présentation finie
sur $U$ et $\mathcal{F}_W$ de présentation finie sur $W$,
d'après Morphismes des espaces, lemmes
\ref{spaces-morphisms-lemma-quasi-compact-local},
\ref{spaces-morphisms-lemma-separated-local}, et
\ref{spaces-morphisms-lemma-finite-presentation-local}.
\end{proof}

\begin{lemma}
\label{lemma-glueing-near-closed-point}
Soient $S$ un schéma et $s \in S$ un point fermé tel que
$U = S \setminus \{s\} \to S$ soit quasi-compact. En posant
$V = \Spec(\mathcal{O}_{S, s}) \setminus \{s\}$, on a
une équivalence de catégories
$$
FP_S \longrightarrow FP_U \times_{FP_V} FP_{\Spec(\mathcal{O}_{S, s})}
$$
où $FP_T$ est la catégorie des espaces algébriques de présentation finie
sur $T$.
\end{lemma}

\begin{proof}
Soit $W \subset S$ un voisinage ouvert de $s$. Le foncteur
$$
FP_S \to FP_U \times_{FP_{W \setminus \{s\}}} FP_W
$$
est une équivalence de catégories d'après le lemme \ref{lemma-relative-glueing}.
On a $\mathcal{O}_{S, s} = \colim \mathcal{O}_W(W)$, où
$W$ parcourt les voisinages ouverts affines de $s$.
Ainsi, $\Spec(\mathcal{O}_{S, s}) = \lim W$, où $W$
parcourt les voisinages ouverts affines de $s$.
La catégorie des espaces algébriques de présentation finie
sur $\Spec(\mathcal{O}_{S, s})$ est donc la limite des
catégories d'espaces algébriques de présentation finie sur
$W$, où $W$ parcourt les voisinages ouverts affines
de $s$, voir le
lemme \ref{lemma-descend-finite-presentation}.
Pour tout ouvert affine $s \in W$, on voit que $U \cap W$
est quasi-compact puisque $U \to S$ est quasi-compact.
Ainsi, $V = \lim W \cap U = \lim W \setminus \{s\}$ est une limite de
schémas quasi-compacts et quasi-séparés (voir
Limites, lemme \ref{limits-lemma-directed-inverse-system-has-limit}).
La catégorie des espaces algébriques de présentation finie
sur $V$ est donc elle aussi la limite des
catégories d'espaces algébriques de présentation finie sur
$W \cap U$, où $W$ parcourt les voisinages ouverts affines
de $s$. Le lemme résulte formellement de la combinaison
de ces résultats.
\end{proof}

\begin{lemma}
\label{lemma-glueing-near-point}
Soient $S$ un schéma et $U \subset S$ un ouvert rétrocompact.
Soit $s \in S$ un point du complémentaire de $U$. En posant
$V = \Spec(\mathcal{O}_{S, s}) \cap U$, on a
une équivalence de catégories
$$
\colim_{s \in U' \supset U\text{ ouvert}} FP_{U'}
\longrightarrow
FP_U \times_{FP_V} FP_{\Spec(\mathcal{O}_{S, s})}
$$
où $FP_T$ est la catégorie des espaces algébriques de présentation finie
sur $T$.
\end{lemma}

\begin{proof}
Soit $W \subset S$ un voisinage ouvert de $s$. D'après le
lemme \ref{lemma-relative-glueing}, le foncteur
$$
FP_{U \cup W}
\longrightarrow
FP_U \times_{FP_{U \cap W}} FP_W
$$
est une équivalence de catégories. On a
$\mathcal{O}_{S, s} = \colim \mathcal{O}_W(W)$, où
$W$ parcourt les voisinages ouverts affines de $s$.
Ainsi, $\Spec(\mathcal{O}_{S, s}) = \lim W$, où $W$
parcourt les voisinages ouverts affines de $s$.
La catégorie des espaces algébriques de présentation finie
sur $\Spec(\mathcal{O}_{S, s})$ est donc la limite des
catégories d'espaces algébriques de présentation finie sur
$W$, où $W$ parcourt les voisinages ouverts affines
de $s$, voir le
lemme \ref{lemma-descend-finite-presentation}.
Pour tout ouvert affine $s \in W$, on voit que $U \cap W$
est quasi-compact puisque $U \to S$ est quasi-compact.
Ainsi, $V = \lim W \cap U$ est une limite de
schémas quasi-compacts et quasi-séparés (voir
Limites, lemme \ref{limits-lemma-directed-inverse-system-has-limit}).
La catégorie des espaces algébriques de présentation finie
sur $V$ est donc elle aussi la limite des
catégories d'espaces algébriques de présentation finie sur
$W \cap U$, où $W$ parcourt les voisinages ouverts affines
de $s$. Le lemme résulte formellement de la combinaison
de ces résultats.
\end{proof}

\begin{lemma}
\label{lemma-glueing-near-multiple-closed-points}
Soient $S$ un schéma et $s_1, \ldots, s_n \in S$ des
points fermés deux à deux distincts tels que
$U = S \setminus \{s_1, \ldots, s_n\} \to S$ soit quasi-compact. En posant
$S_i = \Spec(\mathcal{O}_{S, s_i})$ et $U_i = S_i \setminus \{s_i\}$,
on a une équivalence de catégories
$$
FP_S \longrightarrow
FP_U \times_{(FP_{U_1} \times \ldots \times FP_{U_n})}
(FP_{S_1} \times \ldots \times FP_{S_n})
$$
où $FP_T$ est la catégorie des espaces algébriques de présentation finie
sur $T$.
\end{lemma}

\begin{proof}
Pour $n = 1$, c'est le lemme \ref{lemma-glueing-near-closed-point}.
Pour $n > 1$, le lemme se démontre exactement de la même façon ou
peut s'en déduire. Par exemple, supposons que $f_i : X_i \to S_i$
soient des objets de $FP_{S_i}$ et que $f : X \to U$ soit un objet
de $FP_U$, et qu'on se soit donné des isomorphismes $X_i \times_{S_i} U_i = X \times_U U_i$.
D'après le lemme \ref{lemma-glueing-near-closed-point}, on peut trouver
un morphisme $f' : X' \to U' = S \setminus \{s_1, \ldots, s_{n - 1}\}$
qui est de présentation finie, qui est isomorphe à
$X_i$ au-dessus de $S_i$, qui est isomorphe à $X$ au-dessus de $U$, et
ces isomorphismes sont compatibles avec l'isomorphisme donné
$X_i \times_{S_n} U_n = X \times_U U_n$.
On peut alors appliquer l'hypothèse de récurrence à
$f_i : X_i \to S_i$, $i \leq n - 1$,
$f' : X' \to U'$, et aux
isomorphismes induits $X_i \times_{S_i} U_i = X' \times_{U'} U_i$, $i \leq n - 1$.
Cela démontre la surjectivité essentielle. Nous omettons la preuve de la
pleine fidélité.
\end{proof}







\section{Application aux modifications}
\label{section-modifications-at-a-point}

\noindent
À l'aide des limites, on peut décrire la catégorie des modifications d'un
espace algébrique décent en un point fermé au moyen de
l'anneau local hensélien.

\begin{lemma}
\label{lemma-excision-modifications}
Soit $S$ un schéma. Considérons un morphisme séparé et \'etale
$f : V \to W$ d'espaces algébriques sur $S$.
Supposons qu'il existe un
sous-espace fermé $T \subset W$ tel que $f^{-1}T \to T$ soit
un isomorphisme. Alors, en posant $W^0 = W \setminus T$ et
$V^0 = f^{-1}W^0$, le foncteur de changement de base
$$
\left\{
\begin{matrix}
g : X \to W\text{ morphisme d'espaces algébriques} \\
g^{-1}(W^0) \to W^0\text{ est un isomorphisme}
\end{matrix}
\right\}
\longrightarrow
\left\{
\begin{matrix}
h : Y \to V\text{ morphisme d'espaces algébriques} \\
h^{-1}(V^0) \to V^0\text{ est un isomorphisme}
\end{matrix}
\right\}
$$
est une équivalence de catégories.
\end{lemma}

\begin{proof}
Comme $V \to W$ est séparé, on a
$V \times_W V = \Delta(V) \amalg U$ pour un certain sous-espace ouvert et fermé
$U$ de $V \times_W V$. L'hypothèse que $f^{-1}T \to T$ est un
isomorphisme donne $U \times_W T = \emptyset$, c'est-à-dire que les deux
projections $U \to V$ prennent leurs valeurs dans $V^0$.

\medskip\noindent
Étant donné $h : Y \to V$ dans la catégorie de droite, considérons le
foncteur contravariant $X$ sur $(\Sch/S)_{fppf}$ défini par la règle
$$
X(T) = \{(w, y) \mid
w :  T \to W,\ y : T \times_{w, W} V \to Y\text{ morphisme sur }V\}
$$
Notons $g : X \to W$ l'application qui envoie $(w, y) \in X(T)$ sur $w \in W(T)$.
Comme $h^{-1}V^0 \to V^0$ est un isomorphisme, si
$w : T \to W$ prend ses valeurs dans $W^0$, il n'existe qu'un seul choix pour $h$.
Autrement dit, $X \times_{g, W} W^0 = W^0$. D'autre part, considérons
un point à valeurs dans $T$, $(w, y, v)$, de $X \times_{g, W, f} V$.
Alors $w = f \circ v$ et
$$
y : T \times_{f \circ v, W} V \longrightarrow V
$$
est un morphisme sur $V$. Considérons le morphisme
$$
T \times_{f \circ v, W} V \xrightarrow{(v, \text{id}_V)}
V \times_W V = V \amalg U
$$
L'image réciproque de $V$ est $T$, plongé par
$(\text{id}_T, v) : T \to T \times_{f \circ v, W} V$.
Le composé $y' = y \circ (\text{id}_T, v) : T \to Y$
est un morphisme tel que $v = h \circ y'$; il détermine $y$, car la
restriction de $y$ à l'autre composante est uniquement déterminée puisque
$U$ prend ses valeurs dans $V^0$ par la seconde projection. Il s'ensuit que
$X \times_{g, W, f} V \to Y$, $(w, y, v) \mapsto y'$ est un isomorphisme.

\medskip\noindent
Ainsi, il suffit de montrer que $X$ est un espace algébrique.
Puisque $V \to W$ est séparé et \'etale, il est représentable d'après
Morphismes d'espaces, lemme
\ref{spaces-morphisms-lemma-locally-quasi-finite-separated-representable}
(et Morphismes d'espaces, lemme
\ref{spaces-morphisms-lemma-etale-locally-quasi-finite}).
Il est clair que $W^0 \to W$ est représentable et \'etale, puisqu'il s'agit d'une
immersion ouverte. Ainsi
$$
W^0 \amalg Y = X \times_{g, W} W^0 \amalg X \times_{g, W, f} V
= X \times_{g, W} (W^0 \amalg V) \longrightarrow X
$$
est représentable, surjectif et \'etale d'après Espaces, lemmes
\ref{spaces-lemma-base-change-representable-transformations} et
\ref{spaces-lemma-base-change-representable-transformations-property}.
Ainsi, $X$ est un espace
algébrique d'après Espaces, lemme
\ref{spaces-lemma-etale-locally-representable-by-space-gives-space}.
\end{proof}

\begin{lemma}
\label{lemma-excision-modifications-properties}
Notations et hypothèses comme dans le lemme \ref{lemma-excision-modifications}.
Soient $g : X \to W$ et $h : Y \to V$ qui se correspondent par l'équivalence.
Alors $g$ est quasi-compact, quasi-séparé, séparé, localement de présentation
finie, de présentation finie, localement de type fini, de type fini,
propre, entier, fini, et ajouter d'autres propriétés ici si et seulement si
$h$ possède la même propriété.
\end{lemma}

\begin{proof}
Si $g$ est quasi-compact, quasi-séparé, séparé, localement de présentation
finie, de présentation finie, localement de type fini, de type fini,
propre, fini, il en est de même de $h$, comme changement de base de $g$, d'après
Morphismes d'espaces, lemmes
\ref{spaces-morphisms-lemma-base-change-quasi-compact},
\ref{spaces-morphisms-lemma-base-change-separated},
\ref{spaces-morphisms-lemma-base-change-finite-presentation},
\ref{spaces-morphisms-lemma-base-change-finite-type},
\ref{spaces-morphisms-lemma-base-change-proper},
\ref{spaces-morphisms-lemma-base-change-integral}.
Réciproquement, soit $P$ une propriété des morphismes d'espaces
algébriques qui est locale pour la topologie \'etale sur la base et qui est vérifiée par
le morphisme identité de tout espace algébrique.
Puisque $\{W^0 \to W, V \to W\}$ est un recouvrement \'etale,
pour prouver que $g$ possède $P$, il suffit de montrer
que $h$ possède $P$. On conclut donc à l'aide de
Morphismes d'espaces, lemmes
\ref{spaces-morphisms-lemma-quasi-compact-local},
\ref{spaces-morphisms-lemma-separated-local},
\ref{spaces-morphisms-lemma-finite-presentation-local},
\ref{spaces-morphisms-lemma-finite-type-local},
\ref{spaces-morphisms-lemma-proper-local},
\ref{spaces-morphisms-lemma-integral-local}.
\end{proof}

\begin{lemma}
\label{lemma-modifications}
Soit $S$ un schéma. Soit $X$ un espace algébrique décent sur $S$.
Soit $x \in |X|$ un point fermé tel que $U = X \setminus \{x\} \to X$
soit quasi-compact. En posant
$V = \Spec(\mathcal{O}_{X, x}^h) \setminus \{\mathfrak m_x^h\}$,
le foncteur de changement de base
$$
\left\{
\begin{matrix}
f : Y \to X\text{ de présentation finie} \\
f^{-1}(U) \to U\text{ est un isomorphisme}
\end{matrix}
\right\}
\longrightarrow
\left\{
\begin{matrix}
g : Y \to \Spec(\mathcal{O}_{X, x}^h)\text{ de présentation finie} \\
g^{-1}(V) \to V\text{ est un isomorphisme}
\end{matrix}
\right\}
$$
est une équivalence de catégories.
\end{lemma}

\begin{proof}
Soit $a : (W, w) \to (X, x)$ un voisinage \'etale élémentaire de $x$
avec $W$ affine, comme dans
Espaces décents, lemme
\ref{decent-spaces-lemma-decent-space-elementary-etale-neighbourhood}.
Comme $x$ est un point fermé de $X$ et que $w$ est l'unique point de $W$
au-dessus de $x$, on voit que $w$ est un point fermé de $W$. Puisque $a$
est \'etale et identifie les corps résiduels en $x$ et $w$, il
s'ensuit que $a$ induit un isomorphisme $a^{-1}x \to x$ (comme sous-espaces
fermés de $X$ et $W$). On peut donc appliquer
les lemmes \ref{lemma-excision-modifications} et
\ref{lemma-excision-modifications-properties}
pour se ramener au cas où $X$ est un schéma affine.

\medskip\noindent
Supposons $X$ affine. Rappelons que $\mathcal{O}_{X, x}^h$
est la limite inductive des $\Gamma(U, \mathcal{O}_U)$ sur les
voisinages \'etales élémentaires affines $(U, u) \to (X, x)$.
Rappelons que la catégorie de ces voisinages est
cofiltrante; voir Espaces décents, lemme
\ref{decent-spaces-lemma-elementary-etale-neighbourhoods} ou
Compléments sur les morphismes, lemme
\ref{more-morphisms-lemma-elementary-etale-neighbourhoods}.
Alors $\Spec(\mathcal{O}_{X, x}^h) = \lim U$ et
$V = \lim U \setminus \{u\}$
(lemme \ref{lemma-directed-inverse-system-has-limit}),
où les limites sont prises sur la même catégorie. Ainsi, d'après
le lemme \ref{lemma-descend-finite-presentation},
la catégorie de droite est la limite inductive des catégories
associées aux couples $(U, u)$. D'après le premier
paragraphe, chacune de ces catégories est équivalente à la
catégorie associée au couple $(X, x)$. Cela achève la démonstration.
\end{proof}






\section{Morphismes universellement fermés}
\label{section-universally-closed}

\noindent
Dans cette section, nous étudions à quelle condition un morphisme quasi-compact
(mais non nécessairement séparé) est universellement fermé. Nous démontrons d'abord un lemme qui
permettra de vérifier cette propriété après un changement de base
localement de présentation finie.

\begin{lemma}
\label{lemma-separate}
Soit $S$ un schéma. Soient $f : X \to Y$ et $g : Z \to Y$ des
morphismes d'espaces algébriques sur $S$. Soit $z \in |Z|$ et soit
$T \subset |X \times_Y Z|$ une partie fermée
telle que $z \not \in \Im(T \to |Z|)$.
Si $f$ est quasi-compact, il existe
un voisinage \'etale $(V, v) \to (Z, z)$,
un diagramme commutatif
$$
\xymatrix{
V \ar[d] \ar[r]_a & Z' \ar[d]^b \\
Z \ar[r]^g & Y,
}
$$
et une partie fermée $T' \subset |X \times_Y Z'|$ tels que
\begin{enumerate}
\item le morphisme $b : Z' \to Y$ soit localement de présentation finie,
\item en posant $z' = a(v)$, on ait $z' \not \in \Im(T' \to |Z'|)$, et
\item l'image réciproque de $T$ dans $|X \times_Y V|$
ait son image dans $T'$ par $|X \times_Y V| \to |X \times_Y Z'|$.
\end{enumerate}
De plus, on peut supposer que $V$ et $Z'$ sont des schémas affines et, si $Z$
est un schéma, que $V$ est un voisinage ouvert affine de $z$.
\end{lemma}

\begin{proof}
Nous allons déduire ce résultat du résultat analogue pour les morphismes de schémas.
Soit $y \in |Y|$ l'image de $z$. Choisissons d'abord un voisinage \'etale affine
$(U, u) \to (Y, y)$, puis un voisinage \'etale affine
$(V, v) \to (Z, z)$ tel que le morphisme $V \to Y$
se factorise par $U$. On peut alors remplacer
\begin{enumerate}
\item $X \to Y$ par $X \times_Y U \to U$,
\item $Z \to Y$ par $V \to U$,
\item $z$ par $v$, et
\item $T$ par son image réciproque dans
$|(X \times_Y U) \times_U V| = |X \times_Y V|$.
\end{enumerate}
En effet, nous montrerons ci-dessous qu'après avoir remplacé $V$ par un voisinage
ouvert affine de $v$, il existe un morphisme $a : V \to Z'$ vers
un certain $Z' \to U$ de présentation finie et une partie fermée $T'$
de $|(X \times_Y U) \times_U Z'| = |X \times_Y Z'|$ tels que
$T$ ait son image dans $T'$ et que $a(v) \not \in \Im(T' \to |Z'|)$.
On peut donc supposer, et l'on suppose désormais, que $Z$ et $Y$ sont des schémas affines,
étant entendu qu'il faut trouver une solution où $V$
est un voisinage ouvert de $z$.

\medskip\noindent
Puisque $f$ est quasi-compact et $Y$ affine, l'espace algébrique
$X$ est quasi-compact. Choisissons un schéma affine $W$ et un morphisme
\'etale surjectif $W \to X$. Soit $T_W \subset |W \times_Y Z|$
l'image réciproque de $T$. Alors $z$ n'appartient pas à l'image de
$T_W$. D'après le cas des schémas (Limites, lemme \ref{limits-lemma-separate}),
on peut trouver un voisinage ouvert $V \subset Z$ de $z$,
un diagramme commutatif de schémas
$$
\xymatrix{
V \ar[d] \ar[r]_a & Z' \ar[d]^b \\
Z \ar[r]^g & Y,
}
$$
et une partie fermée $T' \subset |W \times_Y Z'|$ tels que
\begin{enumerate}
\item le morphisme $b : Z' \to Y$ soit localement de présentation finie,
\item en posant $z' = a(z)$, on ait $z' \not \in \Im(T' \to Z')$, et
\item $T_1 = T_W \cap |W \times_Y V|$ ait son image dans $T'$ par
$|W \times_Y V| \to |W \times_Y Z'|$.
\end{enumerate}
Le diagramme commutatif
$$
\xymatrix{
W \times_Y Z \ar[d] &
W \times_Y V \ar[l] \ar[rr]_{a_1} \ar[d]_c & &
W \times_Y Z' \ar[d]^q \\
X \times_Y Z &
X \times_Y V \ar[l] \ar[rr]^{a_2} & &
X \times_Y Z'
}
$$
a des carrés cartésiens, et les morphismes verticaux sont surjectifs, \'etales
et, a fortiori, ouverts. En considérant le carré de gauche, on
voit que $T_1 = T_W \cap |W \times_Y V|$ est l'image réciproque de
$T_2 = T \cap |X \times_Y V|$ par $c$. D'après Propriétés des espaces, lemme
\ref{spaces-properties-lemma-points-cartesian}, on a
$a_1(T_1) = q^{-1}(a_2(T_2))$.
D'après Topologie, lemme \ref{topology-lemma-open-morphism-quotient-topology},
on a
$$
q^{-1}\left(\overline{a_2(T_2)}\right) =
\overline{q^{-1}(a_2(T_2))} =
\overline{a_1(T_1)} \subset T'
$$
Comme $q$ est surjectif, l'image de $\overline{a_2(T_2)} \to |Z'|$
ne contient pas $z'$, puisque cela est vrai pour $T'$.
On peut donc prendre le diagramme ci-dessus avec $Z', V, a, b$ et la
partie fermée $\overline{a_2(T_2)} \subset |X \times_Y Z'|$ comme
solution du problème posé par le lemme.
\end{proof}

\begin{lemma}
\label{lemma-test-universally-closed}
Soit $S$ un schéma.
Soit $f : X \to Y$ un morphisme quasi-compact d'espaces algébriques sur $S$.
Les conditions suivantes sont équivalentes :
\begin{enumerate}
\item $f$ est universellement fermé,
\item pour tout morphisme $Z \to Y$ localement de présentation finie,
l'application $|X \times_Y Z| \to |Z|$ est fermée, et
\item il existe un schéma $V$ et un morphisme \'etale surjectif $V \to Y$
de sorte que l'application $|\mathbf{A}^n \times (X \times_Y V)| \to |\mathbf{A}^n \times V|$
soit fermée pour tout $n \geq 0$.
\end{enumerate}
\end{lemma}

\begin{proof}
Il est clair que (1) implique (2).
Supposons que $|X \times_Y Z| \to |Z|$ ne soit pas fermée pour un certain
morphisme d'espaces algébriques $Z \to Y$ sur $S$. Cela signifie qu'il
existe une partie fermée $T \subset |X \times_Y Z|$
telle que $\Im(T \to |Z|)$ ne soit pas fermée. Choisissons $z \in |Z|$
dans l'adhérence de l'image de $T$, mais non dans cette image.
Appliquons le lemme \ref{lemma-separate}.
On obtient un voisinage \'etale $(V, v) \to (Z, z)$, un diagramme commutatif
$$
\xymatrix{
V \ar[d] \ar[r]_a & Z' \ar[d]^b \\
Z \ar[r]^g & Y,
}
$$
et une partie fermée $T' \subset |X \times_Y Z'|$ tels que
\begin{enumerate}
\item le morphisme $b : Z' \to Y$ soit localement de présentation finie,
\item en posant $z' = a(v)$, on ait $z' \not \in \Im(T' \to |Z'|)$, et
\item l'image réciproque de $T$ dans $|X \times_Y V|$ ait son image dans $T'$ par
$|X \times_Y V| \to |X \times_Y Z'|$.
\end{enumerate}
Nous affirmons que $z'$ appartient à l'adhérence de $\Im(T' \to |Z'|)$,
ce qui implique que $|X \times_Y Z'| \to |Z'|$ n'est pas fermée.
Cette affirmation montre que (2) implique (1).
Pour la vérifier, considérons
le diagramme commutatif suivant :
$$
\xymatrix{
X \times_Y Z \ar[d] &
X \times_Y V \ar[l] \ar[d] \ar[r] &
X \times_Y Z' \ar[d] \\
Z & V \ar[l] \ar[r]^a & Z'
}
$$
Soit $T_V \subset |X \times_Y V|$ l'image réciproque de $T$.
D'après Propriétés des espaces, lemme \ref{spaces-properties-lemma-points-cartesian},
l'image de $T_V$ dans $|V|$ est l'image réciproque de l'image
de $T$ dans $|Z|$. Puisque $z$ appartient à l'adhérence de l'image de
$T \to |Z|$ et que $|V| \to |Z|$ est ouverte, on voit que $v$ appartient à
l'adhérence de l'image de $T_V \to |V|$. Comme l'image de
$T_V$ dans $|X \times_Y Z'|$ est contenue dans $|T'|$, il s'ensuit
immédiatement que $z' = a(v)$ appartient à l'adhérence de l'image de $T'$.

\medskip\noindent
Il est clair que (1) implique (3). Soit $V \to Y$ comme dans (3).
Si l'on montre que $X \times_Y V \to V$ est universellement fermé,
alors $f$ est universellement fermé d'après
Morphismes d'espaces, lemme
\ref{spaces-morphisms-lemma-universally-closed-local}.
Il suffit donc de montrer que $f : X \to Y$ vérifie (2)
si $f$ est un morphisme quasi-compact d'espaces algébriques,
si $Y$ est un schéma et si $|\mathbf{A}^n \times X| \to |\mathbf{A}^n \times Y|$
est fermée pour tout $n$. Soit $Z \to Y$ localement de présentation finie.
Il faut montrer que l'application $|X \times_Y Z| \to |Z|$ est fermée.
Cette question est locale pour la topologie \'etale sur $Z$; on peut donc supposer $Z$
affine (certains détails sont omis). Puisque $Y$ est un schéma, $Z$ est affine
et $Z \to Y$ est localement de présentation finie, on peut trouver
une immersion $Z \to \mathbf{A}^n \times Y$; voir
Morphismes, lemme \ref{morphisms-lemma-quasi-affine-finite-type-over-S}.
Considérons le diagramme cartésien
$$
\vcenter{
\xymatrix{
X \times_Y Z \ar[d] \ar[r] & \mathbf{A}^n \times X \ar[d] \\
Z \ar[r] & \mathbf{A}^n \times Y
}
}
\quad
\begin{matrix}
\text{qui induit le} \\
\text{carré cartésien}
\end{matrix}
\quad
\vcenter{
\xymatrix{
|X \times_Y Z| \ar[d] \ar[r] & |\mathbf{A}^n \times X| \ar[d] \\
|Z| \ar[r] & |\mathbf{A}^n \times Y|
}
}
$$
d'espaces topologiques, dont les flèches horizontales sont des homéomorphismes
sur des parties localement fermées (Propriétés des espaces, lemme
\ref{spaces-properties-lemma-subspace-induced-topology}).
Ainsi, toute partie fermée $T$
de $|X \times_Y Z|$ est l'image réciproque d'une partie fermée $T'$ de
$|\mathbf{A}^n \times Y|$. Comme l'hypothèse affirme que l'image
de $T'$ dans $|\mathbf{A}^n \times X|$ est fermée, on conclut que
l'image de $T$ dans $|Z|$ est fermée, comme voulu.
\end{proof}

\begin{lemma}
\label{lemma-limited-base-change}
Soit $S$ un schéma. Soit $f : X \to Y$ un
morphisme d'espaces algébriques sur $S$.
Supposons $f$ séparé et de type fini.
Les conditions suivantes sont équivalentes :
\begin{enumerate}
\item Le morphisme $f$ est propre.
\item Pour tout morphisme $Y \to Z$ localement de présentation finie,
l'application $|X \times_Y Z| \to |Z|$ est fermée, et
\item il existe un schéma $V$ et un morphisme \'etale surjectif $V \to Y$
de sorte que l'application $|\mathbf{A}^n \times (X \times_Y V)| \to |\mathbf{A}^n \times V|$
soit fermée pour tout $n \geq 0$.
\end{enumerate}
\end{lemma}

\begin{proof}
Puisqu'un morphisme propre est exactement un morphisme
séparé, de type fini et universellement fermé, ce
lemme est un cas particulier du lemme \ref{lemma-test-universally-closed}.
\end{proof}


\section{Critère valuatif noethérien}
\label{section-Noetherian-valuative-criterion}

\noindent
Nous avons déjà démontré certains résultats dans Cohomologie des espaces, section
\ref{spaces-cohomology-section-Noetherian-valuative-criterion}.
La section correspondante pour les schémas est
Limites, section \ref{limits-section-Noetherian-valuative-criterion}.

\medskip\noindent
Bon nombre des résultats de cette section peuvent (et devraient peut-être)
se démontrer en invoquant le lemme suivant, bien que nous ne l'ayons pas
toujours fait.

\begin{lemma}
\label{lemma-reach-point-closure-Noetherian}
Soit $S$ un schéma. Soit $f : X \to Y$ un morphisme d'espaces algébriques
sur $S$. Supposons $f$ de type fini et $Y$ localement noethérien.
Soit $y \in |Y|$ un point de l'adhérence de l'image de $|f|$.
Alors il existe un diagramme commutatif
$$
\xymatrix{
\Spec(K) \ar[r] \ar[d] & X \ar[d]^f \\
\Spec(A) \ar[r] & Y
}
$$
où $A$ est un anneau de valuation discrète et $K$ son corps des fractions,
et où le point fermé de $\Spec(A)$ a pour image $y$. En outre, nous pouvons supposer
que le point $x \in |X|$ correspondant à $\Spec(K) \to X$ est un point de
codimension $0$\footnote{Voir la discussion dans
Propriétés des espaces, section \ref{spaces-properties-section-generic-points}.}
et que $K$ est le corps résiduel d'un point
d'un schéma \'etale sur $X$.
\end{lemma}

\begin{proof}
Choisissons un schéma affine $V$, un point $v \in V$ et un morphisme \'etale
$V \to Y$ qui envoie $v$ sur $y$. L'application $|V| \to |Y|$ est ouverte et, d'après
Propriétés des espaces, lemme \ref{spaces-properties-lemma-points-cartesian},
l'image de $|X \times_Y V| \to |V|$ est l'image inverse de
celle de $|f|$. Nous en concluons que le point $v$ appartient à l'adhérence de
l'image de $|X \times_Y V| \to |V|$. Si nous démontrons le lemme pour
$X \times_Y V \to V$ et le point $v$, alors le lemme en résulte pour
$f$ et $y$. Nous nous ramenons ainsi à la situation décrite au
paragraphe suivant.

\medskip\noindent
Supposons donnés $f : X \to Y$ et $y \in |Y|$ comme dans le lemme, où
$Y$ est un schéma affine. Comme $f$ est quasi-compact, nous en déduisons que
$X$ est quasi-compact. Nous pouvons donc choisir un schéma affine $W$ et
un morphisme \'etale surjectif $W \to X$. Alors l'image de
$|f|$ est égale à celle de $W \to Y$. Nous nous ramenons ainsi
au cas des schémas, traité dans
Limites, lemme \ref{limits-lemma-reach-point-closure-Noetherian}.
\end{proof}

\noindent
Énonçons d'abord le résultat concernant
la séparation. Nous utiliserons souvent des diagrammes commutatifs à flèches pleines de
morphismes d'espaces algébriques sur un schéma de base $S$, de la forme suivante :
\begin{equation}
\label{equation-valuative}
\vcenter{
\xymatrix{
\Spec(K) \ar[r] \ar[d] & X \ar[d] \\
\Spec(A) \ar[r] \ar@{-->}[ru] & Y
}
}
\end{equation}
où $A$ est un anneau de valuation et $K$ son corps des fractions.

\begin{lemma}
\label{lemma-Noetherian-dvr-valuative-separation}
Soit $S$ un schéma. Soit $f : X \to Y$ un morphisme d'espaces algébriques
sur $S$. Supposons $f$ quasi-séparé et localement de type fini, et
$Y$ localement noethérien. Les conditions suivantes sont équivalentes :
\begin{enumerate}
\item Le morphisme $f$ est séparé.
\item Dans tout diagramme (\ref{equation-valuative}), il existe au plus
une flèche en tirets.
\item Dans tout diagramme (\ref{equation-valuative}) où $A$ est un anneau de
valuation discrète, il existe au plus une flèche en tirets.
\item Dans tout diagramme (\ref{equation-valuative}) où $A$ est un anneau de
valuation discrète et où l'image de $\Spec(K) \to X$ est un point de
codimension $0$ de $X$, il existe au plus une flèche en tirets.
\end{enumerate}
\end{lemma}

\begin{proof}
Nous avons (1) $\Rightarrow$ (2) d'après
Morphismes d'espaces, lemme
\ref{spaces-morphisms-lemma-separated-implies-valuative}.
Les implications (2) $\Rightarrow$ (3) et (3) $\Rightarrow$ (4)
sont immédiates. Il reste à montrer que (4) implique (1).

\medskip\noindent
Supposons (4). Nous devons montrer que la diagonale
$\Delta : X \to X \times_Y X$ est une immersion
fermée. Nous savons déjà que $\Delta$ est représentable, séparé et
localement de type fini, et que c'est un monomorphisme ; voir Morphismes d'espaces, lemme
\ref{spaces-morphisms-lemma-properties-diagonal}.
Choisissons un schéma affine $U$ et un morphisme \'etale
$U \to X \times_Y X$. Posons $V = X \times_{\Delta, X \times_Y X} U$.
Il suffit de montrer que $V \to U$ est une immersion fermée
(Morphismes d'espaces, lemme
\ref{spaces-morphisms-lemma-closed-immersion-local}).
Comme $X \times_Y X$ est localement de type fini sur $Y$, nous voyons que
$U$ est noethérien (utiliser Morphismes d'espaces, lemmes
\ref{spaces-morphisms-lemma-composition-finite-type},
\ref{spaces-morphisms-lemma-base-change-finite-type} et
\ref{spaces-morphisms-lemma-locally-finite-type-locally-noetherian}).
Notons que $V$ est un schéma puisque $\Delta$ est représentable.
De plus, $V$ est quasi-compact parce que $f$ est quasi-séparé.
Ainsi, $V \to U$ est séparé et de type fini. Considérons un diagramme commutatif
$$
\xymatrix{
\Spec(K) \ar[r] \ar[d] & V \ar[d] \\
\Spec(A) \ar[r] \ar@{-->}[ru] & U
}
$$
de morphismes de schémas, où $A$ est un anneau de valuation discrète
de corps des fractions $K$, et où $K$ est le corps résiduel
d'un point générique du schéma noethérien $V$. Comme $V \to X$
est \'etale (en tant que changement de base du morphisme \'etale $U \to X \times_Y X$),
nous voyons que l'image de $\Spec(K) \to V \to X$ est un point de
codimension $0$ ; voir Propriétés des espaces, section
\ref{spaces-properties-section-dimension-local-ring}.
Nous pouvons interpréter le composé $\Spec(A) \to U \to X \times_Y X$
comme une paire de morphismes $a, b : \Spec(A) \to X$ qui coïncident comme morphismes
vers $Y$ et sont égaux après restriction à $\Spec(K)$, leur restriction commune
ayant pour image un point de codimension $0$. Notre hypothèse
(4) garantit donc $a = b$, et nous obtenons la flèche en tirets du diagramme.
D'après Limites, lemme \ref{limits-lemma-Noetherian-dvr-valuative-proper},
nous concluons que $V \to U$ est propre. Autrement dit, $\Delta$ est propre.
Comme $\Delta$ est un monomorphisme, nous en déduisons que $\Delta$ est une
immersion fermée (Morphismes \'etales, lemme
\ref{etale-lemma-characterize-closed-immersion}), comme voulu.
\end{proof}

\begin{lemma}
\label{lemma-Noetherian-dvr-valuative-proper}
Soit $S$ un schéma. Soit $f : X \to Y$ un morphisme d'espaces algébriques
sur $S$. Supposons $f$ quasi-séparé et de type fini, et $Y$
localement noethérien. Les conditions suivantes sont équivalentes :
\begin{enumerate}
\item $f$ est propre,
\item $f$ satisfait au critère valuatif, voir
Morphismes d'espaces, définition
\ref{spaces-morphisms-definition-valuative-criterion},
\item dans tout diagramme (\ref{equation-valuative}), il existe exactement
une flèche en tirets,
\item dans tout diagramme (\ref{equation-valuative}) où $A$ est un anneau de
valuation discrète, il existe exactement une flèche en tirets, et
\item dans tout diagramme (\ref{equation-valuative}) où $A$ est un anneau de
valuation discrète et où l'image de $\Spec(K) \to X$ est un point de
codimension $0$ de $X$, il existe exactement une flèche en tirets\footnote{Il existe
une formulation plus précise où, pour la partie d'existence, on demande seulement que
la flèche en tirets existe après une extension d'anneaux de valuation discrète.}.
\end{enumerate}
\end{lemma}

\begin{proof}
Nous avons (1) $\Leftrightarrow$ (2) $\Leftrightarrow$ (3) d'après 
Morphismes d'espaces, lemme \ref{spaces-morphisms-lemma-characterize-proper}.
Il est clair que (3) $\Rightarrow$ (4) $\Rightarrow$ (5).
Pour achever la démonstration, montrons que (5) implique (1).

\medskip\noindent
Supposons (5). D'après le lemme \ref{lemma-Noetherian-dvr-valuative-separation},
nous voyons que $f$ est séparé. Pour achever la démonstration, il suffit de montrer
que $f$ est universellement fermé. Soit $V \to Y$ un morphisme \'etale,
où $V$ est un schéma affine. Il suffit de montrer que le changement de
base $V \times_Y X \to V$ est universellement fermé ; voir
Morphismes d'espaces, lemme
\ref{spaces-morphisms-lemma-universally-closed-local}.
Considérons le diagramme
$$
\xymatrix{
\Spec(K) \ar[r] \ar[d] & V \times_Y X \ar[d] \ar[r] & X \ar[d] \\
\Spec(A) \ar[r] \ar@{-->}[ru] \ar@{..>}[rru] & V \ar[r] & Y
}
$$
ci-dessus, formé de morphismes d'espaces algébriques sur $S$,
qui est commutatif, où $A$ est un anneau de valuation discrète
de corps des fractions $K$ et où $\Spec(K) \to V \times_Y X$
a pour image un point de codimension $0$ de l'espace algébrique
$V \times_Y X$. Puisque $V \times_Y X \to X$ est \'etale, il s'ensuit
que l'image de $\Spec(K) \to X$ est un point de codimension $0$
de $X$. Ainsi, (5) nous donne la plus longue des deux flèches non pleines,
qui est pointillée et s'insère dans le diagramme. Nous obtenons alors bien sûr aussi
la plus courte, qui est en tirets. Il s'ensuit que nos hypothèses valent pour
le morphisme $V \times_Y X \to V$, et nous sommes ramenés au cas traité
au paragraphe suivant.

\medskip\noindent
Supposons que $Y$ soit un schéma affine noethérien. Dans ce cas, $X$ est un
espace algébrique noethérien séparé (nous savons déjà que $f$ est séparé)
de type fini sur $Y$. (En particulier, l'espace algébrique $X$
possède un sous-espace ouvert dense qui est un schéma d'après Propriétés des espaces, proposition
\ref{spaces-properties-proposition-locally-quasi-separated-open-dense-scheme},
bien qu'à strictement parler nous n'en ayons pas besoin.)
Choisissons un schéma quasi-projectif $X'$ sur $Y$ et un morphisme propre surjectif
$X' \to X$ comme dans la forme faible du lemme de Chow
(Cohomologie des espaces, lemme \ref{spaces-cohomology-lemma-weak-chow}).
Nous pouvons remplacer $X'$ par la réunion disjointe des composantes irréductibles
qui dominent une composante irréductible de $X$ ; nous omettons les détails.
En particulier, nous pouvons supposer que les points génériques du schéma $X'$ ont pour image
des points de codimension $0$ de $X$ (dans ce cas, ce sont exactement
les points génériques de $X$). Nous affirmons que $X' \to Y$ est propre.
Cette affirmation implique que $X$ est propre sur $Y$ d'après
Morphismes d'espaces, lemme \ref{spaces-morphisms-lemma-image-proper-is-proper}.
Pour la démontrer, d'après
Limites, lemme \ref{limits-lemma-Noetherian-dvr-valuative-proper},
il suffit de prouver que, dans tout diagramme commutatif à flèches pleines
$$
\xymatrix{
\Spec(K) \ar[r] \ar[d] & X' \ar[r] & X \ar[d] \\
\Spec(A) \ar[rr] \ar@{-->}[ru]^a \ar@{-->}[rru]_b & & Y
}
$$
où $A$ est un anneau de valuation discrète de corps des fractions $K$, et où $K$ est le corps résiduel
d'un point générique de $X'$, nous pouvons trouver la
flèche en tirets $a$ (nous savons déjà qu'elle est unique puisque $X'$ est séparé).
L'hypothèse (5) nous fournit la flèche en tirets $b$.
Alors le morphisme $X' \times_{X, b} \Spec(A) \to \Spec(A)$
est un morphisme propre de schémas et, par le critère valuatif
pour les morphismes de schémas, le relèvement de $b$ fournit le morphisme $a$ recherché.
\end{proof}

\begin{lemma}
\label{lemma-check-universally-closed-Noetherian}
Soit $S$ un schéma. Soit $f : X \to Y$ un morphisme d'espaces algébriques
sur $S$. Supposons $Y$ localement noethérien et $f$
de type fini. Les conditions suivantes sont équivalentes :
\begin{enumerate}
\item $f$ est universellement fermé,
\item $f$ satisfait à la partie d'existence du critère valuatif,
\item il existe un schéma $V$ et un morphisme \'etale surjectif
$V \to Y$ tels que
$|\mathbf{A}^n \times X \times_Y V| \to |\mathbf{A}^n \times V|$ soit fermée
pour tout $n \geq 0$,
\item dans tout diagramme (\ref{equation-valuative}) où $A$ est un anneau de
valuation discrète, il existe une extension de corps finie séparable $K'/K$,
un anneau de valuation discrète $A' \subset K'$ qui domine $A$, et
un morphisme $\Spec(A') \to X$ tel que le diagramme suivant soit commutatif :
$$
\xymatrix{
\Spec(K') \ar[r] \ar[d] & \Spec(K) \ar[r] & X \ar[d] \\
\Spec(A') \ar[r] \ar[rru] & \Spec(A) \ar[r] & Y
}
$$
\item dans tout diagramme (\ref{equation-valuative}) où $A$ est un anneau de
valuation discrète, il existe une extension de corps $K'/K$,
un anneau de valuation $A' \subset K'$ qui domine $A$, et
un morphisme $\Spec(A') \to X$ tel que le diagramme suivant soit commutatif :
$$
\xymatrix{
\Spec(K') \ar[r] \ar[d] & \Spec(K) \ar[r] & X \ar[d] \\
\Spec(A') \ar[r] \ar[rru] & \Spec(A) \ar[r] & Y
}
$$
\end{enumerate}
\end{lemma}

\begin{proof}
Les conditions (1), (2) et (3) sont équivalentes d'après
le lemme \ref{lemma-test-universally-closed} et
Morphismes d'espaces, lemme
\ref{spaces-morphisms-lemma-quasi-compact-existence-universally-closed}.
Ces conditions équivalentes impliquent (4), car
Morphismes d'espaces, lemme \ref{spaces-morphisms-lemma-finite-separable-enough},
nous dit que nous pouvons toujours choisir l'extension $K'/K$ finie séparable dans
la partie d'existence du critère valuatif, ce qui force automatiquement
$A'$ à être un anneau de valuation discrète par le théorème de Krull-Akizuki
(Algèbre, lemme \ref{algebra-lemma-krull-akizuki}).
L'implication (4) $\Rightarrow$ (5) est immédiate.
Dans le reste de la démonstration, nous montrons que (5) implique (1).

\medskip\noindent
Supposons (5). Choisissons un schéma affine $V$ et un morphisme \'etale $V \to Y$.
Il suffit de montrer que le changement de base de $f$ à $V$ est universellement fermé ;
voir Morphismes d'espaces, lemme
\ref{spaces-morphisms-lemma-universally-closed-local}.
Exactement comme dans la démonstration du lemme \ref{lemma-Noetherian-dvr-valuative-proper},
nous voyons que l'hypothèse (5) se transmet à ce changement de base ; nous omettons
les détails. Nous sommes ainsi ramenés au cas traité au paragraphe suivant.

\medskip\noindent
Supposons que $Y$ soit un schéma affine noethérien et que (5) soit satisfaite.
Pour prouver que $f$ est universellement fermé, il suffit de montrer que
$|X \times \mathbf{A}^n| \to |Y \times \mathbf{A}^n|$ est fermée
pour tout $n$ (d'après la discussion ci-dessus). Comme l'hypothèse (5)
se transmet au morphisme produit
$X \times \mathbf{A}^n \to Y \times \mathbf{A}^n$ (détails omis),
nous sommes ramenés à prouver que $|X| \to |Y|$ est fermée.

\medskip\noindent
Supposons que $Y$ soit un schéma affine noethérien et que (5) soit satisfaite.
Soit $T \subset |X|$ une partie fermée. Nous devons montrer que
l'image de $T$ dans $|Y|$ est fermée. Nous pouvons remplacer $X$
par la structure réduite de sous-espace fermé induite sur $T$ ; nous
omettons de vérifier que la propriété (5) est préservée par
ce remplacement. Nous sommes ainsi ramenés à prouver que l'image
de $|X| \to |Y|$ est fermée.

\medskip\noindent
Soit $y \in |Y|$ un point de l'adhérence de l'image de
$|X| \to |Y|$. D'après le lemme \ref{lemma-reach-point-closure-Noetherian},
nous pouvons choisir un diagramme commutatif
$$
\xymatrix{
\Spec(K) \ar[r] \ar[d] & X \ar[d]^f \\
\Spec(A) \ar[r] & Y
}
$$
où $A$ est un anneau de valuation discrète et $K$ son corps des fractions,
et où le point fermé de $\Spec(A)$ a pour image $y$. Il résulte immédiatement
de la propriété (5) que $y$ appartient à l'image de
$|X| \to |Y|$, ce qui achève la démonstration.
\end{proof}


















\section{Critères valuatifs noethériens raffinés}
\label{section-refined-valuative-criteria}

\noindent
Cette section est l'analogue de
Limites, section \ref{limits-section-refined-valuative-criteria}.
Pour v\'erifier les crit\`eres valuatifs, il n'est en g\'en\'eral pas n\'ecessaire
de consid\'erer tous les diagrammes possibles faisant intervenir des anneaux de valuation.

\begin{lemma}
\label{lemma-refined-valuative-criterion-proper}
Soit $S$ un sch\'ema. Soient $f : X \to Y$ et $h : U \to X$ des
morphismes d'espaces alg\'ebriques sur $S$. Supposons que $Y$ soit
localement noeth\'erien, que $f$ et $h$ soient de type fini,
que $f$ soit s\'epar\'e et que l'image de $|h| : |U| \to |X|$
soit dense dans $|X|$. Si, pour tout diagramme commutatif de fl\`eches pleines
$$
\xymatrix{
\Spec(K) \ar[r] \ar[d] & U \ar[r]^h & X \ar[d]^f \\
\Spec(A) \ar[rr] \ar@{-->}[rru] & & Y
}
$$
o\`u $A$ est un anneau de valuation discr\`ete de corps des fractions $K$, il
existe une fl\`eche pointill\'ee rendant le diagramme commutatif, alors $f$ est propre.
\end{lemma}

\begin{proof}
Il suffit de montrer que $f$ est universellement ferm\'e.
Soit $V \to Y$ un morphisme \'etale, o\`u $V$ est un sch\'ema affine.
Par Morphismes d'espaces, lemme
\ref{spaces-morphisms-lemma-universally-closed-local}
il suffit de montrer que le changement de base $X \times_Y V \to V$ est
universellement ferm\'e. Par Propri\'et\'es des espaces, lemme
\ref{spaces-properties-lemma-points-cartesian}
l'image $I$ de $|U \times_Y V| \to |X \times_Y V|$
est l'image r\'eciproque de l'image de $|h|$. Puisque
$|X \times_Y V| \to |X|$ est ouverte
(Propri\'et\'es des espaces, lemme \ref{spaces-properties-lemma-etale-open}),
on en d\'eduit que $I$ est dense dans $|X \times_Y V|$.
Les hypoth\`eses du lemme sont donc
satisfaites pour les morphismes $U \times_Y V \to X \times_Y V \to V$.
On peut ainsi supposer que $Y$ est un sch\'ema affine.

\medskip\noindent
Supposons que $Y$ soit un sch\'ema affine. Alors $U$ est quasi-compact. Choisissons
un sch\'ema affine et un morphisme \'etale surjectif $W \to U$.
On remplace alors $U$ par $W$ et l'on suppose que $U$ est affine.
Par la version faible du lemme de Chow
(Cohomologie des espaces, lemme \ref{spaces-cohomology-lemma-weak-chow})
on peut choisir un morphisme propre surjectif $X' \to X$
o\`u $X'$ est un sch\'ema. Alors $U' = X' \times_X U$ est un sch\'ema
et $U' \to X'$ est de type fini. On peut remplacer $X'$ par
l'image sch\'ematique de $h' : U' \to X'$; ainsi $h'(U')$
est dense dans $X'$. Nous affirmons que, pour tout diagramme
$$
\xymatrix{
\Spec(K) \ar[r] \ar[d] & U' \ar[r]^h & X' \ar[d]^{f'} \\
\Spec(A) \ar[rr] \ar@{-->}[rru] & & Y
}
$$
o\`u $A$ est un anneau de valuation discr\`ete de corps des fractions $K$, il
existe une fl\`eche pointill\'ee rendant le diagramme commutatif. En effet, l'hypoth\`ese du
lemme fournit d'abord une fl\`eche $\Spec(A) \to X$, que l'on rel\`eve ensuite
en une fl\`eche $\Spec(A) \to X'$ au moyen du crit\`ere valuatif de propret\'e
(Morphismes d'espaces, lemme \ref{spaces-morphisms-lemma-characterize-proper}).
Le morphisme $X' \to Y$ est s\'epar\'e
comme compos\'e d'un morphisme propre et d'un morphisme s\'epar\'e.
Le cas des sch\'emas montre donc que le morphisme $X' \to Y$ est propre
(Limites, lemme
\ref{limits-lemma-refined-valuative-criterion-proper}).
Par Morphismes d'espaces, lemme
\ref{spaces-morphisms-lemma-image-proper-is-proper}
on en conclut que $X \to Y$ est propre.
\end{proof}

\begin{lemma}
\label{lemma-refined-valuative-criterion-separated}
Soit $S$ un sch\'ema. Soient $f : X \to Y$ et $h : U \to X$ des
morphismes d'espaces alg\'ebriques sur $S$. Supposons que $Y$ soit
localement noeth\'erien, que $f$ soit localement de type fini et quasi-s\'epar\'e,
que $h$ soit de type fini et que l'image de $|h| : |U| \to |X|$
soit dense dans $|X|$.
Si, pour tout diagramme commutatif de fl\`eches pleines
$$
\xymatrix{
\Spec(K) \ar[r] \ar[d] & U \ar[r]^h & X \ar[d]^f \\
\Spec(A) \ar[rr] \ar@{-->}[rru] & & Y
}
$$
o\`u $A$ est un anneau de valuation discr\`ete de corps des fractions $K$, il
existe au plus une fl\`eche pointill\'ee rendant le diagramme commutatif, alors $f$ est
s\'epar\'e.
\end{lemma}

\begin{proof}
Nous appliquerons le Lemme \ref{lemma-refined-valuative-criterion-proper}
aux morphismes $U \to X$ et $\Delta : X \to X \times_Y X$.
V\'erifions les conditions. Remarquons que $\Delta$ est quasi-compact puisque
$f$ est quasi-s\'epar\'e. Bien entendu, $\Delta$ est localement de type fini et
s\'epar\'e (comme tout morphisme diagonal).
Enfin, supposons donn\'e un diagramme commutatif de fl\`eches pleines
$$
\xymatrix{
\Spec(K) \ar[r] \ar[d] & U \ar[r]^h & X \ar[d]^\Delta \\
\Spec(A) \ar[rr]^{(a, b)} \ar@{-->}[rru] & & X \times_Y X
}
$$
o\`u $A$ est un anneau de valuation discr\`ete de corps des fractions $K$.
Alors $a$ et $b$ fournissent deux fl\`eches pointill\'ees dans le diagramme du lemme
et doivent \^etre \`egales. On peut donc prendre $a = b$ comme fl\`eche pointill\'ee,
ce qui donne l'existence et ach\`eve la d\'emonstration.
\end{proof}

\begin{lemma}
\label{lemma-refined-valuative-criterion-universally-closed}
Soit $S$ un sch\'ema.
Soient $f : X \to Y$ et $h : U \to X$ des morphismes d'espaces alg\'ebriques sur $S$.
Supposons que $Y$ soit localement noeth\'erien, que $f$ et $h$ soient de type fini,
que $f$ soit quasi-s\'epar\'e et
que $h(U)$ soit dense dans $X$. Si, pour tout diagramme commutatif de fl\`eches pleines
$$
\xymatrix{
\Spec(K) \ar[r] \ar[d] & U \ar[r]^h & X \ar[d]^f \\
\Spec(A) \ar[rr] \ar@{-->}[rru] & & Y
}
$$
o\`u $A$ est un anneau de valuation discr\`ete de corps des fractions $K$, il
existe une unique fl\`eche pointill\'ee rendant le diagramme commutatif, alors $f$ est propre.
\end{lemma}

\begin{proof}
Combiner les Lemmes \ref{lemma-refined-valuative-criterion-separated} et
\ref{lemma-refined-valuative-criterion-proper}.
\end{proof}







\section{Descente des espaces de type fini}
\label{section-finite-type-quasi-separated}

\noindent
Cette section poursuit le th\`eme de la
section \ref{section-finite-type-closed-in-finite-presentation}
dans l'esprit des r\'esultats examin\'es dans la
section \ref{section-descending-relative}.
Elle est aussi l'analogue de
Limites, section \ref{limits-section-finite-type-quasi-separated},
pour les espaces alg\'ebriques.

\begin{situation}
\label{situation-limit-noetherian}
Soit $S$ un sch\'ema, par exemple $\Spec(\mathbf{Z})$.
Soit $B = \lim_{i \in I} B_i$ la limite
d'un syst\`eme projectif filtrant d'espaces noeth\'eriens sur $S$
\`a morphismes de transition affines
$B_{i'} \to B_i$ pour $i' \geq i$.
\end{situation}

\begin{lemma}
\label{lemma-good-diagram}
Dans la Situation \ref{situation-limit-noetherian}.
Soit $X \to B$ un morphisme d'espaces alg\'ebriques
quasi-s\'epar\'e et de type fini.
Alors il existe $i \in I$ et un diagramme
\begin{equation}
\label{equation-good-diagram}
\vcenter{
\xymatrix{
X \ar[r] \ar[d] & W \ar[d] \\
B \ar[r] & B_i
}
}
\end{equation}
tels que $W \to B_i$ soit de type fini et que
le morphisme induit $X \to B \times_{B_i} W$ soit une immersion
ferm\'ee.
\end{lemma}

\begin{proof}
Par le Lemme \ref{lemma-finite-type-closed-in-finite-presentation},
il existe une immersion ferm\'ee $X \to X'$
sur $B$, o\`u $X'$ est un espace alg\'ebrique de pr\'esentation finie sur $B$.
Par le Lemme \ref{lemma-descend-finite-presentation},
il existe $i$ et un morphisme de pr\'esentation finie
$X'_i \to B_i$ dont le changement de base est $X'$. Posons $W = X'_i$.
\end{proof}

\begin{lemma}
\label{lemma-limit-from-good-diagram}
Dans la Situation \ref{situation-limit-noetherian}.
Soit $X \to B$ un morphisme d'espaces alg\'ebriques
quasi-s\'epar\'e et de type fini. \`Etant donn\'es $i \in I$ et un diagramme
$$
\vcenter{
\xymatrix{
X \ar[r] \ar[d] & W \ar[d] \\
B \ar[r] & B_i
}
}
$$
comme dans (\ref{equation-good-diagram}), pour $i' \geq i$, soit
$X_{i'}$ l'image sch\'ematique de $X \to B_{i'} \times_{B_i} W$.
Alors $X = \lim_{i' \geq i} X_{i'}$.
\end{lemma}

\begin{proof}
Puisque $X$ est quasi-compact et quasi-s\'epar\'e, la formation de
l'image sch\'ematique de $X \to B_{i'} \times_{B_i} W$
commute \`a la localisation \'etale
(Morphismes d'espaces, lemme
\ref{spaces-morphisms-lemma-quasi-compact-scheme-theoretic-image}).
On peut donc supposer que $W$ est affine et s'envoie dans un sch\'ema affine
$U_i$ \'etale sur $B_i$. Alors
$$
B_{i'} \times_{B_i} W =
B_{i'} \times_{B_i} U_i \times_{U_i} W =
U_{i'} \times_{U_i} W
$$
o\`u $U_{i'} = B_{i'} \times_{B_i} U_i$ est affine puisque les morphismes de
transition sont affines. Le lemme r\'esulte donc du cas des sch\'emas,
qui est
Limites, lemme \ref{limits-lemma-limit-from-good-diagram}.
\end{proof}

\begin{lemma}
\label{lemma-morphism-good-diagram}
Dans la Situation \ref{situation-limit-noetherian}.
Soit $f : X \to Y$ un morphisme d'espaces alg\'ebriques quasi-s\'epar\'es
et de type fini sur $B$. Soient
$$
\vcenter{
\xymatrix{
X \ar[r] \ar[d] & W \ar[d] \\
B \ar[r] & B_{i_1}
}
}
\quad\text{et}\quad
\vcenter{
\xymatrix{
Y \ar[r] \ar[d] & V \ar[d] \\
B \ar[r] & B_{i_2}
}
}
$$
des diagrammes comme dans (\ref{equation-good-diagram}). Soient
$X = \lim_{i \geq i_1} X_i$ et
$Y = \lim_{i \geq i_2} Y_i$ les descriptions correspondantes
comme limites, donn\'ees par le Lemme \ref{lemma-limit-from-good-diagram}.
Alors il existe $i_0 \geq \max(i_1, i_2)$ et un morphisme
$$
(f_i)_{i \geq i_0} : (X_i)_{i \geq i_0} \to (Y_i)_{i \geq i_0}
$$
de syst\`emes projectifs sur $(B_i)_{i \geq i_0}$ tel que
$f = \lim_{i \geq i_0} f_i$.
Si $(g_i)_{i \geq i_0} : (X_i)_{i \geq i_0} \to (Y_i)_{i \geq i_0}$
est un second morphisme de syst\`emes projectifs sur $(B_i)_{i \geq i_0}$ tel que
$f = \lim_{i \geq i_0} g_i$,
alors $f_i = g_i$ pour tout $i \gg i_0$.
\end{lemma}

\begin{proof}
Puisque $V \to B_{i_2}$ est de pr\'esentation finie et que
$X = \lim_{i \geq i_1} X_i$, on peut invoquer la Proposition
\ref{proposition-characterize-locally-finite-presentation}
dans la forme pr\'ecis\'ee par le Lemme \ref{lemma-better-characterize-relative-limit-preserving}
pour trouver $i_0 \geq \max(i_1, i_2)$ et un morphisme $h : X_{i_0} \to V$
sur $B_{i_2}$ tel que $X \to X_{i_0} \to V$ soit \`egal \`a $X \to Y \to V$.
Pour $i \geq i_0$, on obtient un diagramme commutatif de fl\`eches pleines
$$
\xymatrix{
X \ar[d] \ar[r] &
X_i \ar[r] \ar@{..>}[d] \ar@/_2pc/[dd] |!{[d];[ld]}\hole &
X_{i_0} \ar[d]^h \\
Y \ar[r] \ar[d] & Y_i \ar[r] \ar[d] & V \ar[d] \\
B \ar[r] & B_i \ar[r] & B_{i_0}
}
$$
Puisque l'image de $X \to X_i$ est sch\'ematiquement dense
et que $Y_i$ est l'image sch\'ematique de
$Y \to B_i \times_{B_{i_2}} V$,
le morphisme $X_i \to B_i \times_{B_{i_2}} V$
induit par le diagramme se factorise par $Y_i$
(Morphismes d'espaces, lemme \ref{spaces-morphisms-lemma-factor-factor}).
Ceci prouve l'existence.

\medskip\noindent
Unicit\'e. Soit $E_i \to X_i$ le noyau du couple $f_i$, $g_i$
pour $i \geq i_0$. On a
$E_i = Y_i \times_{\Delta, Y_i \times_{B_i} Y_i, (f_i, g_i)} X_i$.
Ainsi $E_i \to X_i$ est un monomorphisme de pr\'esentation finie, comme changement
de base de la diagonale de $Y_i$ sur $B_i$; voir
Morphismes d'espaces, lemmes \ref{spaces-morphisms-lemma-properties-diagonal} et
\ref{spaces-morphisms-lemma-diagonal-morphism-finite-type}.
Puisque $X_i$ est un sous-espace ferm\'e de $B_i \times_{B_{i_0}} X_{i_0}$
et qu'il en va de m\^eme pour $Y_i$, on a
$$
E_i =
X_i \times_{(B_i \times_{B_{i_0}} X_{i_0})} (B_i \times_{B_{i_0}} E_{i_0}) =
X_i \times_{X_{i_0}} E_{i_0}
$$
De m\^eme, on a $X = X \times_{X_{i_0}} E_{i_0}$. Le Lemme
\ref{lemma-descend-isomorphism} donne donc $E_i = X_i$ pour $i$ assez grand.
\end{proof}

\begin{remark}
\label{remark-finite-type-gives-well-defined-system}
Dans la Situation \ref{situation-limit-noetherian},
les Lemmes \ref{lemma-good-diagram}, \ref{lemma-limit-from-good-diagram} et
\ref{lemma-morphism-good-diagram}
montrent que la cat\'egorie des espaces alg\'ebriques quasi-s\'epar\'es et
de type fini sur $B$ est \`equivalente \`a certains types de
syst\`emes projectifs d'espaces alg\'ebriques sur $(B_i)_{i \in I}$, \`a savoir
ceux obtenus en appliquant le Lemme \ref{lemma-limit-from-good-diagram}
\`a un diagramme de la forme (\ref{equation-good-diagram}).
Par exemple, soient $X \to B$ de type fini et quasi-s\'epar\'e,
et deux diagrammes distincts $X \to V_1 \to B_{i_1}$
et $X \to V_2 \to B_{i_2}$ comme dans (\ref{equation-good-diagram}). En
appliquant le Lemme \ref{lemma-morphism-good-diagram} \`a $\text{id}_X$
(dans les deux sens),
on voit que les descriptions correspondantes de
$X$ comme limite sont canoniquement isomorphes (quitte \`a restreindre
l'ensemble filtrant $I$). Et ainsi de suite.
\end{remark}

\begin{lemma}
\label{lemma-morphism-good-diagram-flat}
Notations et hypoth\`eses comme dans le Lemme \ref{lemma-morphism-good-diagram}.
Si $f$ est plat et de pr\'esentation finie, alors
il existe $i_3 > i_0$ tel que, pour $i \geq i_3$, le morphisme
$f_i$ soit plat, $X_i = Y_i \times_{Y_{i_3}} X_{i_3}$ et
$X = Y \times_{Y_{i_3}} X_{i_3}$.
\end{lemma}

\begin{proof}
Par le Lemme \ref{lemma-descend-finite-presentation},
on peut choisir $i \geq i_2$ et un morphisme
$U \to Y_i$ de pr\'esentation finie tel que $X = Y \times_{Y_i} U$
(c'est ici que l'on utilise que $f$ est de pr\'esentation finie).
Quitte \`a augmenter $i$, on peut supposer que $U \to Y_i$ est plat; voir le
Lemme \ref{lemma-descend-flat}.
Comme expliqu\'e dans la Remarque \ref{remark-finite-type-gives-well-defined-system},
on remplace le diagramme initial ayant servi \`a d\'efinir le syst\`eme
$(X_i)_{i \geq i_1}$ par le syst\`eme correspondant \`a
$X \to U \to B_i$. Ainsi, $X_{i'}$, pour $i' \geq i$, est d\'efini comme
l'image sch\'ematique de $X \to B_{i'} \times_{B_i} U$.

\medskip\noindent
Puisque $U \to Y_i$ est plat (c'est ici que l'on utilise que $f$ est plat),
que $X = Y \times_{Y_i} U$ et
que l'image sch\'ematique de $Y \to Y_i$ est $Y_i$,
on voit que l'image sch\'ematique de $X \to U$ est $U$
(Morphismes d'espaces, lemme
\ref{spaces-morphisms-lemma-flat-base-change-scheme-theoretic-image}).
Remarquons que $Y_{i'} \to B_{i'} \times_{B_i} Y_i$ est une
immersion ferm\'ee pour $i' \geq i$ par construction du syst\`eme des $Y_j$.
Le m\^eme argument que ci-dessus montre alors que l'image sch\'ematique
de $X \to B_{i'} \times_{B_i} U$
est le sous-espace ferm\'e $Y_{i'} \times_{Y_i} U$.
Ainsi $X_{i'} = Y_{i'} \times_{Y_i} U$ pour tout $i' \geq i$,
et le lemme vaut donc avec $i_3 = i$.
\end{proof}

\begin{lemma}
\label{lemma-morphism-good-diagram-smooth}
Notations et hypoth\`eses comme dans le Lemme \ref{lemma-morphism-good-diagram}.
Si $f$ est lisse, alors il existe $i_3 > i_0$ tel que, pour
$i \geq i_3$, le morphisme $f_i$ soit lisse.
\end{lemma}

\begin{proof}
Combiner les Lemmes \ref{lemma-morphism-good-diagram-flat} et
\ref{lemma-descend-smooth}.
\end{proof}

\begin{lemma}
\label{lemma-morphism-good-diagram-proper}
Notations et hypoth\`eses comme dans le Lemme \ref{lemma-morphism-good-diagram}.
Si $f$ est propre, alors il existe $i_3 \geq i_0$ tel que, pour
$i \geq i_3$, le morphisme $f_i$ soit propre.
\end{lemma}

\begin{proof}
Il r\'esulte de la discussion de la 
Remarque \ref{remark-finite-type-gives-well-defined-system}
que le choix de $i_1$ et de $W$ s'inscrivant dans un diagramme comme
(\ref{equation-good-diagram}) est sans incidence sur la validit\'e du
lemme. On choisit donc $W$ comme suit.
On choisit d'abord une immersion ferm\'ee $X \to X'$
avec $X' \to Y$ propre et de pr\'esentation finie; voir le
Lemme \ref{lemma-proper-limit-of-proper-finite-presentation}.
On choisit ensuite $i_3 \geq i_2$ et un morphisme propre $W \to Y_{i_3}$
tel que $X' = Y \times_{Y_{i_3}} W$. Cela est possible puisque
$Y = \lim_{i \geq i_2} Y_i$ et en vertu des
Lemmes \ref{lemma-relative-approximation} et \ref{lemma-eventually-proper}.
Avec ce choix de $W$, il r\'esulte imm\'ediatement de la construction que,
pour $i \geq i_3$, l'espace alg\'ebrique $X_i$ est un sous-espace ferm\'e de
$Y_i \times_{Y_{i_3}} W \subset B_i \times_{B_{i_3}} W$
et est donc propre sur $Y_i$.
\end{proof}

\begin{lemma}
\label{lemma-good-diagram-fibre-product}
Dans la Situation \ref{situation-limit-noetherian}, supposons donn\'e un
diagramme cart\'esien
$$
\xymatrix{
X^1 \ar[r]_p \ar[d]_q & X^3 \ar[d]^a \\
X^2 \ar[r]^b & X^4
}
$$
d'espaces alg\'ebriques quasi-s\'epar\'es et de type fini sur $B$.
Pour chaque $j = 1, 2, 3, 4$, choisissons $i_j \in I$ et un diagramme
$$
\xymatrix{
X^j \ar[r] \ar[d] & W^j \ar[d] \\
B \ar[r] & B_{i_j}
}
$$
comme dans (\ref{equation-good-diagram}). Soient
$X^j = \lim_{i \geq i_j} X^j_i$ les descriptions correspondantes comme limites,
comme dans le Lemme \ref{lemma-morphism-good-diagram}.
Soient $(a_i)_{i \geq i_5}$, $(b_i)_{i \geq i_6}$, $(p_i)_{i \geq i_7}$ et
$(q_i)_{i \geq i_8}$ les morphismes correspondants de syst\`emes projectifs
construits dans le Lemme \ref{lemma-morphism-good-diagram}. Alors il existe
$i_9 \geq \max(i_5, i_6, i_7, i_8)$ tel que, pour $i \geq i_9$, on ait
$a_i \circ p_i = b_i \circ q_i$ et que
$$
(q_i, p_i) : X^1_i \longrightarrow X^2_i \times_{b_i, X^4_i, a_i} X^3_i
$$
soit une immersion ferm\'ee.
Si $a$ et $b$ sont plats et de pr\'esentation finie, alors il existe
$i_{10} \geq \max(i_5, i_6, i_7, i_8, i_9)$ tel que, pour $i \geq i_{10}$,
le dernier morphisme affich\'e soit un isomorphisme.
\end{lemma}

\begin{proof}
D'apr\`es la discussion de la
Remarque \ref{remark-finite-type-gives-well-defined-system},
le choix de $W^1$ s'inscrivant dans un diagramme comme
(\ref{equation-good-diagram}) est sans incidence sur la validit\'e du
lemme. On peut donc choisir $W^1 = W^2 \times_{W^4} W^3$.
Il r\'esulte alors imm\'ediatement de la construction de $X^1_i$ que 
$a_i \circ p_i = b_i \circ q_i$ et que
$$
(q_i, p_i) : X^1_i \longrightarrow X^2_i \times_{b_i, X^4_i, a_i} X^3_i
$$
est une immersion ferm\'ee.

\medskip\noindent
Si $a$ et $b$ sont plats et de pr\'esentation finie, il en va de m\^eme de
$p$ et $q$, changements de base de $a$ et $b$. On peut donc appliquer le
Lemme \ref{lemma-morphism-good-diagram-flat}
\`a chacun de $a$, $b$, $p$, $q$ et $a \circ p = b \circ q$.
Il en r\'esulte qu'il existe $i_9 \in I$ tel que
$$
(q_i, p_i) : X^1_i \to X^2_i \times_{X^4_i} X^3_i
$$
soit le changement de base de $(q_{i_9}, p_{i_9})$ par le morphisme
$X^4_i \to X^4_{i_9}$ pour tout $i \geq i_9$.
On conclut que $(q_i, p_i)$ est un isomorphisme pour tout $i$ suffisamment
grand par le Lemme \ref{lemma-descend-isomorphism}.
\end{proof}












\begin{multicols}{2}[\section{Autres chapitres}]
\noindent
Préliminaires
\begin{enumerate}
\item \hyperref[introduction-section-phantom]{Introduction}
\item \hyperref[conventions-section-phantom]{Conventions}
\item \hyperref[sets-section-phantom]{Théorie des ensembles}
\item \hyperref[categories-section-phantom]{Catégories}
\item \hyperref[topology-section-phantom]{Topologie}
\item \hyperref[sheaves-section-phantom]{Faisceaux sur les espaces}
\item \hyperref[sites-section-phantom]{Sites et faisceaux}
\item \hyperref[stacks-section-phantom]{Champs}
\item \hyperref[fields-section-phantom]{Corps}
\item \hyperref[algebra-section-phantom]{Algèbre commutative}
\item \hyperref[brauer-section-phantom]{Groupes de Brauer}
\item \hyperref[homology-section-phantom]{Algèbre homologique}
\item \hyperref[derived-section-phantom]{Catégories dérivées}
\item \hyperref[simplicial-section-phantom]{Méthodes simpliciales}
\item \hyperref[more-algebra-section-phantom]{Compléments d'algèbre}
\item \hyperref[smoothing-section-phantom]{Lissage des morphismes d'anneaux}
\item \hyperref[modules-section-phantom]{Faisceaux de modules}
\item \hyperref[sites-modules-section-phantom]{Modules sur les sites}
\item \hyperref[injectives-section-phantom]{Objets injectifs}
\item \hyperref[cohomology-section-phantom]{Cohomologie des faisceaux}
\item \hyperref[sites-cohomology-section-phantom]{Cohomologie sur les sites}
\item \hyperref[dga-section-phantom]{Algèbre différentielle graduée}
\item \hyperref[dpa-section-phantom]{Algèbre à puissances divisées}
\item \hyperref[sdga-section-phantom]{Faisceaux différentiels gradués}
\item \hyperref[hypercovering-section-phantom]{Hyperrecouvrements}
\end{enumerate}
Schémas
\begin{enumerate}
\setcounter{enumi}{25}
\item \hyperref[schemes-section-phantom]{Schémas}
\item \hyperref[constructions-section-phantom]{Constructions de schémas}
\item \hyperref[properties-section-phantom]{Propriétés des schémas}
\item \hyperref[morphisms-section-phantom]{Morphismes de schémas}
\item \hyperref[coherent-section-phantom]{Cohomologie des schémas}
\item \hyperref[divisors-section-phantom]{Diviseurs}
\item \hyperref[limits-section-phantom]{Limites de schémas}
\item \hyperref[varieties-section-phantom]{Variétés}
\item \hyperref[topologies-section-phantom]{Topologies sur les schémas}
\item \hyperref[descent-section-phantom]{Descente}
\item \hyperref[perfect-section-phantom]{Catégories dérivées de schémas}
\item \hyperref[more-morphisms-section-phantom]{Compléments sur les morphismes}
\item \hyperref[flat-section-phantom]{Compléments sur la platitude}
\item \hyperref[groupoids-section-phantom]{Schémas en groupoïdes}
\item \hyperref[more-groupoids-section-phantom]{Compléments sur les schémas en groupoïdes}
\item \hyperref[etale-section-phantom]{Morphismes étales de schémas}
\end{enumerate}
Thèmes de théorie des schémas
\begin{enumerate}
\setcounter{enumi}{41}
\item \hyperref[chow-section-phantom]{Homologie de Chow}
\item \hyperref[intersection-section-phantom]{Théorie de l'intersection}
\item \hyperref[pic-section-phantom]{Schémas de Picard des courbes}
\item \hyperref[weil-section-phantom]{Théories cohomologiques de Weil}
\item \hyperref[adequate-section-phantom]{Modules adéquats}
\item \hyperref[dualizing-section-phantom]{Complexes dualisants}
\item \hyperref[duality-section-phantom]{Dualité pour les schémas}
\item \hyperref[discriminant-section-phantom]{Discriminants et différentes}
\item \hyperref[derham-section-phantom]{Cohomologie de de Rham}
\item \hyperref[local-cohomology-section-phantom]{Cohomologie locale}
\item \hyperref[algebraization-section-phantom]{Géométrie algébrique et formelle}
\item \hyperref[curves-section-phantom]{Courbes algébriques}
\item \hyperref[resolve-section-phantom]{Résolution des surfaces}
\item \hyperref[models-section-phantom]{Réduction semi-stable}
\item \hyperref[functors-section-phantom]{Foncteurs et morphismes}
\item \hyperref[equiv-section-phantom]{Catégories dérivées de variétés}
\item \hyperref[pione-section-phantom]{Groupes fondamentaux de schémas}
\item \hyperref[etale-cohomology-section-phantom]{Cohomologie étale}
\item \hyperref[crystalline-section-phantom]{Cohomologie cristalline}
\item \hyperref[proetale-section-phantom]{Cohomologie pro-étale}
\item \hyperref[relative-cycles-section-phantom]{Cycles relatifs}
\item \hyperref[more-etale-section-phantom]{Compléments de cohomologie étale}
\item \hyperref[trace-section-phantom]{La formule des traces}
\end{enumerate}
Espaces algébriques
\begin{enumerate}
\setcounter{enumi}{64}
\item \hyperref[spaces-section-phantom]{Espaces algébriques}
\item \hyperref[spaces-properties-section-phantom]{Propriétés des espaces algébriques}
\item \hyperref[spaces-morphisms-section-phantom]{Morphismes d'espaces algébriques}
\item \hyperref[decent-spaces-section-phantom]{Espaces algébriques décents}
\item \hyperref[spaces-cohomology-section-phantom]{Cohomologie des espaces algébriques}
\item \hyperref[spaces-limits-section-phantom]{Limites d'espaces algébriques}
\item \hyperref[spaces-divisors-section-phantom]{Diviseurs sur les espaces algébriques}
\item \hyperref[spaces-over-fields-section-phantom]{Espaces algébriques sur des corps}
\item \hyperref[spaces-topologies-section-phantom]{Topologies sur les espaces algébriques}
\item \hyperref[spaces-descent-section-phantom]{Descente et espaces algébriques}
\item \hyperref[spaces-perfect-section-phantom]{Catégories dérivées d'espaces}
\item \hyperref[spaces-more-morphisms-section-phantom]{Compléments sur les morphismes d'espaces}
\item \hyperref[spaces-flat-section-phantom]{Platitude sur les espaces algébriques}
\item \hyperref[spaces-groupoids-section-phantom]{Groupoïdes dans les espaces algébriques}
\item \hyperref[spaces-more-groupoids-section-phantom]{Compléments sur les groupoïdes dans les espaces}
\item \hyperref[bootstrap-section-phantom]{Amorçage}
\item \hyperref[spaces-pushouts-section-phantom]{Sommes amalgamées d'espaces algébriques}
\end{enumerate}
Thèmes de géométrie
\begin{enumerate}
\setcounter{enumi}{81}
\item \hyperref[spaces-chow-section-phantom]{Groupes de Chow des espaces}
\item \hyperref[groupoids-quotients-section-phantom]{Quotients de groupoïdes}
\item \hyperref[spaces-more-cohomology-section-phantom]{Compléments sur la cohomologie des espaces}
\item \hyperref[spaces-simplicial-section-phantom]{Espaces simpliciaux}
\item \hyperref[spaces-duality-section-phantom]{Dualité pour les espaces}
\item \hyperref[formal-spaces-section-phantom]{Espaces algébriques formels}
\item \hyperref[restricted-section-phantom]{Algébrisation des espaces formels}
\item \hyperref[spaces-resolve-section-phantom]{Résolution des surfaces revisitée}
\end{enumerate}
Théorie des déformations
\begin{enumerate}
\setcounter{enumi}{89}
\item \hyperref[formal-defos-section-phantom]{Théorie formelle des déformations}
\item \hyperref[defos-section-phantom]{Théorie des déformations}
\item \hyperref[cotangent-section-phantom]{Le complexe cotangent}
\item \hyperref[examples-defos-section-phantom]{Problèmes de déformation}
\end{enumerate}
Champs algébriques
\begin{enumerate}
\setcounter{enumi}{93}
\item \hyperref[algebraic-section-phantom]{Champs algébriques}
\item \hyperref[examples-stacks-section-phantom]{Exemples de champs}
\item \hyperref[stacks-sheaves-section-phantom]{Faisceaux sur les champs algébriques}
\item \hyperref[criteria-section-phantom]{Critères de représentabilité}
\item \hyperref[artin-section-phantom]{Axiomes d'Artin}
\item \hyperref[quot-section-phantom]{Espaces de Quot et de Hilbert}
\item \hyperref[stacks-properties-section-phantom]{Propriétés des champs algébriques}
\item \hyperref[stacks-morphisms-section-phantom]{Morphismes de champs algébriques}
\item \hyperref[stacks-limits-section-phantom]{Limites de champs algébriques}
\item \hyperref[stacks-cohomology-section-phantom]{Cohomologie des champs algébriques}
\item \hyperref[stacks-perfect-section-phantom]{Catégories dérivées de champs}
\item \hyperref[stacks-introduction-section-phantom]{Introduction aux champs algébriques}
\item \hyperref[stacks-more-morphisms-section-phantom]{Compléments sur les morphismes de champs}
\item \hyperref[stacks-geometry-section-phantom]{La géométrie des champs}
\end{enumerate}
Thèmes de théorie des espaces de modules
\begin{enumerate}
\setcounter{enumi}{107}
\item \hyperref[moduli-section-phantom]{Champs de modules}
\item \hyperref[moduli-curves-section-phantom]{Espaces de modules de courbes}
\end{enumerate}
Divers
\begin{enumerate}
\setcounter{enumi}{109}
\item \hyperref[examples-section-phantom]{Exemples}
\item \hyperref[exercises-section-phantom]{Exercices}
\item \hyperref[guide-section-phantom]{Guide bibliographique}
\item \hyperref[desirables-section-phantom]{Desiderata}
\item \hyperref[coding-section-phantom]{Style de programmation}
\item \hyperref[obsolete-section-phantom]{Obsolète}
\item \hyperref[fdl-section-phantom]{Licence de documentation libre GNU}
\item \hyperref[index-section-phantom]{Index généré automatiquement}
\end{enumerate}
\end{multicols}


\bibliography{my}
\bibliographystyle{amsalpha}

\end{document}
 
